\chapter{Modules, Zhu's algebra, and the rationality hierarchy}\label{ch:L12}
\begingroup\setlength{\emergencystretch}{3em}

\providecommand{\ltwelvevac}{\mathbf{1}}
\providecommand{\ltwelveY}{\mathcal{Y}}
\providecommand{\ltwelveO}{O}
\providecommand{\ltwelveA}{A}
\providecommand{\ltwelvecirc}{\circ}
\providecommand{\ltwelveHP}{\mathrm{HP}_0}

This chapter documents the release of the LeanVOA library named in the front matter; it was first written against an earlier state of the library (``the snapshot'' below), and checks that were
run at an earlier state of the library and not re-run for this edition are marked as such. Every Lean quotation below is taken from the files of the release and is
machine-checked by \texttt{scripts/textbook/sync.py} to occur verbatim (up to whitespace and
comments) in the source. The chapter covers the files assigned to it in
\texttt{Notes/Textbook/assignment.txt}: all of \leanpath{VOA/\allowbreak Modules/\allowbreak } except the
roughly three hundred files \leanpath{VOA/\allowbreak Modules/\allowbreak ZhuVerma*.lean} (the Frenkel--Zhu
tower and the hard direction of Zhu's correspondence, Chapter~\ref{ch:L13}), the super-module files
\leanpath{VOA/\allowbreak Modules/\allowbreak Super*.lean} (Chapter~\ref{ch:L17}), the three Poisson-trace files
(Chapter~\ref{ch:L16}) and \leanpath{VOA/\allowbreak Modules/\allowbreak FZBimodule.lean} (Chapter~\ref{ch:14}). About fifty files
and sixteen thousand lines of Lean remain. Where a statement of this chapter needs a definition
that lives in one of those other files (the Frenkel--Zhu spaces $A_n(V)$, the quotient
$\ltwelveHP(R_V)$ of Arakawa--Kawasetsu, the subspaces $C_m(V)$), we quote the definition here, with
its file, so that the chapter is self-contained.

\emph{Release revision.} The chapter was revised for the first public release, applying the release audit \texttt{Notes/\allowbreak Textbook/\allowbreak referee/\allowbreak
RELEASE\_AUDIT\_ZHU.md} (2026-10-27). The changes concern the use made of this chapter's rationality predicates by Zhu's
modularity theorem (Chapter~\ref{ch:L15}) and the status of the two halves of Zhu's correspondence (Chapter~\ref{ch:L13}); no
quotation of this chapter changed.

\emph{Second release revision (2026-10-28).} The chapter was revised again, applying finding F2 of \texttt{Notes/\allowbreak Textbook/\allowbreak referee/\allowbreak
RELEASE\_AUDIT\_NEW\_DECLS.md}. The first release revision said that the implication ``rational $\Rightarrow A(V)$
semisimple'', this chapter's named predicate \lean{ZhuAdmissibleSemisimple}, was unproved. That was false. The library proves
it (\lean{isSemisimpleRing_}\allowbreak\lean{A_of_}\allowbreak\lean{admissiblyRational} and \lean{zhuAdmissibleSemisimple_of_Pr},
in \leanpath{VOA/\allowbreak Huang/\allowbreak ZhuSemisimple.lean}), and Zhu's modularity theorem is now stated for
$C_2$-cofinite rational $V$ (Chapter~\ref{ch:L15}). The new quotation (Theorem~\ref{l12:thm:zhuss}) was copied at
that revision; no older quotation changed.

\emph{Third release revision (2026-10-28).} The chapter was revised once more, for the file \leanpath{VOA/\allowbreak Classification/\allowbreak COne/\allowbreak
LiteratureInputsProved.lean}, which this chapter now documents. Earlier revisions said that the library does not prove the
structure theorems of Dong--Li--Mason, Abe--Buhl--Dong and Li, and carries all seven of them as named hypotheses. That is no
longer true. Of the seven named literature predicates of Section~\ref{l12:sec:litnamed}, six are theorems of the library: four for
every VOA (\lean{ABDWeakIsAdmissible}, \lean{ABDIrreducibleWeakIsOrdinary}, \lean{ABDRegularOfC2AndRational},
\lean{ZhuAdmissibleSemisimple}) and two for every $C_2$-cofinite VOA (\lean{DLMSimpleAdmissibleOrdinary},
\lean{DLMFinitelyManyIrredAdmissible}); only \lean{LiC2CofiniteOfRegular} (H.~Li's ``regular $\Rightarrow C_2$-cofinite'')
is unproved, together with the two DLM predicates for VOAs that are not $C_2$-cofinite. The new material is
Theorem~\ref{l12:thm:litproved}; the status boxes of the rationality lattice and of strong rationality, and the checklist,
were corrected accordingly. No older quotation changed.

Chapter~\ref{ch:01} introduced, briefly, weak modules, ordinary modules, Zhu's product, the space
$\ltwelveO(V)$, Zhu's algebra $\ltwelveA(V)$, the zero-mode action on a lowest-weight space and the
library's notion of intertwining operator. We do not repeat those definitions in full; we recall
them in Section~\ref{l12:sec:recall} and refer back. This chapter goes further in five directions:
\begin{enumerate}
\item the \emph{full ladder of module categories} used in the literature --- weak, admissible
  ($\N$-gradable), logarithmic (grading-restricted generalised) and ordinary --- and the theorems
  relating them (Section~\ref{l12:sec:categories});
\item the \emph{associative algebra} $\ltwelveA(V)$, which the library proves associative
  unconditionally, its action on lowest-weight spaces, the forward half of Zhu's correspondence, and
  a Zhu-algebra simplicity criterion (Section~\ref{l12:sec:zhu});
\item \emph{$C_2$-cofiniteness}, the Poisson algebra $V/C_2(V)$, finite generation, the Frenkel--Zhu
  and Dong--Li--Mason higher Zhu algebras $A_n(V)$ and their finiteness under $C_2$
  (Sections~\ref{l12:sec:c2} and~\ref{l12:sec:an});
\item the \emph{contragredient module} $M'$, which the library constructs in full
  (Section~\ref{l12:sec:contra}), and the library's fusion coefficient, which we prove to be
  identically zero (Section~\ref{l12:sec:fusion});
\item the \emph{rationality hierarchy}: regular, admissibly rational (``rational''),
  Gannon-rational, holomorphic, strongly finite and strongly rational, and the implications among
  them, together with the structure theorems of Dong--Li--Mason, Abe--Buhl--Dong and Li, stated as named predicates;
  all but one of these (Li's theorem) are now proved in the library, the two DLM ones for $C_2$-cofinite $V$
  (Section~\ref{l12:sec:rationality}).
\end{enumerate}

\paragraph{The findings of this chapter, in brief.} Most of what is documented here is faithful to
the literature and is proved without hypotheses, axiom-clean, and without \lean{VOA.pbw} (the PBW theorem of Chapter~\ref{ch:L02}; formerly a project axiom, now a theorem). The fidelity checks did, however, turn up several definitions that do not mean what
their names say. We list them here so that a reader skimming the chapter does not miss them; each
is explained in a red warning box where it occurs.
\begin{itemize}
\item \lean{VOA.fusionCoeff} is \emph{identically zero}. Chapter~\ref{ch:01} argued this on paper; we
  have now proved it in Lean (library theorem \lean{fusionCoeff_eq_zero}, Theorem~\ref{l12:thm:fusionzero}).
  No theorem of the library uses \lean{fusionCoeff}, so nothing proved becomes false or vacuous.
\item \lean{VOA.Unitary} (``unitary module'') is \emph{unsatisfiable} on every module of dimension at
  least two (library theorem \lean{old_unitary_isEmpty},
  Proposition~\ref{l12:prop:unitary}), and \lean{VOA.SelfContragredient} is satisfied by
  \emph{every} module (the library itself proves this). Neither has a consumer.
\item \lean{VOA.InducedAssignment}, the interface through which the file
  \leanpath{VOA/\allowbreak Modules/\allowbreak ZhuInduced.lean} expresses ``the hard direction of Zhu's correspondence'',
  is \emph{uninhabited} (the library proves \lean{not_}\allowbreak\lean{nonempty_}\allowbreak\lean{inducedAssignment}), so the
  round-trip theorems stated with it are vacuous. The genuine bijection is proved elsewhere
  (Chapter~\ref{ch:L13}).
\item \lean{VOA.IsLowestWeight N h} does not require $M_h\neq0$, and requires \emph{all} weights to lie
  in the single coset $h+\N$. Theorems with this hypothesis are true but become statements about the
  zero space when $M_h = 0$.
\item \textbf{Rationality implies that $A(V)$ is semisimple, in Lean.} The named predicate \lean{ZhuAdmissibleSemisimple}
  (``rational $\Rightarrow A(V)$ semisimple'') is a theorem, for every $V$: \lean{isSemisimpleRing_}\allowbreak\lean{A_of_}\allowbreak\lean{admissiblyRational}
  (Theorem~\ref{l12:thm:zhuss}). So Zhu's modularity theorem (Chapter~\ref{ch:L15}) holds in Lean for every $C_2$-cofinite
  $V$ that satisfies the library's rationality predicate \lean{AdmissiblyRational}, in particular for the strongly rational
  lattice VOAs $V_{\sqrt{2m}\Z}$ (Chapter~\ref{ch:L22}). For $C_2$-cofinite $V$, Dong--Li--Mason
  rationality implies \lean{AdmissiblyRational} (\lean{admissiblyRational_of_dlmRational}, Theorem~\ref{l12:thm:dlmrat});
  what remains open is only the converse (Section~\ref{l12:warn:admcr}). \emph{Correction:} the first release
  revision of this chapter said that the implication was unproved and that $V_{\sqrt{2m}\Z}$ was not covered; both claims were
  wrong.
\item \textbf{Six of the seven named literature predicates are theorems.} Abe--Buhl--Dong's Proposition~5.7, Corollary~5.8
  and Theorem~4.5 (hard direction) and Zhu's semisimplicity theorem hold in Lean for every $V$; Dong--Li--Mason's
  Theorem~8.1(b),(c) hold for every $C_2$-cofinite $V$ (Theorem~\ref{l12:thm:litproved}). Consequently a strongly rational
  VOA is regular and Gannon-rational with no literature hypothesis (\lean{StronglyRational.regular},
  \lean{StronglyRational.gannonRational'}). Only H.~Li's ``regular $\Rightarrow C_2$-cofinite'' (\lean{LiC2CofiniteOfRegular})
  remains a named, unproved hypothesis, and the one theorem of this chapter that uses it is the equivalence
  \lean{regular_iff_c2Cofinite_and_admissiblyRational}.
\item The library's $A_n(V)$ (file \leanpath{VOA/\allowbreak Modules/\allowbreak ZhuVermaAn.lean}) omits the generators
  $(L_{-1}+L_0)u$ that Dong--Li--Mason include in $O_n(V)$, so it surjects onto, and may be larger
  than, the literature's $A_n(V)$; and the library's \lean{A₂ W} is \emph{not} the second higher Zhu
  algebra but Zhu's $C_2$-quotient $V/C_2(V)$.
\end{itemize}

%======================================================================
\section{Conventions recalled}\label{l12:sec:recall}
%======================================================================

Throughout, $V$ denotes a vertex operator algebra in the sense of the library (the Lean
\lean{W : VertexOperatorAlgebra 𝒱} with underlying vector space \lean{𝒱}); Chapter~\ref{ch:01},
Section~\ref{l03:sec:axioms}, and Chapter~\ref{ch:L03} give the definition. We recall the features that
matter in this chapter.
\begin{itemize}
\item $V = \bigoplus_{n\in\Z}V_n$ with each $V_n$ finite-dimensional, $V_n = 0$ for $n<0$ and
  $V_0 = \C\ltwelvevac$. These are \emph{structure fields} of every library VOA
  (\lean{positivity}, \lean{vacuum_grade}, \lean{finrank}); so every VOA in the library is of
  \emph{CFT type}. Many literature theorems about modules assume CFT type; here it is automatic. In
  particular the zero vector space with $\ltwelvevac = 0$ is a library VOA (the ``zero VOA''), which
  is why simplicity hypotheses often have to be carried to exclude it.
\item The $n$-th product is $u_{(n)}v$, written \lean{W.smul u n v}. For a formal series
  $f = \sum_k f_k z^k$ the library writes \lean{mode f n} for the coefficient
  $f_{-n-1}$, so that \lean{mode (Y u) n} is the operator $u_{(n)}$ in the expansion
  $Y(u,z) = \sum_n u_{(n)}z^{-n-1}$. The Virasoro operators are
  $L_n = \omega_{(n+1)}$, written \lean{W.L n}; in particular $L_0 = \omega_{(1)}$ and
  $L_{-1} = \omega_{(0)}$.
\item \lean{cbz m i} is the generalised binomial coefficient $\binom mi$ for $m\in\Z$ and $i\in\Z$, equal to
  $0$ when $i<0$; \lean{negOnePow i} is $(-1)^i$. A sum \lean{∑ᶠ i, f i} is a sum over the (finite)
  support of $f$; Mathlib assigns it the value $0$ if the support is infinite, so every time such a
  sum appears in an \emph{axiom} one must check that the support is finite (the dictionary of
  Definition~\ref{l12:def:weakmodule} does this for the module Jacobi identity).
\item \lean{Submodule ℂ X} is a linear subspace of $X$; \lean{⊥} and \lean{⊤} are the zero subspace
  and the whole space; \lean{DirectSum.IsInternal w} says that $X$ is the internal direct sum of the
  family of subspaces $w(\alpha)$; \lean{FiniteDimensional ℂ X} and \lean{Module.finrank ℂ X} are
  finite-dimensionality and dimension (Mathlib's convention: \lean{finrank} of an
  \emph{infinite}-dimensional space is $0$). \lean{X ⧸ S} is the quotient space.
\end{itemize}
From Chapter~\ref{ch:01} we use: \lean{WeakModule} (Definition~\ref{l12:def:weakmodule}), with
action $u^M_{(n)}x$ = \lean{A.act u n x} = \lean{mode (A.YM u) n x}; \lean{OrdVModule}
(Definition~\ref{l12:def:ordmodule}) with weight spaces \lean{N.weight α} $=M_\alpha$; the adjoint
module (Definition~\ref{l12:def:adjoint}); Zhu's product \lean{zhuMul}, \lean{O W}, \lean{A W}
(Definition~\ref{ch12:def:zhumul}); the zero mode \lean{oMod} and \lean{IsLowestWeight}
(Definition~\ref{ch12:def:omod}); Zhu's multiplicativity theorem on the lowest-weight space
(Theorem~\ref{l12:thm:zhuaction}); and the library's intertwiners
(Definition~\ref{l12:def:intertwiner}). Weak submodules \lean{WeakModule.Submodule'} are subspaces
stable under every $u^M_{(n)}$, and \lean{WeakModule.Hom} are linear maps commuting with every
$u^M_{(n)}$. A \emph{sub-VOA} \lean{VOASubalgebra W} (Chapter~\ref{ch:01},
Section~\ref{l04:sec:quot}) is a graded subspace containing $\ltwelvevac$ and $\omega$ and closed
under all products; \lean{S.toVOA} is the VOA it carries.

\begin{physics}
Throughout, the physical dictionary is the one of Chapter~\ref{ch:01}, Section~\ref{ch01:sec:physics}: $V$ is the
\emph{chiral algebra} of a two-dimensional conformal field theory (the holomorphic fields and their operator
product expansions), $u_{(n)}$ are the Laurent modes of the field $Y(u,z)$, $L_0$ measures conformal weight
(energy on the cylinder is $L_0-c/24$), and CFT type says that the vacuum is the unique state of weight $0$ and
that there are no states of negative energy. A \emph{module} is a \emph{sector} (a superselection sector, a
representation of the chiral algebra on a space of states created by a non-vacuum primary field). This chapter
is about the representation theory of the chiral algebra, i.e.\ about the \emph{spectrum of sectors} of a CFT.
\end{physics}

%======================================================================
\section{The ladder of module categories}\label{l12:sec:categories}
%======================================================================

The literature on vertex operator algebras uses (at least) four notions of ``module'', which form
a chain of inclusions
\begin{equation*}
 \text{ordinary}\ \subseteq\ \text{logarithmic}\ \subseteq\ \text{admissible}\ \subseteq\ \text{weak}.
\end{equation*}
They differ only in what is demanded of a grading. A weak module has none. An admissible module has
an auxiliary $\N$-grading compatible with the modes, but its pieces may be infinite-dimensional and
$L_0$ need not act semisimply. A logarithmic module is graded by \emph{generalised} eigenspaces of
$L_0$, each finite-dimensional. An ordinary module is graded by honest eigenspaces of $L_0$, each
finite-dimensional. Rationality conditions are statements of the form ``every module of such-and-such
kind is completely reducible'', and the kind matters, so we must be precise about all four.

\subsection{Weak modules and the derivative property}

Physically, a module over a VOA $V$ is a \emph{sector} of the chiral CFT: a space of states on which all fields of
$V$ act, with the same operator product expansions, but which is not itself created from the vacuum. We begin with the
weakest notion, which has no grading at all.

\begin{definition}[Weak module; Gannon Def.~5.3.1, FHL]\label{l12:def:weakmodule}
Let $V$ be a VOA. A \emph{weak $V$-module} is a vector space $M$ together with a linear map
$Y_M:V\to\End(M)[[z^{\pm1}]]$, $Y_M(u,z) = \sum_{n\in\Z}u^M_{(n)}z^{-n-1}$, such that
\begin{enumerate}[label=(M\arabic*)]
\item $Y_M(\ltwelvevac,z) = \mathrm{id}_M$;
\item each $Y_M(u,z)$ is a field: for all $u\in V$ and $x\in M$, $u^M_{(n)}x = 0$ for $n\gg0$;
\item the Borcherds (Jacobi) identity holds: for all $u,v\in V$, $x\in M$ and $\ell,m,n\in\Z$,
\begin{equation*}
 \sum_{i\ge0}(-1)^i\binom\ell i\Big(u^M_{(\ell+m-i)}v^M_{(n+i)}x - (-1)^\ell v^M_{(\ell+n-i)}u^M_{(m+i)}x\Big)
 = \sum_{i\ge0}\binom mi(u_{(\ell+i)}v)^M_{(m+n-i)}x .
\end{equation*}
\end{enumerate}
\end{definition}
\noindent File \leanpath{VOA/\allowbreak Modules/\allowbreak WeakModule.lean}.
\begin{leancode}
structure WeakModule (W : VertexOperatorAlgebra 𝒱) (M : Type v)
    [AddCommGroup M] [Module ℂ M] where
  YM : 𝒱 →ₗ[ℂ] (End ℂ M)⟦z^±⟧
  vacuum_YM : YM W.vacuum = 1
  truncation : ∀ (u : 𝒱) (x : M), ∃ N : ℤ, ∀ n : ℤ, N ≤ n → mode (YM u) n x = 0
  jacobi : ∀ (u v : 𝒱) (x : M) (ℓ m n : ℤ),
    (∑ᶠ i : ℤ, (negOnePow i * cbz ℓ i) •
        (mode (YM u) (ℓ + m - i) (mode (YM v) (n + i) x)
          - negOnePow ℓ • mode (YM v) (ℓ + n - i) (mode (YM u) (m + i) x)))
      = ∑ᶠ i : ℤ, cbz m i • mode (YM (W.smul u (ℓ + i) v)) (m + n - i) x

def act (A : WeakModule W M) (u : 𝒱) (n : ℤ) (x : M) : M := mode (A.YM u) n x
\end{leancode}
(\lean{VOA.WeakModule}, \lean{VOA.WeakModule.act}; \lean{act} lies in the namespace \lean{WeakModule}.)

\begin{dictionary}
\lean{(End ℂ M)⟦z^±⟧} is the space \lean{FormalDist (End ℂ M)} of doubly infinite series of endomorphisms
(Definition~\ref{l01:def:formaldist}); \lean{mode (YM u) n} is $u^M_{(n)}$, the coefficient of $z^{-n-1}$
(Definition~\ref{l01:def:mode}); \lean{YM W.vacuum = 1} uses the unit distribution $\mathrm{id}_Mz^0$, i.e.\ (M1);
\lean{truncation} is (M2), the field condition (Definition~\ref{l01:def:field}) in mode indexing: ``there is $N$
such that $u_{(n)}x = 0$ for all $n\ge N$''; \lean{jacobi} is (M3), word for word the same formula as
\lean{VertexOperatorAlgebra.jacobi} (Theorem~\ref{l03:thm:borcherds}), where the \emph{algebra} product
$u_{(\ell+i)}v$ on the right is \lean{W.smul u (ℓ + i) v}. \lean{negOnePow i * cbz ℓ i} is $(-1)^i\binom\ell i$ and
\lean{cbz m i} the binomial $\binom mi$ with the convention that it vanishes for $i<0$ (Definition~\ref{l01:def:gbinom}).
\lean{A.act u n x} is the module action $u^M_{(n)}x$.

\textbf{The \lean{∑ᶠ} convention in an axiom.} Here the Jacobi identity is an \emph{axiom}, and an axiom stated with
\lean{∑ᶠ} could be weaker than intended if a sum had infinite support (the sum would then be read as $0$, and the axiom
would assert $0=\cdots$). It does not happen: both supports are finite for every weak module. On the left, $i<0$ gives
$\binom\ell i = 0$, for $i\ge0$ large $v^M_{(n+i)}x = 0$ and $u^M_{(m+i)}x = 0$ by (M2). On the right, $i<0$ gives
$\binom mi=0$ and $u_{(\ell+i)}v = 0$ for $i\gg0$ by the regularity axiom of $V$. So the axiom is the intended one.

\textbf{Why Jacobi rather than locality.} For a module, locality of the fields $Y_M(u,z)$ together with (M1)--(M2) does
\emph{not} imply the Jacobi identity: the relation between $Y_M(u_{(n)}v,z)$ and $Y_M(u,z)$, $Y_M(v,z)$ is extra
information, and there is no vacuum in $M$ to run the Goddard argument (Theorem~\ref{l03:thm:goddard}). The library
therefore takes the Jacobi identity as the axiom, as do Frenkel--Huang--Lepowsky, Lepowsky--Li (Def.~4.1.1) and Gannon.
This is the standard definition.

\textbf{What is not required, and why that is harmless.} No Virasoro axiom and no $L_{-1}$-derivative axiom is imposed on
$M$. The Virasoro relations for $L^M_n = \omega^M_{(n+1)}$ follow from (M3) and the VOA axioms of $V$, and the
$L_{-1}$-derivative property is proved below (Theorem~\ref{l12:thm:covariance}). \textbf{Could the definition be
vacuous?} The zero space is a weak module, and so is $V$ itself (Definition~\ref{l12:def:adjoint}); the axioms are
consistent and have the intended examples. Weak homomorphisms (\lean{WeakModule.Hom}: linear maps commuting with every
$u_{(n)}$) and weak submodules (\lean{WeakModule.Submodule'}: subspaces stable under every $u_{(n)}$) are defined in the
same file in the evident way.
\end{dictionary}

\begin{physics}
A weak module is a space of states $M$ on which the fields $Y_M(u,z)$ act with the same Jacobi identity as in the
vacuum sector: an algebraic shadow of a superselection sector, namely the space generated from the origin by an
operator that is not itself in the vacuum module. The Jacobi identity contains both the commutation relations of the modes and the
operator product expansion (the right side involves the OPE coefficients $u_{(\ell+i)}v$ of $V$). It is
``weak'' because no energy grading is demanded, so it also admits unphysical sectors with energy unbounded below;
the physical sectors are the graded modules of Definition~\ref{l12:def:ordmodule}.
\end{physics}

\begin{status}
Definitions; axiom-clean (\texttt{\#print axioms} of \lean{VOA.WeakModule} and \lean{VOA.WeakModule.act} is
\lean{[propext, Classical.}\allowbreak\lean{choice, Quot.}\allowbreak\lean{sound]}, run for this chapter with
\texttt{lake env lean -DmaxSynthPendingDepth=3}). Satisfiable: the adjoint module is an instance.
\end{status}

\begin{theorem}[Module commutator formula]\label{l12:thm:modcomm}
For a weak module $M$, all $u,v\in V$, $x\in M$ and $m,n\in\Z$,
\begin{equation*}
 u^M_{(m)}v^M_{(n)}x - v^M_{(n)}u^M_{(m)}x \;=\; \sum_{i\ge0}\binom mi\,(u_{(i)}v)^M_{(m+n-i)}x ,
\end{equation*}
a finite sum.
\end{theorem}
\noindent File \leanpath{VOA/\allowbreak Modularity/\allowbreak ZhuRecursionProof.lean}.
\begin{leancode}
theorem mode_commutator (A : WeakModule W M) (u v : 𝒱) (x : M) (m n : ℤ) :
    mode (A.YM u) m (mode (A.YM v) n x) - mode (A.YM v) n (mode (A.YM u) m x)
      = ∑ᶠ i : ℤ, cbz m i • mode (A.YM (W.smul u i v)) (m + n - i) x
\end{leancode}

\begin{dictionary}
This is the formula $[u_{(m)},v_{(n)}] = \sum_i\binom mi(u_{(i)}v)_{(m+n-i)}$ of Theorem~\ref{l03:thm:comm} on $M$ instead of
$V$ (the declaration is \lean{VOA.mode_commutator}, to be distinguished from the VOA-level
\lean{VOA.VertexOperatorAlgebra.mode_commutator}). Unlike the VOA-level statement it is stated as an identity on a vector $x$ with a
\lean{∑ᶠ}, not with an existential bound $N$. The \lean{finsum} is a genuine finite sum for the same reasons as in
Definition~\ref{l12:def:weakmodule}: $\binom mi = 0$ for $i<0$ and $u_{(i)}v = 0$ for $i\gg0$. \textbf{Hypotheses:} $M$ is a weak
module; nothing else. The statement is the textbook one, FHL, Lepowsky--Li (\S4.1) and Gannon (\S5.3.1).
A version with an explicit finite range is \lean{act_comm} in \leanpath{VOA/\allowbreak Moonshine/\allowbreak VnatBlock.lean}.
\end{dictionary}

\begin{physics}
The modes of the fields satisfy the same commutation relations on every sector as on the vacuum module: the
current algebra $[J_m,J_n] = km\delta_{m+n}$ or the Virasoro algebra with the same central charge $c$ acts on every
representation. This is the statement that the operator product expansion of chiral fields is a property of the algebra
of local operators and not of the state on which they act, and it is the origin of the universal Ward identities that
hold in all sectors.
\end{physics}

\textbf{Proof on paper.} Put $\ell = 0$ in (M3). Since $\binom0i = \delta_{i,0}$, the left side collapses to its $i=0$ term,
$u^M_{(m)}v^M_{(n)}x - v^M_{(n)}u^M_{(m)}x$, and the right side is $\sum_{i}\binom mi(u_{(i)}v)^M_{(m+n-i)}x$. \qed

\begin{status}
Unconditional. \texttt{\#print axioms} \lean{VOA.mode\_commutator} $=$
\lean{[propext, Classical.}\allowbreak\lean{choice, Quot.}\allowbreak\lean{sound]} (run for this chapter).
\end{status}

Recall the definition of a weak module (Definition~\ref{l12:def:weakmodule}): a linear map
$Y_M:V\to\End(M)[[z^{\pm1}]]$ with $Y_M(\ltwelvevac,z) = \mathrm{id}$, lower truncation, and the Jacobi
(Borcherds) identity in all indices. No $L_{-1}$-derivative axiom is imposed, because it follows. The
library proves this.

\begin{theorem}[Derivative property on a weak module]\label{l12:thm:covariance}
Let $M$ be a weak $V$-module. For all $a\in V$, $n\in\Z$, $x\in M$,
\begin{equation*}
 (L_{-1}a)^M_{(n)}x = -n\, a^M_{(n-1)}x ,
\end{equation*}
i.e.\ $Y_M(L_{-1}a,z) = \partial_z Y_M(a,z)$.
\end{theorem}
\noindent File \leanpath{VOA/\allowbreak Modules/\allowbreak Translation.lean}.
\begin{leancode}
theorem act_T_covariance (A : WeakModule W M) (a : 𝒱) (n : ℤ) (x : M) :
    A.act (W.L (-1) a) n x = (-(n : ℂ)) • A.act a (n - 1) x := by
\end{leancode}

\begin{dictionary}
\lean{A.act u n x} is $u^M_{(n)}x$; \lean{W.L (-1) a} is $L_{-1}a$; \lean{(-(n : ℂ)) • y} is $-n\,y$
(the integer $n$ cast to $\C$). There are no hypotheses beyond ``$M$ is a weak module''. The
generating-function form: $\partial_z\sum_n a_{(n)}z^{-n-1} = \sum_n(-n-1)a_{(n)}z^{-n-2}
= \sum_n (-n)a_{(n-1)}z^{-n-1}$, whose coefficient of $z^{-n-1}$ is exactly the right-hand side. So the
statement is the textbook one (Lepowsky--Li, Prop.~4.1.5; Frenkel--Huang--Lepowsky, §4.1), with no
weakening.
\end{dictionary}

\textbf{Proof on paper.} Put $m=0$ in the module Jacobi identity. Since $\binom0i = \delta_{i,0}$, the
right-hand side becomes $(u_{(\ell)}v)^M_{(n)}x$ and the identity is the \emph{iterate formula}
\begin{equation*}
 (u_{(\ell)}v)^M_{(n)}x = \sum_{i\ge0}(-1)^i\binom\ell i\Big(u^M_{(\ell-i)}v^M_{(n+i)}x - (-1)^\ell v^M_{(\ell+n-i)}u^M_{(i)}x\Big).
\end{equation*}
In $V$ we have $L_{-1}a = a_{(-2)}\ltwelvevac$ (creation property). Take $u = a$, $v=\ltwelvevac$,
$\ell = -2$. On a module $\ltwelvevac^M_{(k)} = \delta_{k,-1}$, so of the two infinite families only the
term $i = -1-n$ of the first and $i = n-1$ of the second survive (when they have $i\ge0$). Their
coefficients are $(-1)^{-1-n}\binom{-2}{-1-n}$ and $-(-1)^{n-1}\binom{-2}{n-1}$, and using
$\binom{-2}{j} = (-1)^j(j+1)$ one checks that in every case (whether $n\le-1$, $n\ge1$ or $n=0$) the total
coefficient of $a^M_{(n-1)}x$ is $-n$. \qed

\begin{status}
Proved, no hypotheses. \texttt{\#print axioms VOA.act\_T\_covariance} gives
\lean{[propext, Classical.choice, Quot.sound]} (run for this chapter).
\end{status}

\begin{physics}
The derivative property says that the field $Y_M(L_{-1}a,z)$ of the descendant $L_{-1}a$ is the derivative of the
field of $a$ in \emph{every} sector: $L_{-1}$ generates translations $z\mapsto z+\epsilon$ in all
representations, not only in the vacuum sector. Physically it is the statement that the stress tensor's
$L_{-1}$ is the momentum operator, and a sector that violated it would not be translation covariant. It is
the module half of the conformal Ward identity.
\end{physics}

\subsection{Admissible (\texorpdfstring{$\N$}{N}-gradable) modules}

\begin{definition}[Admissible module; Zhu 1996, §1.4; Dong--Li--Mason 1998, Def.~2.2]\label{l12:def:admissible}
An \emph{admissible} $V$-module is a weak module $M$ together with a decomposition
$M = \bigoplus_{k\in\N}M(k)$ such that for every homogeneous $a\in V_w$, every $p\in\Z$ and every
$k$,
\begin{equation*}
 a^M_{(p)}M(k)\subseteq M(k + w - p - 1)
\end{equation*}
(with $M(j) = 0$ for $j<0$). It is \emph{irreducible} if $M\neq0$ and its only weak submodules are
$0$ and $M$; \emph{completely reducible} if every weak submodule has a complementary weak submodule.
\end{definition}
\noindent File \leanpath{VOA/\allowbreak Modules/\allowbreak AdmissibleModule.lean}.
\begin{leancode}
structure AdmissibleModule (W : VertexOperatorAlgebra 𝒱) (M : Type v)
    [AddCommGroup M] [Module ℂ M] extends WeakModule W M where
  lev : ℤ → Submodule ℂ M
  lev_eq_bot_of_neg : ∀ k : ℤ, k < 0 → lev k = ⊥
  isInternal : DirectSum.IsInternal lev
  act_lev : ∀ ⦃w p k : ℤ⦄ ⦃a : 𝒱⦄, a ∈ W.grading w → ∀ ⦃x : M⦄, x ∈ lev k →
    mode (YM a) p x ∈ lev (k + (w - p - 1))

def IsIrreducible (N : AdmissibleModule W M) : Prop :=
  Nontrivial M ∧ ∀ P : WeakModule.Submodule' N.toWeakModule, P.carrier = ⊥ ∨ P.carrier = ⊤

def IsCompletelyReducible (N : AdmissibleModule W M) : Prop :=
  ∀ P : WeakModule.Submodule' N.toWeakModule,
    ∃ Q : WeakModule.Submodule' N.toWeakModule, IsCompl P.carrier Q.carrier

def adjointAdmissibleModule (W : VertexOperatorAlgebra 𝒱) : AdmissibleModule W 𝒱 where
  toWeakModule := adjointWeakModule W
  lev := W.grading
\end{leancode}

\begin{dictionary}
\textbf{Decoding.} \lean{extends WeakModule W M} means an admissible module \emph{is} a weak module
(its field \lean{toWeakModule}) plus data. \lean{lev k} is $M(k)$; the grading is indexed by $\Z$ with
the negative pieces forced to be zero (\lean{lev_eq_bot_of_neg}), which is an $\N$-grading in $\Z$
bookkeeping (it avoids truncated subtraction in $k+w-p-1$). \lean{a ∈ W.grading w} says $a\in V_w$.
The binders \lean{⦃ ⦄} are ``strict implicit'' arguments: the statement is the same universally
quantified statement. \lean{IsCompl P Q} says $P\cap Q = 0$ and $P+Q = M$. \lean{Nontrivial M}
says $M\neq0$.

\textbf{Same as the literature?} The structure is DLM's Definition~2.2 (= Zhu's ``$\N$-gradable weak
module'') with two conventions to note.
\begin{enumerate}
\item DLM sometimes add $M(0)\neq0$ when $M\neq0$ (a normalisation of the grading). The library does
  not; this only enlarges the class by allowing shifts of the grading, and changes nothing about which
  weak modules are admissible.
\item \emph{Irreducible} and \emph{completely reducible} quantify over \emph{weak} submodules (any
  subspace stable under all modes), not over \emph{graded} submodules. DLM's ``simple admissible''
  means \emph{graded}-simple (``$0$ and $M$ are the only graded submodules''), and their complete
  reducibility means ``direct sum of graded-simple admissible submodules''. So the library's
  irreducibility is \emph{stronger} than DLM's (fewer modules count as irreducible) and the
  relationship between the two complete-reducibility notions is discussed in
  the warning in Section~\ref{l12:warn:admcr} below.
\end{enumerate}
The adjoint module $V$ with $M(k) = V_k$ is admissible (\lean{adjointAdmissibleModule}; the mode shift
is the grading axiom of $V$), so the category is never empty.
\end{dictionary}

\begin{physics}
An admissible module is a sector whose states come in levels $M(0),M(1),\dots$ above a ground floor, with modes
$a_{(p)}$ of a weight-$w$ field raising the level by $w-p-1$ (so $a_{(w-1)}$, the zero mode, preserves it). It
is the minimal structure one needs to talk about ``the ground states'' of a sector without yet assuming that
the Hamiltonian $L_0$ is diagonalisable or that each level has finite degeneracy. In logarithmic CFT, where
$L_0$ has Jordan blocks, admissible modules are the natural home of the indecomposable sectors.
\end{physics}

\begin{status}
Definitions; the adjoint instance is proved. \texttt{\#print axioms} of the file's declarations:
\lean{[propext, Classical.choice, Quot.sound]} (file docstring; we re-ran it for the downstream
theorems below).
\end{status}

\subsection{Ordinary modules, and the global lower bound}

\begin{definition}[Ordinary module; irreducible; completely reducible; Gannon Def.~5.3.1(c)]\label{l12:def:ordmodule}
A \emph{$V$-module} (an \emph{ordinary} module) is a weak module $M$ together with a decomposition
$M = \bigoplus_{\alpha\in\C}M_\alpha$ such that $L^M_0 = \omega^M_{(1)}$ acts on $M_\alpha$ as the scalar $\alpha$, each
$M_\alpha$ is finite-dimensional, and the real parts of the weights $\alpha$ with $M_\alpha\neq0$ are bounded below. It is
\emph{irreducible} if $M\neq0$ and its only weak submodules are $0$ and $M$, and \emph{completely reducible} if every weak
submodule has a complementary weak submodule.
\end{definition}
\noindent File \leanpath{VOA/\allowbreak Modules/\allowbreak WeakModule.lean}.
\begin{leancode}
structure OrdVModule (W : VertexOperatorAlgebra 𝒱) (M : Type v)
    [AddCommGroup M] [Module ℂ M] extends WeakModule W M where
  weight : ℂ → Submodule ℂ M
  isInternal : DirectSum.IsInternal weight
  L0_weight : ∀ (α : ℂ) (x : M), x ∈ weight α → mode (YM W.conformalVector) 1 x = α • x
  finrank : ∀ α, FiniteDimensional ℂ (weight α)
  boundedBelow : ∃ B : ℝ, ∀ α : ℂ, weight α ≠ ⊥ → B ≤ α.re

def IsIrreducible (N : OrdVModule W M) : Prop :=
  Nontrivial M ∧ ∀ P : WeakModule.Submodule' N.toWeakModule, P.carrier = ⊥ ∨ P.carrier = ⊤

def IsCompletelyReducible (N : OrdVModule W M) : Prop :=
  ∀ P : WeakModule.Submodule' N.toWeakModule,
    ∃ Q : WeakModule.Submodule' N.toWeakModule, IsCompl P.carrier Q.carrier
\end{leancode}
(\lean{VOA.OrdVModule}, \lean{VOA.OrdVModule.IsIrreducible}, \lean{VOA.OrdVModule.IsCompletelyReducible}; the last two lie in
the namespace \lean{OrdVModule}.)

\begin{dictionary}
\textbf{Name.} \lean{OrdVModule} (``ordinary $V$-module'') avoids a clash with Mathlib's \lean{Module}. \lean{extends WeakModule W M}
means that an ordinary module \emph{is} a weak module (the field \lean{toWeakModule}) together with the extra data.
\textbf{Fields.} \lean{weight α} is the subspace $M_\alpha$; \lean{isInternal} says that $M$ is the internal direct sum of the
$M_\alpha$; \lean{L0_weight} says $L_0^M=\alpha$ on $M_\alpha$ (\lean{mode (YM ω) 1} is $\omega^M_{(1)} = L^M_0$); together these make
$M_\alpha$ exactly the $\alpha$-eigenspace of $L^M_0$, so that $L_0^M$ is diagonalisable (if $L_0x=\alpha x$ and
$x=\sum x_\beta$, then $(\beta-\alpha)x_\beta=0$ for all $\beta$, so $x=x_\alpha$); \lean{finrank} gives finite-dimensional weight
spaces; \lean{boundedBelow} is a \emph{global} real bound $B\le\mathrm{Re}\,\alpha$ on all occurring weights (\lean{α.re} is the real
part). \textbf{Predicates.} \lean{Nontrivial M} says that $M$ has two distinct elements, i.e.\ $M\neq0$; \lean{P.carrier = ⊥ ∨
P.carrier = ⊤} says that the underlying subspace of the weak submodule $P$ is $0$ or $M$; \lean{IsCompl P Q} says
$P\cap Q = 0$ and $P+Q = M$. Both \lean{IsIrreducible} and \lean{IsCompletelyReducible} quantify over \emph{weak} submodules
(\lean{WeakModule.Submodule'}: subspaces stable under every $u_{(n)}$), not merely graded ones. For ordinary modules this is the
standard meaning, since every weak submodule of an ordinary module is automatically graded (it is stable under $L_0=\omega_{(1)}$, and
$L_0$ is diagonalisable on $M$).

\textbf{Fidelity.} The same as the standard definition (FHL, Lepowsky--Li Def.~4.1.6, Gannon Def.~5.3.1(c)) with complex weights; some authors (Dong--Li--Mason) require weights in $h+\N$,
which is a special case. The one deviation is the global lower bound (the warning in Section~\ref{l12:sec:categories}): the library demands one real
bound for \emph{all} weights, which is stronger than the coset-wise condition of FHL. A twisted sector with weights in
$\tfrac32+\tfrac12\Z$ is \emph{not} an \lean{OrdVModule} of the untwisted algebra; it is a \lean{TwistedModule} (Chapter~\ref{ch:06}).
\textbf{Could the definition be vacuous?} The adjoint module (Definition~\ref{l12:def:adjoint}) is an \lean{OrdVModule}, and the
Moonshine sectors are further examples, so the class is non-empty; and the notions \lean{IsIrreducible},
\lean{IsCompletelyReducible} are satisfiable (the adjoint module is irreducible exactly when $V$ is simple, see Definition~\ref{l12:def:phi}). The distinction from the admissible versions of the same names is in Definition~\ref{l12:def:admissible}.
\end{dictionary}

\begin{physics}
An ordinary module is a physical sector: a graded Hilbert space of states $M=\bigoplus_\alpha M_\alpha$ with $M_\alpha$ the
eigenspace of $L_0$ of conformal weight $\alpha$ and finite degeneracy. The weights are the scaling dimensions of the
primary and descendant states, and on the cylinder the energy of a state is $L_0-c/24$. Irreducibility means that no proper
subspace of states is closed under all the modes, i.e.\ the sector is a single irreducible representation of the chiral
algebra and cannot be split into independent superselection sectors. For a rational CFT the irreducible modules are the finitely many
superselection sectors, they label the primary fields, and their characters (Definition~\ref{l12:def:char}) are the building blocks of the partition
function. Finite dimensionality of each $M_\alpha$ is the statement of a discrete spectrum with finite degeneracy.
\end{physics}

\textbf{Proof on paper.} There is no theorem to prove: these are definitions. The only claim made in the dictionary, that every weak
submodule of an ordinary module is graded, follows because a weak submodule $P$ is stable under $L_0^M=\omega^M_{(1)}$; for
$x=\sum_\alpha x_\alpha\in P$ the vectors $(L_0^M)^kx = \sum_\alpha\alpha^kx_\alpha$ lie in $P$, and a Vandermonde argument on the
finitely many nonzero components shows that each $x_\alpha\in P$. \qed

\begin{status}
Definitions; axiom-clean (\texttt{\#print axioms} of \lean{VOA.OrdVModule}, \lean{VOA.OrdVModule.IsIrreducible} and
\lean{VOA.OrdVModule.IsCompletelyReducible} is \lean{[propext, Classical.}\allowbreak\lean{choice, Quot.}\allowbreak\lean{sound]}, run for this chapter).
The global lower bound is a recorded design choice (moderate); see the warning in Section~\ref{l12:sec:categories}.
\end{status}

\begin{definition}[Adjoint module]\label{l12:def:adjoint}
Every VOA $V$ is a module over itself, the \emph{adjoint module}, with $Y_V = Y$ and weight spaces $V_\alpha = V_n$ if $\alpha = n\in\Z$ and
$V_\alpha = 0$ if $\alpha\notin\Z$.
\end{definition}
\noindent File \leanpath{VOA/\allowbreak Modules/\allowbreak WeakModule.lean}.
\begin{leancode}
def adjWeight (W : VertexOperatorAlgebra 𝒱) (α : ℂ) : Submodule ℂ 𝒱 :=
  ⨆ (n : ℤ) (_ : (n : ℂ) = α), W.grading n

def adjointWeakModule (W : VertexOperatorAlgebra 𝒱) : WeakModule W 𝒱 where
  YM := W.Y
  vacuum_YM := W.vacuum_Y
  truncation := fun u x => W.regularity u x
  jacobi := W.jacobi

def adjointModule (W : VertexOperatorAlgebra 𝒱) : OrdVModule W 𝒱 where
  toWeakModule := adjointWeakModule W
  weight := adjWeight W
\end{leancode}
(\lean{VOA.adjWeight}, \lean{VOA.adjointWeakModule}, \lean{VOA.adjointModule}. The remaining fields of \lean{adjointModule}
(\lean{isInternal}, \lean{L0_weight}, \lean{finrank}, \lean{boundedBelow}) are proofs and are omitted.)

\begin{dictionary}
\lean{adjWeight W α} is the supremum, over all integers $n$ with $(n:\C)=\alpha$, of $V_n$: it equals $V_n$ at $\alpha=n$ and $0$ at every
non-integer $\alpha$ (the library proves \lean{adjWeight_cast} and \lean{adjWeight_eq_bot}). This re-indexes the integer grading of $V$ by the complex
weights that an ordinary module requires. Here the one-line proof terms are informative: the Jacobi axiom of the adjoint weak module is
\emph{literally} the theorem \lean{W.jacobi} (Theorem~\ref{l03:thm:borcherds}), the truncation axiom is VA5 (regularity), and
\lean{vacuum_YM} is the vacuum axiom $Y(\ltwelvevac,z)=\mathrm{id}$. The fields of \lean{adjointModule} that were omitted use the VOA data:
\lean{isInternal} is the direct-sum decomposition of $V$, \lean{L0_weight} is $L_0v=nv$ on $V_n$, \lean{finrank} is $\dim V_n<\infty$, and
\lean{boundedBelow} follows from $V_n=0$ for $n<0$ (CFT type). So the adjoint module is the honest one, and
$u^V_{(n)}x = u_{(n)}x$.
\end{dictionary}

\begin{physics}
The adjoint module is the algebra acting on itself: the vacuum sector, whose states are those created from the vacuum by the fields of the
algebra itself (the state--field correspondence). It is the identity sector among the superselection sectors (the vacuum representation),
and in the fusion product it plays the role of the unit: $V\times M = M$.
\end{physics}

\textbf{Proof on paper.} (That $V$ is an ordinary module.) (M1) is the vacuum axiom of $V$, (M2) is regularity (VA5), (M3) is the Borcherds identity
of Theorem~\ref{l03:thm:borcherds}, for which $Y_M=Y$ makes both sides literally the same. $L_0^V=\omega_{(1)}$ acts as $n$ on $V_n$ by the
axiom (VOA2), the $V_n$ are finite-dimensional by the VOA axioms, $V_n=0$ for $n<0$ gives the lower bound $B=0$, and the sum is direct. \qed

\begin{status}
Unconditional. \texttt{\#print axioms} of \lean{VOA.adjointWeakModule} and \lean{VOA.adjointModule} is
\lean{[propext, Classical.}\allowbreak\lean{choice, Quot.}\allowbreak\lean{sound]} (run for this chapter).
\end{status}

Ordinary modules (Lean \lean{OrdVModule}) were defined in Definition~\ref{l12:def:ordmodule}: a weak
module with $M = \bigoplus_{\alpha\in\C}M_\alpha$, $L_0 = \alpha$ on $M_\alpha$, $\dim M_\alpha<\infty$,
and $\exists B\in\mathbb R$ with $\mathrm{Re}\,\alpha\ge B$ whenever $M_\alpha\neq0$. We add one
fidelity remark that Chapter~\ref{ch:01} did not make.

\begin{warning}
\textbf{The lower bound is global, not per coset (moderate, design choice).} Frenkel--Huang--Lepowsky
and Lepowsky--Li require, for a module with complex weights, only that \emph{for each} $\alpha$ one has
$M_{\alpha+n} = 0$ for all sufficiently negative integers $n$ (``truncation in each coset of $\Z$''). The
library's field \lean{boundedBelow} demands one real bound $B$ for \emph{all} weights. The two agree for
modules whose weights lie in finitely many cosets $\alpha+\Z$ --- in particular for every irreducible
module and every finite direct sum of them --- but an infinite direct sum
$\bigoplus_{j\ge1}M^{(j)}$ of irreducibles whose lowest weights $h_j$ lie in distinct cosets with
$\mathrm{Re}\,h_j\to-\infty$ is an FHL module and \emph{not} an \lean{OrdVModule}. Consequences: a
statement ``every \lean{OrdVModule} is completely reducible'' (Gannon's clause (1) below) quantifies over
\emph{fewer} modules than the FHL statement, hence is formally weaker. The library records the same
choice for logarithmic modules (docstring of \leanpath{VOA/\allowbreak Modules/\allowbreak LogVModule.lean}); for
ordinary modules it is not mentioned. No theorem in this chapter depends on the difference.
\end{warning}

\begin{lemma}[Modes shift weights; lowering modes kill the bottom]\label{l12:lem:modeweight}
Let $M$ be an ordinary module, $v\in V_\beta$, $x\in M_\alpha$, $n\in\Z$. Then
$v^M_{(n)}x\in M_{\alpha+\beta-n-1}$. If moreover every weight of $M$ lies in $h+\N$, $x\in M_h$ and
$n\ge\beta$, then $v^M_{(n)}x = 0$.
\end{lemma}
\noindent File \leanpath{VOA/\allowbreak Modules/\allowbreak ZeroModeMult.lean}.
\begin{leancode}
theorem mode_mem_weight (N : OrdVModule W M) {β : ℤ} {v : 𝒱} (hv : v ∈ W.grading β)
    {α : ℂ} {x : M} (hx : x ∈ N.weight α) (n : ℤ) :
    mode (N.YM v) n x ∈ N.weight (α + (β : ℂ) - n - 1) := by

theorem mode_lowering_vanish (N : OrdVModule W M) {h : ℂ} (hcw : IsLowestWeight N h)
    {β : ℤ} {v : 𝒱} (hv : v ∈ W.grading β) {x : M} (hx : x ∈ N.weight h) {n : ℤ}
    (hn : β ≤ n) : mode (N.YM v) n x = 0 := by
\end{leancode}

\begin{dictionary}
The first statement is the commutator $[L_0, v_{(n)}] = (\beta-n-1)v_{(n)}$ on $M$ in eigenvector form;
the second uses \lean{IsLowestWeight N h} (Definition~\ref{l12:def:lowest} below): the target weight
$h+\beta-n-1$ is not of the form $h+k$ with $k\in\N$ when $n\ge\beta$, so the target space is zero.
Same as the textbook (it is the standard consequence of the commutator formula with $\omega$; see Lepowsky--Li, §4.1, and the proof of Zhu 1996, Theorem~2.1.2).
\end{dictionary}

\textbf{Proof on paper.} The module commutator formula (Theorem~\ref{l12:thm:modcomm}) with $u=\omega$,
$m = 1$ gives $[L_0,v_{(n)}] = \sum_i\binom1i(\omega_{(i)}v)_{(n+1-i)} = (L_{-1}v)_{(n+1)} + (L_0v)_{(n)}
= (-(n+1)+\beta)v_{(n)}$ by Theorem~\ref{l12:thm:covariance}. Hence $L_0v_{(n)}x = (\alpha+\beta-n-1)v_{(n)}x$,
and since $L_0$ is diagonal with eigenspaces $M_\gamma$, an eigenvector of eigenvalue $\gamma$ lies in
$M_\gamma$ (library lemma \lean{eigenvector_mem_weight}). The vanishing statement follows. \qed

\begin{status}
Proved, no hypotheses beyond those displayed; axiom-clean (part of the dependency cone of
\lean{oMod_zhuMul_lowest_all}, whose \texttt{\#print axioms} we ran).
\end{status}

\begin{physics}
A mode $v_{(n)}$ of a field of weight $\beta$ changes the energy by $\beta-n-1$: modes with $n<\beta-1$ create
excitations, $n=\beta-1$ is the zero mode (a conserved charge on each energy level) and $n>\beta-1$ annihilate. On
the ground states of a sector, all annihilation modes vanish: this is the defining property of a \emph{primary}
state, and it is what makes Zhu's algebra act on the ground states alone.
\end{physics}

\subsection{Logarithmic modules}

In logarithmic conformal field theory $L_0$ is not diagonalisable on some sectors: it has Jordan
blocks, and correlation functions acquire logarithms $\log z$. The appropriate module notion is the
\emph{grading-restricted generalised module} of Huang--Lepowsky--Zhang.

\begin{definition}[Logarithmic module; Huang--Lepowsky--Zhang; McRae Def.~2.4]\label{l12:def:log}
A \emph{logarithmic} $V$-module is a weak module $M$ with $M = \bigoplus_{\alpha\in\C}M_{[\alpha]}$ such
that every $x\in M_{[\alpha]}$ is annihilated by some power of $L_0-\alpha$, every $M_{[\alpha]}$ is
finite-dimensional, and $\mathrm{Re}\,\alpha$ is bounded below on the nonzero pieces.
\end{definition}
\noindent File \leanpath{VOA/\allowbreak Modules/\allowbreak LogVModule.lean}.
\begin{leancode}
structure LogVModule (W : VertexOperatorAlgebra 𝒱) (M : Type v)
    [AddCommGroup M] [Module ℂ M] extends WeakModule W M where
  weight : ℂ → Submodule ℂ M
  isInternal : DirectSum.IsInternal weight
  weight_le_maxGenEigenspace :
    ∀ α : ℂ, weight α ≤ (mode (YM W.conformalVector) 1).maxGenEigenspace α
  finrank : ∀ α, FiniteDimensional ℂ (weight α)
  boundedBelow : ∃ B : ℝ, ∀ α : ℂ, weight α ≠ ⊥ → B ≤ α.re
\end{leancode}

\begin{dictionary}
\lean{mode (YM W.conformalVector) 1} is $L_0^M = \omega^M_{(1)}$. For an endomorphism $T$,
\lean{T.maxGenEigenspace α} is Mathlib's \emph{generalised eigenspace}
$\{x : (T-\alpha)^kx = 0\text{ for some }k\}$. The only difference from \lean{OrdVModule} is that
\lean{weight_le_maxGenEigenspace} replaces ``$L_0 = \alpha$ on $M_\alpha$''. Only one inclusion is
demanded; the library proves the other (Theorem~\ref{l12:thm:loggrading}), so the grading is not extra
data. The lower bound is global, with the same caveat as for ordinary modules (the library's docstring
records it): HLZ and McRae truncate per coset.

\textbf{Same as the literature?} Same as HLZ's grading-restricted generalised module and McRae's
Definition~2.4, up to the global-versus-per-coset lower bound. The Jacobi identity is required in its
usual form; the $L_{-1}$-derivative property, which HLZ list, is automatic
(Theorem~\ref{l12:thm:covariance}).

\textbf{Could it be vacuous, or collapse to ordinary modules?} Every ordinary module is logarithmic
(\lean{OrdVModule.toLog}), so the class is non-empty. Whether the library contains a module with
$L_0^{\mathrm{nil}}\neq0$ is the subject of the triplet chapter (Chapter~\ref{ch:L18}); in this chapter we
only use the definition.
\end{dictionary}

\begin{physics}
In a logarithmic CFT (the triplet models $W(p)$, symplectic fermions, critical percolation and polymers) the
Hamiltonian $L_0$ has Jordan blocks: there are pairs of states $(\phi,\psi)$ with $L_0\phi = h\phi$ and
$L_0\psi = h\psi+\phi$. The two-point function then contains $\log z$, and the partition function must be
completed by \emph{pseudo-traces}. A logarithmic module is a sector of such a theory; the generalised
eigenspace $M_{[h]}$ collects all states whose energy is $h$ ``up to nilpotent mixing''.
\end{physics}

\begin{theorem}[The logarithmic grading is canonical; Jordan decomposition of $L_0$]\label{l12:thm:loggrading}
For a logarithmic module $M$: (i) $M_{[\alpha]}$ equals the full generalised $\alpha$-eigenspace of $L_0$;
(ii) writing $L_0^{\mathrm{ss}}$ for the operator acting as $\alpha$ on $M_{[\alpha]}$ and
$L_0^{\mathrm{nil}} = L_0 - L_0^{\mathrm{ss}}$, the two commute, $L_0^{\mathrm{nil}}$ is locally nilpotent,
and $M$ is an ordinary module (with the same grading) if and only if $L_0^{\mathrm{nil}} = 0$.
\end{theorem}
\noindent File \leanpath{VOA/\allowbreak Modules/\allowbreak LogVModule.lean}.
\begin{leancode}
theorem weight_eq_maxGenEigenspace (N : LogVModule W M) (α : ℂ) :
    N.weight α = N.L0op.maxGenEigenspace α :=

noncomputable def L0ss (N : LogVModule W M) : Module.End ℂ M :=
  letI := DirectSum.IsInternal.chooseDecomposition N.weight N.isInternal
  (DirectSum.toModule ℂ ℂ M (fun β : ℂ => β • (N.weight β).subtype)) ∘ₗ
    (DirectSum.decomposeLinearEquiv N.weight).toLinearMap

noncomputable def L0nil (N : LogVModule W M) : Module.End ℂ M := N.L0op - N.L0ss

theorem L0ss_comm_L0nil (N : LogVModule W M) : N.L0ss * N.L0nil = N.L0nil * N.L0ss := by

theorem L0nil_locNilpotent (N : LogVModule W M) {α : ℂ} {x : M} (hx : x ∈ N.weight α) :
    ∃ k : ℕ, (N.L0nil ^ k) x = 0 := by

def IsL0Semisimple (N : LogVModule W M) : Prop :=
  ∀ (α : ℂ) (x : M), x ∈ N.weight α → N.L0op x = α • x

theorem isL0Semisimple_iff_L0nil_eq_zero (N : LogVModule W M) :
    N.IsL0Semisimple ↔ N.L0nil = 0 := by

def LogVModule.toOrd (N : LogVModule W M) (h : N.IsL0Semisimple) : OrdVModule W M where
\end{leancode}

\begin{dictionary}
\lean{N.L0op} is $L_0^M$. \lean{L0ss} is defined by decomposing $x = \sum_\beta x_\beta$ (the internal direct
sum gives the components, by choice, hence \lean{noncomputable}) and returning $\sum_\beta\beta x_\beta$: this
is $L_0^{\mathrm{ss}}$. \lean{*} on \lean{Module.End} is composition. \lean{toOrd} turns an
$L_0$-semisimple logarithmic module into an ordinary one with the same weak module and grading, and the
library proves the round trips \lean{toOrd_toLog}, \lean{toLog_toOrd} by \lean{rfl}. So ordinary modules
are \emph{exactly} the logarithmic modules with $L_0^{\mathrm{nil}} = 0$; this is the intended
statement.
\end{dictionary}

\textbf{Proof on paper.} (i) If $y$ is a generalised $\alpha$-eigenvector, decompose $y = \sum_\beta y_\beta$
along the grading; $y - y_\alpha$ is both a generalised $\alpha$-eigenvector and a sum of generalised
$\beta$-eigenvectors with $\beta\neq\alpha$. Generalised eigenspaces for distinct eigenvalues are
independent, so $y = y_\alpha\in M_{[\alpha]}$. (ii) $L_0$ preserves each generalised eigenspace and
$L_0^{\mathrm{ss}}$ is a scalar on it, so they commute; on $M_{[\alpha]}$, $L_0^{\mathrm{nil}} = L_0-\alpha$,
which is locally nilpotent there; and $L_0^{\mathrm{nil}}=0$ says precisely that $L_0 = \alpha$ on each
$M_{[\alpha]}$. \qed

\begin{status}
Proved, no hypotheses. \texttt{\#print axioms} of \lean{VOA.LogVModule.weight_eq_maxGenEigenspace} and
\lean{VOA.LogVModule.isL0Semisimple_iff_L0nil_eq_zero}: \lean{[propext, Classical.choice, Quot.sound]}.
\end{status}

\begin{physics}
$L_0 = L_0^{\mathrm{ss}} + L_0^{\mathrm{nil}}$ is the Jordan decomposition of the Hamiltonian. The torus
amplitude of a logarithmic sector is $\tr\, q^{L_0^{\mathrm{ss}}-c/24}e^{2\pi i\tau L_0^{\mathrm{nil}}}$-type
expressions; the nilpotent part is what produces the factors of $\tau$ (equivalently $\log q$) in logarithmic
characters, and its vanishing is exactly the statement that the sector is an ordinary, non-logarithmic one.
\end{physics}

\subsection{Every ordinary module is admissible}

An ordinary module comes with a $\C$-grading, not an $\N$-grading. To compare rationality
conditions over the two categories one needs to manufacture an admissible grading on the same weak
module.

\begin{theorem}[Ordinary $\subseteq$ admissible; Dong--Li--Mason, Lemma~3.4 at $g=1$]\label{l12:thm:ordadm}
Every ordinary module $M$ admits an $\N$-grading $M = \bigoplus_kM(k)$, on the same weak module,
satisfying the admissibility axiom of Definition~\ref{l12:def:admissible}.
\end{theorem}
\noindent Files \leanpath{VOA/\allowbreak Modules/\allowbreak AdmissibleRationality.lean} (the predicate) and
\leanpath{VOA/\allowbreak Modules/\allowbreak RationalityLattice.lean} (the theorem).
\begin{leancode}
def OrdVModule.HasAdmissibleRefinement {W : VertexOperatorAlgebra 𝒱} {M : Type v}
    [AddCommGroup M] [Module ℂ M] (N : OrdVModule W M) : Prop :=
  ∃ lev : ℤ → Submodule ℂ M,
    (∀ k : ℤ, k < 0 → lev k = ⊥) ∧ DirectSum.IsInternal lev ∧
      ∀ ⦃w p k : ℤ⦄ ⦃a : 𝒱⦄, a ∈ W.grading w → ∀ ⦃x : M⦄, x ∈ lev k →
        mode (N.YM a) p x ∈ lev (k + (w - p - 1))

theorem OrdVModule.hasAdmissibleRefinement {W : VertexOperatorAlgebra 𝒱} {M : Type u}
    [AddCommGroup M] [Module ℂ M] (N : OrdVModule W M) : N.HasAdmissibleRefinement := by
\end{leancode}

\begin{dictionary}
\lean{HasAdmissibleRefinement} lists exactly the three admissible-module fields for the weak module
\lean{N.toWeakModule}; it is stated as a proposition about the same weak module so that the two notions of
complete reducibility are literally about the same subspaces. The carrier is required to live in the
universe \lean{Type u} of $V$ (a bookkeeping restriction of Lean's type theory; every module of interest is
there). Same as DLM Lemma~3.4 (untwisted case).
\end{dictionary}

\textbf{Proof on paper.} Let $B$ be the global lower bound. Put $M(k) = \bigoplus\{M_\alpha :
\lfloor\mathrm{Re}\,\alpha - B\rfloor = k\}$. Nonzero pieces have $\mathrm{Re}\,\alpha\ge B$, so $M(k) = 0$ for
$k<0$. Regrouping an internal direct sum along any map of index sets is again an internal direct sum. If
$a\in V_w$ and $x\in M_\alpha$, then $a_{(p)}x\in M_{\alpha+w-p-1}$ (Lemma~\ref{l12:lem:modeweight}), and since
$w-p-1$ is an \emph{integer}, $\lfloor\mathrm{Re}(\alpha+w-p-1)-B\rfloor = \lfloor\mathrm{Re}\,\alpha-B\rfloor
+ w-p-1$. \qed

(This is where the \emph{global} bound is convenient: with a per-coset bound one would choose a base point
in each coset of $\C/\Z$.) A related, weaker statement is in the same file:
\lean{OrdVModule.}\allowbreak\lean{hasAdmissibleRefinement_}\allowbreak\lean{of_}\allowbreak\lean{isLowestWeight}, which takes $M(k) = M_{h+k}$ when all weights
lie in $h+\N$.

\begin{status}
Proved, no hypotheses; \texttt{\#print axioms VOA.OrdVModule.hasAdmissibleRefinement} =
\lean{[propext, Classical.choice, Quot.sound]}. The library also packages it as a definition
\lean{OrdVModule.toAdmissible : AdmissibleModule W M}.
\end{status}

\begin{physics}
Every physical sector, with its energy grading, can be viewed as a sector with ``levels'': group the energy
eigenspaces by the integer part of their energy above the vacuum floor. Since fields shift energies by integers,
they shift levels by the same integers. This is why statements about all admissible sectors (rationality)
automatically say something about all ordinary sectors.
\end{physics}

\subsection{Submodules, morphisms, exact sequences}

Three small pieces of infrastructure are used later.
\begin{itemize}
\item \textbf{Submodules of ordinary modules are graded.} If $P$ is a weak submodule of an ordinary module
  $M$, then $P = \bigoplus_\alpha(P\cap M_\alpha)$ and $P$ is an ordinary module with these weight spaces
  (\lean{OrdVModule.submodule} in \leanpath{VOA/\allowbreak Modules/\allowbreak OrdSubmodule.lean}, proved,
  axiom-clean). The reason: $P$ is $L_0$-stable, and an $L_0$-stable subspace of a space on which $L_0$
  is diagonalisable with given eigenspaces contains the eigencomponents of each of its vectors (Vandermonde).
  Consequently, for \emph{ordinary} modules ``irreducible'' and ``completely reducible'' mean the same
  whether one uses weak or graded submodules --- the subtlety of the warning in Section~\ref{l12:warn:admcr} below is
  specific to admissible modules.
\item \textbf{Morphisms and short exact sequences} are the subject of the next proposition.
\end{itemize}

\begin{proposition}[Graded dimension is additive on short exact sequences]\label{l12:prop:ses}
If $0\to A\xrightarrow{f}B\xrightarrow{g}C\to0$ is a short exact sequence of ordinary modules with
weight-preserving module maps, then $\dim B_\alpha = \dim A_\alpha + \dim C_\alpha$ for every $\alpha$, i.e.\
$\chi_B = \chi_A+\chi_C$ as formal $q$-series.
\end{proposition}
\noindent File \leanpath{VOA/\allowbreak Modules/\allowbreak OrdinaryExact.lean}.
\begin{leancode}
structure Hom (A : OrdVModule W M) (B : OrdVModule W N) where
  toWeakHom : A.toWeakModule.Hom B.toWeakModule
  map_weight : ∀ (α : ℂ) {m : M}, m ∈ A.weight α →
    toWeakHom.toLinearMap m ∈ B.weight α

structure ShortExact (f : Hom A B) (g : Hom B C) : Prop where
  injective : Function.Injective f.toLinearMap
  exactAt : ExactAt f g
  surjective : Function.Surjective g.toLinearMap

theorem gradedDim_eq_add (h : ShortExact f g) :
    gradedDim B = gradedDim A + gradedDim C := by
\end{leancode}

\begin{dictionary}
\lean{ExactAt f g} is ``range $f$ = kernel $g$'' pointwise (\lean{exactAt_iff_range_eq_ker}).
\lean{gradedDim} is the formal character (Definition~\ref{l12:def:char}). Weight preservation is explicit
data in \lean{Hom}; it is in fact automatic for module maps between ordinary modules (they commute with
$L_0$; compare \lean{ModuleEquiv.maps_weight} below), so carrying it costs nothing. Same as the standard
statement. The companion file \leanpath{VOA/\allowbreak Modules/\allowbreak OrdinaryComplex.lean} develops
cochain complexes of ordinary modules, weight homology, and a coefficientwise Euler--Poincar\'e identity
(\lean{EulerWindowSupport.}\allowbreak\lean{eulerCharacter_}\allowbreak\lean{eq}); exactness of any concrete resolution is \emph{not} claimed there.
\end{dictionary}

\textbf{Proof on paper.} Restricted to weight $\alpha$ the sequence is a short exact sequence of
finite-dimensional vector spaces (weight preservation and the internal decompositions); apply rank--nullity.
\qed

\begin{status}
Proved; \texttt{\#print axioms VOA.OrdVModule.Hom.ShortExact.gradedDim\_eq\_add} = \lean{[propext,
Classical.choice, Quot.sound]}.
\end{status}

\begin{physics}
Characters count states level by level, so they are additive in exact sequences: a sector built as an extension
of two others has the sum of their partition functions. This is why characters cannot detect whether an
extension splits, and why in logarithmic and non-rational theories the characters of irreducibles must be
supplemented by finer data (pseudo-traces, or Euler characteristics of resolutions such as Felder's).
\end{physics}

\subsection{Restriction to a sub-VOA}

\begin{definition}[Restriction; Gannon §5.3]\label{l12:def:restrict}
Let $S\subseteq V$ be a sub-VOA (containing $\ltwelvevac$ and the conformal vector $\omega$). Any weak
(resp.\ ordinary) $V$-module $M$ becomes a weak (resp.\ ordinary) $S$-module by restricting $Y_M$ to $S$,
with the same weight spaces.
\end{definition}
\noindent Files \leanpath{VOA/\allowbreak Modules/\allowbreak ModuleRestrict.lean},
\leanpath{VOA/\allowbreak Modules/\allowbreak ModuleRestrictOrd.lean}.
\begin{leancode}
noncomputable def restrictWeakModule {M : Type v} [AddCommGroup M] [Module ℂ M]
    (A : WeakModule W M) : WeakModule S.toVOA M where
  YM := A.YM.comp S.carrier.subtype
  vacuum_YM := A.vacuum_YM
  truncation := fun u x => A.truncation (u : 𝒱) x
  jacobi := fun u v x ℓ m n => A.jacobi (u : 𝒱) (v : 𝒱) x ℓ m n

noncomputable def restrictOrdVModule {M : Type v} [AddCommGroup M] [Module ℂ M]
    (N : OrdVModule W M) : OrdVModule S.toVOA M where
  toWeakModule := S.restrictWeakModule N.toWeakModule
  weight := N.weight
  isInternal := N.isInternal
  L0_weight := N.L0_weight
  finrank := N.finrank
  boundedBelow := N.boundedBelow
\end{leancode}

\begin{dictionary}
These are full definitions, and every field is informative: \lean{A.YM.comp S.carrier.subtype} is
$Y_M\circ(S\hookrightarrow V)$; each axiom is the $V$-axiom evaluated on elements of $S$ (written
\lean{(u : 𝒱)}, the coercion of $u\in S$ into $V$). The Jacobi identity type-checks without change because the
product of $S$ is that of $V$. The $L_0$ field transfers because $S$ has the \emph{same} conformal vector as $V$
(a field of \lean{VOASubalgebra}). So this is restriction along a \emph{conformal} sub-VOA, as in the
literature; restriction along a vertex subalgebra with a different conformal vector (a coset construction) is
not covered. The library also provides \lean{adjointRestrict}, $V$ as an $S$-module, and the inclusion
$S\hookrightarrow V$ as a module map (\lean{adjointInclusion}).
\end{dictionary}

\begin{status}
Definitions with all proofs inlined above; \texttt{\#print axioms} of
\lean{VOA.VOASubalgebra.restrictWeakModule} and \lean{...restrictOrdVModule}: \lean{[propext,
Classical.choice, Quot.sound]}.
\end{status}

\begin{physics}
Restriction is ``forgetting symmetry'': a sector of a CFT with chiral algebra $V$ is in particular a
representation of any smaller chiral algebra $S$ with the same stress tensor. For the lattice VOA $V_L$ and its
Heisenberg subalgebra, the vacuum sector of $V_L$ decomposes into the charged Fock sectors $F_\alpha$, $\alpha\in L$
--- the physicist's statement that the compactified boson's Hilbert space is a sum of momentum/winding sectors of
the free boson. For an orbifold $V^G\subseteq V$, restriction gives the untwisted sectors of the orbifold theory.
\end{physics}

%======================================================================
\section{Characters, zero modes, and lowest weights}\label{l12:sec:char}
%======================================================================

\begin{definition}[Character, zero mode, one-point function; Zhu 1996, §4]\label{l12:def:char}
Let $M$ be an ordinary module and $c$ the central charge of $V$.
\begin{enumerate}
\item The \emph{character} (graded dimension) is the formal series
  $\chi_M(q) = \tr_M q^{L_0-c/24} = \sum_\alpha(\dim M_\alpha)\,q^{\alpha-c/24}$, recorded as the function
  $e\mapsto\dim M_{e+c/24}$ (coefficient of $q^e$).
\item The \emph{zero mode} of $v\in V_k$ is $o(v) = v^M_{(k-1)}$, extended linearly to all of $V$.
\item The \emph{one-point function} is $\chi_M(v;q) = \tr_M o(v)q^{L_0-c/24}$, the function
  $e\mapsto\tr_{M_{e+c/24}}o(v)$.
\end{enumerate}
\end{definition}
\noindent File \leanpath{VOA/\allowbreak Modules/\allowbreak Character.lean}.
\begin{leancode}
abbrev FormalQSeries : Type := ℂ → ℂ

noncomputable def weightBlock (N : OrdVModule W M) (T : End ℂ M) (α : ℂ) : End ℂ (N.weight α) :=
  (internalProjTo N.isInternal α).comp (T.comp (N.weight α).subtype)

noncomputable def gradedDim (N : OrdVModule W M) : FormalQSeries :=
  fun e => (Module.finrank ℂ (N.weight (e + W.centralCharge / 24)) : ℂ)

noncomputable def onePoint (N : OrdVModule W M) (v : 𝒱) : FormalQSeries :=
  fun e =>
    haveI := N.finrank (e + W.centralCharge / 24)
    LinearMap.trace ℂ _ (weightBlock N (oMod N v) (e + W.centralCharge / 24))

theorem onePoint_vacuum (N : OrdVModule W M) : onePoint N W.vacuum = gradedDim N := by
\end{leancode}

\begin{dictionary}
A \lean{FormalQSeries} is any function $\C\to\C$, read as the formal sum $\sum_e f(e)q^e$ over complex
exponents; no convergence is asserted (that is the business of Chapters~\ref{ch:L14}--\ref{ch:L15}).
\lean{weightBlock N T α} is the diagonal block $M_\alpha\to M\to M_\alpha$ of an operator $T$ (inclusion,
then $T$, then projection along the direct sum). \lean{LinearMap.trace} is the trace, well defined because
$M_\alpha$ is finite-dimensional (the \lean{haveI} supplies that fact). \lean{oMod N} is the zero mode of
Definition~\ref{ch12:def:omod}. \lean{onePoint_vacuum} says $\chi_M(\ltwelvevac;q) = \chi_M(q)$
because $o(\ltwelvevac) = \mathrm{id}$. The $q^{-c/24}$ shift is built into the indexing. Same as Zhu's
definitions.
\end{dictionary}

\begin{status}
Definitions; \texttt{\#print axioms} of the lemmas \lean{onePoint_vacuum}, \lean{onePoint_add} is the
standard three (file docstring).
\end{status}

\begin{physics}
$\chi_M(q)$ is the torus partition function of the chiral sector $M$, $\tr\,q^{L_0-c/24}$ with $q = e^{2\pi i\tau}$;
the shift $-c/24$ is the Casimir energy of the cylinder. The one-point function is the torus one-point
function of the zero mode of a field --- the expectation value of a conserved charge, weighted by the
Boltzmann factor. Zhu's theorem says that for a rational $C_2$-cofinite theory these span a representation of
$SL_2(\Z)$, which is the mathematical form of modular invariance of the torus amplitudes. (The library proves it in
Chapter~\ref{ch:L15} for one-point functions of $C_2$-cofinite $V$ that are rational in its sense
\lean{AdmissiblyRational}, and more generally whenever $A(V)$ is semisimple.)
\end{physics}

\begin{definition}[Lowest weight; conformal weight]\label{l12:def:lowest}
An ordinary module $M$ \emph{has lowest weight} $h\in\C$ (library sense) if every $\alpha$ with
$M_\alpha\neq0$ is of the form $h+n$ with $n\in\N$. It \emph{has conformal weight} $h$ if moreover
$M_h\neq0$.
\end{definition}
\noindent Files \leanpath{VOA/\allowbreak Modules/\allowbreak ZeroModeMult.lean} and
\leanpath{VOA/\allowbreak Modules/\allowbreak Rationality.lean}.
\begin{leancode}
def IsLowestWeight (N : OrdVModule W M) (h : ℂ) : Prop :=
  ∀ α : ℂ, N.weight α ≠ ⊥ → ∃ n : ℕ, α = h + n

def HasConformalWeight (N : OrdVModule W M) (h : ℂ) : Prop :=
  N.weight h ≠ ⊥ ∧ ∀ α : ℂ, N.weight α ≠ ⊥ → ∃ n : ℕ, α = h + n
\end{leancode}

\begin{dictionary}
\lean{N.weight α ≠ ⊥} is ``$M_\alpha\neq0$''. The two predicates differ exactly by the conjunct
\lean{N.weight h ≠ ⊥}.
\end{dictionary}

\begin{warning}
\textbf{\lean{IsLowestWeight} does not require $M_h\neq0$ (confirmed; moderate).} This confirms, with a
machine check, a defect suspected earlier in the project. Precisely:
\begin{itemize}
\item the zero module satisfies \lean{IsLowestWeight N h} for \emph{every} $h$;
\item a module with lowest weight $h$ in the usual sense also satisfies \lean{IsLowestWeight N (h-1)},
  \lean{IsLowestWeight N (h-2)}, \ldots, and for those values $M_{h-k} = 0$.
\end{itemize}
Theorems that take \lean{hcw : IsLowestWeight N h} and conclude something about $M_h$ (Zhu's
multiplicativity $o(u*v) = o(u)o(v)$ on $M_h$, $O(V)M_h = 0$, the $A(V)$-module structure
\lean{moduleMh} on $M_h$) remain \emph{true}; when $M_h = 0$ they are statements about the zero space, i.e.\
vacuous. They are therefore no weaker as hypotheses, but a reader must not infer from them that $M_h$ is the
genuine bottom of $M$. Where the library needs a genuine bottom it adds the conjunct: \lean{HasConformalWeight}
above, and the field \lean{realized : isMod.weight h ≠ ⊥} of the bundled \lean{LWModule}
(Definition~\ref{l12:def:lwmodule}).

There is a second, opposite difference from the informal phrase ``$M$ has lowest weight $h$'': the library
predicate requires \emph{all} weights to lie in the single coset $h+\N$. A direct sum of two irreducible modules
whose lowest weights differ by a non-integer satisfies \lean{IsLowestWeight} for no $h$ at all. In the literature
Zhu's algebra acts on the ``top level'' $\Omega(M)$ of any admissible module, coset by coset; the library's
statements cover modules with one coset of weights. For irreducible modules, the case that matters for Zhu's
correspondence, all weights lie in one coset, so nothing is lost there.
\end{warning}

\begin{physics}
The lowest weight $h$ is the conformal dimension of the primary field that creates the sector from the vacuum:
the ground-state energy of the sector is $h-c/24$ on the cylinder. ``Weights in $h+\N$'' says that all other
states are descendants with integer-spaced energies above the ground state, which is the spectrum of a single
irreducible representation (one primary and its descendants). A sector made of several primaries with
non-integer energy differences is a direct sum of representations and has no single ``$h$''.
\end{physics}

\begin{theorem}[$T$-periodicity of the character; Gannon Lemma~5.3.3]\label{l12:thm:Tperiod}
If $M$ has conformal weight $h$ then $\chi_M(\tau+1) = e^{2\pi i(h-c/24)}\chi_M(\tau)$, i.e.\ the coefficient
of $q^e$ in $\chi_M$ is multiplied by $e^{2\pi ie}$ and this factor equals $e^{2\pi i(h-c/24)}$ on every
nonzero coefficient.
\end{theorem}
\noindent File \leanpath{VOA/\allowbreak Modules/\allowbreak Rationality.lean}.
\begin{leancode}
noncomputable def qShift (χ : FormalQSeries) : FormalQSeries :=
  fun e => Complex.exp (2 * (Real.pi : ℂ) * Complex.I * e) * χ e

theorem char_T_transform (N : OrdVModule W M) (h : ℂ) (hcw : HasConformalWeight N h) :
    qShift (gradedDim N)
      = (Complex.exp (2 * (Real.pi : ℂ) * Complex.I * (h - W.centralCharge / 24)))
          • gradedDim N := by

theorem conformalWeight_char_periodicity (N : OrdVModule W M)
    {h : ℂ} (hcw : HasConformalWeight N h) :
\end{leancode}

\begin{dictionary}
\lean{qShift} implements $\tau\mapsto\tau+1$ on a formal series in $q = e^{2\pi i\tau}$: $q^e\mapsto e^{2\pi ie}q^e$.
The only hypothesis of \lean{char_T_transform} is \lean{HasConformalWeight}; in fact only its second clause
(weights in $h+\N$) is used. Same as Gannon's lemma.

\textbf{Decorative hypotheses, since removed.} The variant
\lean{conformalWeight_}\allowbreak\lean{char_}\allowbreak\lean{periodicity} used to carry \lean{_hwr : GannonRational W} and
\lean{_hN : OrdVModule.IsIrreducible N} (the leading underscore is Lean's convention for an unused argument) so
that it read like Gannon's lemma. The library has dropped them; the current statement, quoted above, has only
\lean{hcw}. The theorem is \emph{not} conditional on rationality or irreducibility.
\end{dictionary}

\textbf{Proof on paper.} If $M_{e+c/24}\neq0$ then $e+c/24 = h+n$ with $n\in\N$, so
$e^{2\pi ie} = e^{2\pi i(h-c/24)}e^{2\pi in} = e^{2\pi i(h-c/24)}$; where the coefficient is $0$ both sides vanish.
\qed

\begin{status}
Proved; \texttt{\#print axioms VOA.char\_T\_transform} = \lean{[propext, Classical.choice, Quot.sound]}.
The \emph{rationality} of $h$ (Anderson--Moore, Dong--Li--Mason) is not proved in this file; its docstring calls
it a ``gate'' to be supplied by modularity, and nothing is claimed.
\end{status}

\begin{physics}
$\tau\mapsto\tau+1$ is the Dehn twist of the torus, generated by $e^{2\pi i(L_0-c/24)}$. On a sector of
conformal weight $h$ it acts by the phase $e^{2\pi i(h-c/24)}$, the eigenvalue of the modular $T$-matrix; the
spin $h$ mod $1$ of the primary field is what this phase measures.
\end{physics}

%======================================================================
\section{Zhu's algebra}\label{l12:sec:zhu}
%======================================================================

Recall (Definition~\ref{ch12:def:zhumul}) that for homogeneous $u\in V_k$,
\begin{equation*}
 u*v = \Res_z Y(u,z)v\frac{(1+z)^k}{z} = \sum_{m=0}^k\binom km u_{(m-1)}v,
\end{equation*}
extended bilinearly; $O(V)$ is the span of all $((L_{-1}+L_0)u)*v$, equivalently (Chapter~\ref{ch:01}) of all
$u\circ v = \Res_zY(u,z)v(1+z)^{\wt u}/z^2$; and $A(V) = V/O(V)$. Chapter~\ref{ch:01} proved that
$o(u*v) = o(u)o(v)$ on a lowest-weight space and that $O(V)$ acts there by zero
(Theorem~\ref{l12:thm:zhuaction}). What remains is the ring structure itself.

\begin{theorem}[Zhu: zero modes act on the lowest-weight space through $A(V)$]\label{l12:thm:zhuaction}
Let $M$ be an ordinary $V$-module with lowest weight $h$. For all $u,v\in V$ and $x\in M_h$,
$o(u*v)x = o(u)o(v)x$. Moreover $o(O(V))M_h = 0$, so $u\mapsto o(u)|_{M_h}$ factors through $A(V) = V/O(V)$. For homogeneous
$a\in V_M$ ($M\in\N$) and all $b\in V$, all the residues $\sum_i\binom{M+n}{i}a_{(i-2-n)}b$, $n\ge0$, lie in $O(V)$.
\end{theorem}
\noindent File \leanpath{VOA/\allowbreak Modules/\allowbreak ZhuModule.lean}.
\begin{leancode}
theorem oMod_zhuMul_lowest_all (N : OrdVModule W Mod) {h : ℂ} (hcw : IsLowestWeight N h)
    (u v : 𝒱) {x : Mod} (hx : x ∈ N.weight h) :
    oMod N (zhuMul W u v) x = oMod N u (oMod N v x)

noncomputable def rhoLW (N : OrdVModule W Mod) (h : ℂ) : 𝒱 →ₗ[ℂ] End ℂ (N.weight h) where
  toFun u := weightBlock N (oMod N u) h

theorem O_le_ker_rhoLW (N : OrdVModule W Mod) {h : ℂ} (hcw : IsLowestWeight N h) :
    O W ≤ LinearMap.ker (rhoLW N h)

noncomputable def Pres (W : VertexOperatorAlgebra 𝒱) (M : ℕ) (a b : 𝒱) (n : ℕ) : 𝒱 :=
  ∑ i ∈ Finset.range (M + n + 1), cbz ((M : ℤ) + n) (i : ℤ) • W.smul a ((i : ℤ) - 2 - n) b

theorem Pres_mem_O (W : VertexOperatorAlgebra 𝒱) {M : ℕ} {a : 𝒱}
    (ha : a ∈ W.grading (M : ℤ)) (b : 𝒱) (n : ℕ) : Pres W M a b n ∈ O W
\end{leancode}

\begin{dictionary}
\lean{oMod N u} is the zero mode $o(u)$ of Definition~\ref{l12:def:char}(2) (extended linearly from homogeneous $u$); \lean{zhuMul W u v} is the
Zhu product $u*v$ and \lean{O W} the subspace $O(V)$ (Definition~\ref{ch12:def:zhumul}); \lean{IsLowestWeight N h} is
Definition~\ref{l12:def:lowest}; \lean{N.weight h} is $M_h$. \lean{rhoLW N h u} is the restriction of $o(u)$ to $M_h$:
\lean{weightBlock N T α} takes the diagonal block $M_\alpha\to M\to M_\alpha$ of an operator $T$ (the library proves, as a separate lemma, that $o(u)$
preserves $M_h$, so the block loses nothing). \lean{O W ≤ LinearMap.ker (rhoLW N h)} says every element of $O(V)$ acts as zero on $M_h$.
\lean{Pres W M a b n} is the residue $\Res_zY(a,z)b\,\frac{(1+z)^{M+n}}{z^{2+n}}$, written as a finite sum over $i\in\{0,\dots,M+n\}$
(the terms with larger $i$ vanish since $\binom{M+n}{i}=0$). \textbf{Hypotheses:} lowest weight $h$ for the first three statements, which
are quantified over \emph{all} $u,v\in V$ (inhomogeneous vectors included, by bilinearity), and homogeneity of $a$ in \lean{Pres_mem_O}.
\textbf{The trap of Definition~\ref{l12:def:lowest}:} \lean{IsLowestWeight N h} does not require $M_h\neq0$; if $M_h=0$ the
conclusions are statements about the zero space and are true but empty; the theorem is not weakened as a hypothesis, but it
does not by itself say that $M_h$ is the genuine bottom of $M$. Same as Zhu, \emph{J.~AMS} 9 (1996), Theorem~2.1.2 (the easy direction) and
Frenkel--Zhu Lemma~2.1.2; and Gannon (5.3.9).
\end{dictionary}

\begin{physics}
Zhu's theorem says that the zero modes (charges) acting on the lowest-weight space $M_h$, the space of primary states, form a representation of the
associative algebra $A(V)$: the product of zero modes is the Zhu product, and the null vectors in $O(V)$ act as zero. Physically this is the content
of the null-vector (fusion) constraints on the weights of primary fields; for example, the Virasoro vacuum relation gives the Kac condition on $h$. This is
the main tool for classifying the sectors of a rational CFT: the irreducible modules are produced from the irreducible representations of the
finite-dimensional algebra $A(V)$.
\end{physics}

\textbf{Proof on paper.} For $u\in V_k$, $v\in V_\ell$ and $x\in M_h$, modes $v_{(n)}$ with $n\ge\ell$ kill $x$ (they would lower the weight below $h$; compare
Lemma~\ref{l12:lem:modeweight}). Expand $o(u*v)x = \sum_m\binom km(u_{(m-1)}v)_{(k+\ell-m-1)}x$ by the iterate formula (Theorem~\ref{l03:thm:iterate}, in its
module form obtained by putting $m=0$ in (M3)); all terms with an annihilation mode acting first on $x$ vanish, and the binomial sums collapse to
$u_{(k-1)}v_{(\ell-1)}x = o(u)o(v)x$. For $O(V)$: $o((L_{-1}+L_0)u)=0$ on any module, because $(L_{-1}u)_{(k)} = -k\,u_{(k-1)}$
(Theorem~\ref{l12:thm:covariance}) cancels $L_0u = ku$; combined with the multiplicativity on $M_h$ this kills the generators $((L_{-1}+L_0)u)*v$
of $O(V)$. The residue statement is by induction on $n$: $n=0$ is the case $u\circ v\in O(V)$, and the induction step is the integration-by-parts
identity $(n+2)\,\mathrm{Pres}(n+1) = \mathrm{Pres}_{L_{-1}a}(n) + (M+n+1)\mathrm{Pres}(n)$, applied to $a$ and to $L_{-1}a$. \qed

\begin{status}
Unconditional given the displayed hypotheses. \texttt{\#print axioms} of \lean{VOA.oMod\_zhuMul\_lowest\_all}, \lean{VOA.rhoLW},
\lean{VOA.O\_le\_ker\_rhoLW} and \lean{VOA.Pres\_mem\_O} is \lean{[propext, Classical.}\allowbreak\lean{choice, Quot.}\allowbreak\lean{sound]}
for each (run for this chapter). The PBW theorem \lean{VOA.pbw} is not used.
\end{status}

\subsection{\texorpdfstring{$A(V)$}{A(V)} is an associative algebra}

\begin{theorem}[Zhu 1996, Theorem~2.1.1]\label{l12:thm:zhuring}
For every VOA $V$: $O(V)$ is a two-sided ideal for $*$, the product $*$ is associative modulo $O(V)$, and the
class of $\ltwelvevac$ is a two-sided unit. Hence $A(V)$ is an associative unital $\C$-algebra with product
$[u][v] = [u*v]$.
\end{theorem}
\noindent Files \leanpath{VOA/\allowbreak Modules/\allowbreak ZhuRing.lean} (the predicate) and
\leanpath{VOA/\allowbreak Modules/\allowbreak ZhuRingInstance.lean} (the theorem).
\begin{leancode}
structure ZhuRingData (W : VertexOperatorAlgebra 𝒱) : Prop where
  left_ideal : ∀ {u : 𝒱}, u ∈ O W → ∀ v, zhuMul W u v ∈ O W
  right_ideal : ∀ (u : 𝒱) {v : 𝒱}, v ∈ O W → zhuMul W u v ∈ O W
  assoc : ∀ u v w : 𝒱, zhuMul W (zhuMul W u v) w - zhuMul W u (zhuMul W v w) ∈ O W

theorem zhuRingData_holds (W : VertexOperatorAlgebra 𝒱) : ZhuRingData W :=
  zhuRingData_of_cores W
    (fun _ _ _ _ hu ha b => ZhuHR.right_ideal_core W hu ha b)
    (fun _ _ _ _ hu hv w => associator_core W hu hv w)

@[reducible] noncomputable def instRingA (W : VertexOperatorAlgebra 𝒱) : Ring (A W) :=
  AsRing (zhuRingData_holds W)

@[reducible] noncomputable def instAlgebraA (W : VertexOperatorAlgebra 𝒱) :
    @Algebra ℂ (A W) _ (instRingA W).toSemiring :=
  AsAlgebra (zhuRingData_holds W)
\end{leancode}

\begin{dictionary}
\textbf{Decoding.} \lean{ZhuRingData W} is a proposition (a \lean{structure ... : Prop}) bundling exactly the
three facts of Zhu's theorem. \lean{zhuRingData_holds} proves it for \emph{every} \lean{W} --- there is no
hypothesis. \lean{instRingA} then builds a \lean{Ring} structure on the quotient type \lean{A W} whose
multiplication is the descended Zhu product and whose unit is $[\ltwelvevac]$ (library lemmas
\lean{zhuMul_vacuum_left}, \lean{zhuMul_vacuum_right} give the unit laws already in $V$), reusing the quotient's
additive group (so ``$+$'' on $A(V)$ is the usual one). \lean{instAlgebraA} adds the $\C$-algebra structure.

\textbf{Why the ring is not a global instance.} \lean{instRingA} is a \lean{def}, not an \lean{instance}, and
every statement that multiplies in $A(V)$ opens it locally (\lean{letI := instRingA W}). This is a technicality
of Lean's instance system; mathematically $A(V)$ simply \emph{is} this ring.

\textbf{Same as the literature?} Yes: this is Zhu's Theorem~2.1.1 (and Frenkel--Zhu's), stated for VOAs of CFT
type, which is every library VOA. The three named cores \lean{ZhuHR.right_ideal_core} (right ideal, for
homogeneous $u$, $a$) and \lean{associator_core} (associativity for homogeneous $u$, $v$) are the genuine
content; the left-ideal property is derived from associativity, and the reduction from homogeneous to arbitrary
vectors is bilinearity.

\textbf{Could it be vacuous?} No: there are no hypotheses, and $A(V)$ is a genuine quotient (for the zero
VOA it is the zero ring, which is correct).
\end{dictionary}

\textbf{Proof on paper.} We follow Zhu. Write, for homogeneous $u$,
$u*v = \Res_zY(u,z)v\frac{(1+z)^{\wt u}}{z}$ and $u\circ v = \Res_zY(u,z)v\frac{(1+z)^{\wt u}}{z^2}$.
\begin{enumerate}
\item \emph{More of $O(V)$.} By induction on $n$, every residue
  $\Res_zY(u,z)v\frac{(1+z)^{\wt u+n}}{z^{2+n}}$, $n\ge0$, lies in $O(V)$: the case $n=0$ is the definition, and
  integrating by parts with $Y(L_{-1}u,z) = \partial_zY(u,z)$ relates the residue at $n+1$ to those at $n$ for
  $u$ and $L_{-1}u$ (library: \lean{Pres_mem_O}, Theorem~\ref{l12:thm:zhuaction}).
\item \emph{Skew formula.} Using skew-symmetry $Y(u,z)v = e^{zL_{-1}}Y(v,-z)u$ one shows
  $u*v - v*u \equiv \Res_zY(u,z)v(1+z)^{\wt u-1}$ modulo $O(V)$; in particular $O(V)$ is a left ideal once it is a
  right ideal.
\item \emph{Right ideal.} For $a$ homogeneous, $u*(a\circ b)$ is a residue in two variables; using the Jacobi
  identity to commute $Y(u,z_1)$ past $Y(a,z_2)$ one rewrites it as a combination of $a'\circ b'$ and of the
  residues of step~1, all in $O(V)$ (library: \lean{ZhuHR.right_ideal_core}, off a binomial ``linchpin''
  identity \lean{zhu_right_binom}).
\item \emph{Associativity.} $(u*v)*w - u*(v*w)$ is, by the iterate (associativity) formula of the Jacobi
  identity, a two-variable residue $\Res_{z_2}\Res_{z_1-z_2}$ minus $\Res_{z_1}\Res_{z_2}$; the difference of the
  two expansions is the residue of a function with a pole of order $\ge2$ at $z_2 = 0$ weighted by
  $(1+z)^{\wt}$, which step~1 places in $O(V)$ (library: \lean{associator_core}, the longest file of this
  chapter, \leanpath{VOA/\allowbreak Modules/\allowbreak ZhuAssociator.lean}).
\item \emph{Unit.} $\ltwelvevac*v = \ltwelvevac_{(-1)}v = v$ and $v*\ltwelvevac = \sum_m\binom km v_{(m-1)}\ltwelvevac = v$
  (only $m=0$ survives by the creation property).
\end{enumerate}
Bilinearity of $*$ then lets the quotient inherit a ring structure. \qed

\begin{status}
Proved, unconditional. \texttt{\#print axioms} (run for this chapter) of \lean{VOA.zhuRingData\_holds},
\lean{VOA.instRingA}, \lean{VOA.associator\_core} and \lean{ZhuHR.right\_ideal\_core}:
\lean{[propext, Classical.choice, Quot.sound]}. \lean{VOA.pbw} is not used.
\end{status}

\begin{physics}
$A(V)$ is the algebra of \emph{zero modes acting on ground states}. In any sector, zero modes of fields preserve
the space of ground states, and their products there are computed by the Zhu product --- a deformation of the
normal-ordered product $u_{(-1)}v$ that accounts for the reordering of modes. Associativity of $A(V)$ is the
associativity of operator composition on ground states, which physically is inherited from associativity of the
OPE. For the free boson $A(V) = \C[J_0]$ (the momentum), for an affine algebra $A(V)$ is a quotient of
$U(\mathfrak g)$ (the global charges), and for the Virasoro algebra $A(V)$ is a quotient of $\C[L_0]$ by the
null-vector relation --- the Kac-table equation.
\end{physics}

\begin{theorem}[Ground states form an $A(V)$-module; Zhu 1996, Theorem~2.1.2]\label{l12:thm:moduleMh}
If $M$ is an ordinary module with lowest weight $h$ (library sense), then $[u]\cdot x = o(u)x$ is a well-defined
unital $A(V)$-module structure on $M_h$.
\end{theorem}
\noindent Files \leanpath{VOA/\allowbreak Modules/\allowbreak ZhuModule.lean},
\leanpath{VOA/\allowbreak Modules/\allowbreak ZhuRingInstance.lean}.
\begin{leancode}
noncomputable def actA (N : OrdVModule W Mod) {h : ℂ} (hcw : IsLowestWeight N h) :
    A W →ₗ[ℂ] End ℂ (N.weight h) :=
  (O W).liftQ (rhoLW N h) (O_le_ker_rhoLW N hcw)

@[reducible] noncomputable def moduleMh (W : VertexOperatorAlgebra 𝒱) {N : OrdVModule W Mod} {h : ℂ}
    (hcw : IsLowestWeight N h) :
    letI := instRingA W; letI := instAlgebraA W; Module (A W) (N.weight h) :=
  AsModule (zhuRingData_holds W) hcw
\end{leancode}

\begin{dictionary}
\lean{rhoLW N h u} is $o(u)$ restricted to $M_h$ (Chapter~\ref{ch:01}). \lean{(O W).liftQ f hf} is the map
induced on the quotient $V/O(V)$ by a linear map $f$ vanishing on $O(V)$; the vanishing is
\lean{O_le_ker_rhoLW} (Theorem~\ref{l12:thm:zhuaction}). \lean{moduleMh} upgrades the linear map $A(V)\to\End(M_h)$
to a module structure, using multiplicativity \lean{oMod_zhuMul_lowest_all} and $o(\ltwelvevac) = \mathrm{id}$.
Hypothesis: \lean{IsLowestWeight N h} --- see the warning in Section~\ref{l12:sec:char}; if $M_h = 0$ the module is
zero. Same as Zhu's theorem (``easy direction'' of the correspondence).
\end{dictionary}

\begin{status}
Proved, unconditional given \lean{hcw}; \texttt{\#print axioms VOA.moduleMh} = \lean{[propext, Classical.choice,
Quot.sound]}.
\end{status}

\begin{physics}
The ground states of a sector carry a representation of the zero-mode algebra. For a primary field of weight $h$
in a Virasoro minimal model, $[\omega]$ acts on the one-dimensional ground space by $h$, and the null-vector
relation in $A(V)$ forces $h$ into the Kac table: this is the algebraic version of the BPZ argument that null
states restrict the allowed primaries.
\end{physics}

\subsection{The forward half of Zhu's correspondence}

\begin{definition}[$A(V)$-modules and lowest-weight modules, bundled]\label{l12:def:lwmodule}
An \lean{AMod} is a vector space $U$ with a linear map $\rho:A(V)\to\End U$ such that $\rho[\ltwelvevac] = 1$ and
$\rho[u*v] = \rho[u]\rho[v]$. An \lean{LWModule} is an ordinary module $M$ with a number $h$ such that all weights lie
in $h+\N$ \emph{and} $M_h\neq0$. Isomorphism classes are taken with respect to linear isomorphisms intertwining
the actions.
\end{definition}
\noindent File \leanpath{VOA/\allowbreak Modules/\allowbreak ZhuBijection.lean}.
\begin{leancode}
structure AMod where
  carrier : Type u
  [addCommGroup : AddCommGroup carrier]
  [module : Module ℂ carrier]
  ρ : A W →ₗ[ℂ] End ℂ carrier
  ρ_one : ρ (one W) = 1
  ρ_mul : ∀ u v : 𝒱, ρ (mk W (zhuMul W u v)) = ρ (mk W u) * ρ (mk W v)

structure LWModule where
  carrier : Type u
  [addCommGroup : AddCommGroup carrier]
  [module : Module ℂ carrier]
  isMod : OrdVModule W carrier
  h : ℂ
  isLow : IsLowestWeight isMod h
  realized : isMod.weight h ≠ ⊥

noncomputable def forwardMap : LWModuleIsoClass W → AModIsoClass W :=
  Quotient.lift (fun P => (Quotient.mk (AMod.setoid W) P.toAMod : AModIsoClass W))
    (fun _ _ ⟨e⟩ => Quotient.sound ⟨LWModule.isoToAModEquiv W e⟩)
\end{leancode}

\begin{dictionary}
\textbf{\lean{AMod}.} A ``bundled'' module: the carrier type is a field. \lean{mk W} is the projection
$V\to A(V)$, \lean{one W} is $[\ltwelvevac]$. Multiplicativity is required on classes of the form $[u*v]$, which is
every product in $A(V)$; so an \lean{AMod} is exactly a unital module over the ring $A(V)$, of any dimension,
possibly zero, possibly reducible.

\textbf{\lean{LWModule}.} Here \lean{realized} adds the missing ``$M_h\neq0$''; an \lean{LWModule} is not
required to be irreducible.

\textbf{\lean{forwardMap}.} \lean{LWModuleIsoClass W} and \lean{AModIsoClass W} are the quotients by
isomorphism (\lean{Quotient}); \lean{Quotient.lift} defines a map on classes from a map on representatives,
given a proof that isomorphic inputs give isomorphic outputs (\lean{LWModule.isoToAModEquiv}: a module
isomorphism preserves weight spaces, intertwines zero modes, and the two lowest weights agree because both
bottoms are nonzero). The map is $[M]\mapsto[M_h]$. Zhu's theorem says its restriction to
\emph{irreducible} modules is a bijection onto irreducible $A(V)$-modules; that is Chapter~\ref{ch:L13}: the
uniqueness half for every $V$ (\lean{zhuBijective_of_voaPBW}; the codomain of \lean{zhuBijective} is the image of the
map), and the full bijection onto simple $A(V)$-modules for $C_2$-cofinite $V$ (\lean{zhuCorrespondence_of_C2Cofinite}).
\end{dictionary}

\begin{status}
Definitions; \texttt{\#print axioms VOA.forwardMap} = \lean{[propext, Classical.choice, Quot.sound]}.
\end{status}

\begin{physics}
The forward map sends a sector to its space of ground states with the action of the zero modes --- in physics
language, a representation of the chiral algebra to the representation of the ``finite-dimensional'' symmetry
seen by its primaries. Zhu's correspondence is the statement that a primary field (an irreducible
representation of the zero-mode algebra on ground states) generates exactly one irreducible sector.
\end{physics}

\begin{warning}
\textbf{\lean{InducedAssignment} is uninhabited; the round-trip theorems built on it are vacuous (serious, recorded in
the library; docstrings stale, documentation only).} The file \leanpath{VOA/\allowbreak Modules/\allowbreak ZhuInduced.lean} expresses
the ``hard direction'' of Zhu's correspondence as the structure
\begin{leancode}
structure InducedAssignment where
  realize : AMod W → LWModule W
  realizes : ∀ U : AMod W, AModEquiv W (realize U).toAMod U
  isoInvariant : ∀ U U' : AMod W, Nonempty (AModEquiv W U U') →
    LWModule.Iso W (realize U) (realize U')

theorem forwardMap_invFun (A : InducedAssignment W) (x : AModIsoClass W) :
    forwardMap W (A.invFun x) = x := by
\end{leancode}
It asks for a lowest-weight module realising \emph{every} \lean{AMod}, including the \emph{zero}
$A(V)$-module (carrier a point; $\rho[\ltwelvevac] = 1$ holds because the endomorphism ring of the zero space is the
zero ring). But an \lean{LWModule} has $M_h\neq0$ (\lean{realized}), and no nonzero space is isomorphic to zero.
So no \lean{InducedAssignment} exists, for any $V$. The library proves exactly this, in a file of
Chapter~\ref{ch:L13} (\leanpath{VOA/\allowbreak Modules/\allowbreak ZhuVermaAnArchZhu534.lean}):
\begin{leancode}
theorem not_nonempty_inducedAssignment (W : VertexOperatorAlgebra 𝒱) :
    ¬ Nonempty (InducedAssignment W) := by
\end{leancode}
(\texttt{\#print axioms}: standard three). Even for nonzero $U$ a realisation can fail: $[\omega]$ acts on the bottom
of an \lean{LWModule} by the scalar $h$, so an $A(V)$-module on which $[\omega]$ is not scalar (e.g.\ a direct sum
of two irreducibles with different $h$) has no realisation. Consequently \lean{forwardMap_invFun},
\lean{forwardMap_}\allowbreak\lean{comp_}\allowbreak\lean{invFun} (this file) and \lean{forwardMap_}\allowbreak\lean{invFun_}\allowbreak\lean{ofPointwise}
(\leanpath{VOA/\allowbreak Modules/\allowbreak ZhuInducedPointwise.lean}) are true for the trivial reason that their
hypothesis can never be supplied; the pointwise input \lean{∀ U, HasInducedModule W U} is refuted by
\lean{not_}\allowbreak\lean{forall_}\allowbreak\lean{hasInducedModule}, and the file \leanpath{ZhuInducedPointwise.lean} itself notes that its second
input \lean{BottomIsoLifts} is false (a generalised Verma module and its irreducible quotient have the same bottom).
The docstrings of \leanpath{ZhuInduced.lean} (``remains reusable API'') and of \leanpath{ZhuBijection.lean}
(``the entire remaining content \ldots{} is located in the single precise object \lean{InducedAssignment}'') are
stale. The correct replacements are in Chapter~\ref{ch:L13}: for $C_2$-cofinite $V$, the genuine bijection between
$\Phi(V)$ and isomorphism classes of simple $A(V)$-modules is \lean{zhuCorrespondence_of_C2Cofinite} (conditional only on
$C_2$); for every $V$, the uniqueness half is \lean{zhuBijective}, a bijection onto its image. The latter is conditional on an explicit datum \lean{Pr : BottomProjection ...}, and its
wrapper \lean{zhuBijective_of_voaPBW} supplies \lean{Pr} from the theorem \lean{VOA.pbw} and is unconditional
(\texttt{\#print axioms} in a later run, after \lean{pbw} became a theorem, not re-run for this edition:
\lean{[propext, Classical.choice, Quot.sound]}; at the snapshot, run for this chapter, the list also contained \lean{pbw}, then an axiom).
\end{warning}

\subsection{Two consequences: bottoms are simple, and a Zhu-polynomial simplicity test}

\begin{theorem}[Simple bottoms]\label{l12:thm:simplebottom}
Assume $A(V)$ is a semisimple ring and the PBW datum \lean{Pr} of Chapter~\ref{ch:L13} is given. Then the bottom
$M_h$ of every irreducible lowest-weight module is a simple $A(V)$-module.
\end{theorem}
\noindent File \leanpath{VOA/\allowbreak Modules/\allowbreak ZhuBottomSimple.lean}.
\begin{leancode}
theorem isSimpleAMod_of_irreducible_of_semisimple
    (Pr : BottomProjection (parabolicRepOfAMod W (regularAMod W)))
    (hss : letI := instRingA W; IsSemisimpleRing (A W))
    (P : LWModule W) (hP : P.IsIrreducibleLW) : IsSimpleAMod P.toAMod := by
\end{leancode}

\begin{dictionary}
\lean{IsSemisimpleRing} is Mathlib's semisimplicity (every module is a direct sum of simple modules).
\lean{P.IsIrreducibleLW} is irreducibility of the underlying ordinary module, \lean{IsSimpleAMod} is ``nonzero
with no proper nonzero $A(V)$-submodule''. \lean{Pr} is the canonical-bottom projection of the Dong--Jiang
induced module (Chapter~\ref{ch:L13}); it is an explicit hypothesis, not an axiom. In the literature this
statement holds without semisimplicity and without \lean{Pr} (Zhu 1996, Theorem~2.2.1: the top level of an
irreducible module is an irreducible $A(V)$-module); the library's version is therefore \emph{weaker} (more
hypotheses) than Zhu's. \textbf{Superseded:} the general statement, with neither hypothesis, is now a library theorem,
\lean{VOA.Huang.LWBottomRecovery.isSimpleAMod_toAMod_of_irreducible} (via \lean{allBottomRecovers}), presented in
Chapter~\ref{ch:L13}, Theorem~\ref{l13:thm:simplebottom-all} (standard three axioms, checked at an earlier commit, not re-run for this edition).
\end{dictionary}

\textbf{Proof on paper.} A nonzero $A(V)$-submodule $S\subseteq M_h$ contains a simple submodule $T$ (semisimplicity).
The universal property of the induced module $\bar M(T)$ gives a module map $\bar M(T)\to M$ which is the inclusion on the
bottom; its image is a nonzero submodule, hence all of $M$ by irreducibility, and comparing bottoms gives $T = M_h$. \qed

\begin{status}
Proved, \emph{conditional on} \lean{Pr} and \lean{hss}; \texttt{\#print axioms} of
\texttt{VOA.\allowbreak isSimpleAMod\_\allowbreak of\_\allowbreak irreducible\_\allowbreak of\_\allowbreak semisimple} = \lean{[propext, Classical.choice, Quot.sound]}.
\end{status}

\begin{physics}
The ground states of an irreducible sector form an irreducible multiplet of the zero-mode algebra: if they split, the
descendants of one piece would form a smaller sector. This is the statement that each irreducible sector has
``one primary'' up to the zero-mode symmetry.
\end{physics}

\begin{theorem}[Zhu-polynomial simplicity criterion]\label{l12:thm:zhupoly}
Let $V\neq0$ with $V_1 = 0$, and let $P\in\C[x]$ satisfy $P([\omega]) = 0$ in $A(V)$. If the only non-negative
integer roots of $P$ are among $\{0,1\}$, then $V$ is simple.
\end{theorem}
\noindent File \leanpath{VOA/\allowbreak Modules/\allowbreak ZhuPolynomialSimplicity.lean}.
\begin{leancode}
def ZhuPolynomialRelation (P : Polynomial ℂ) : Prop :=
  letI := instRingA W
  letI := instAlgebraA W
  Polynomial.aeval (mk W W.conformalVector) P = 0

theorem isSimple_of_zhuPolynomial
    (hneW : (⊤ : Submodule ℂ 𝒱) ≠ ⊥)
    (hV1 : W.grading 1 = ⊥)
    (P : Polynomial ℂ)
    (hpoly : ZhuPolynomialRelation W P)
    (hroots : ∀ n : ℕ, P.eval (n : ℂ) = 0 → n = 0 ∨ n = 1) :
    W.IsSimple := by
\end{leancode}

\begin{dictionary}
\lean{Polynomial.aeval x P} is $P(x)$ evaluated in the algebra $A(V)$; \lean{mk W W.conformalVector} is $[\omega]$.
\lean{W.IsSimple} is simplicity of the VOA: $V\neq0$ and every ideal is $0$ or $V$ (Chapter~\ref{ch:L04};
an ideal \lean{W.IsIdeal I} is a subspace with $u_{(n)}I\subseteq I$). Hypotheses: $V\neq0$, $V_1 = 0$, the
polynomial relation, the root condition. The full Lean name is
\lean{VOA.}\allowbreak\lean{ZhuPolynomialSimplicity.}\allowbreak\lean{isSimple_}\allowbreak\lean{of_}\allowbreak\lean{zhuPolynomial}. This is a recognition principle formulated in the
project (it is used for the triplet algebra $W(2)$, Chapter~\ref{ch:L18}); it is a correct consequence of Zhu's
theory, not a literature theorem to be compared.
\end{dictionary}

\textbf{Proof on paper.} Let $I$ be a nonzero proper ideal. It is a $V$-submodule of the adjoint module, graded,
and has a least occupied weight $n\in\N$; then $I$ is an ordinary module with lowest weight $n$ and $I_n\neq0$
(\lean{idealOrdModule}, \lean{idealBottomWeight}). By Theorem~\ref{l12:thm:moduleMh}, $A(V)$ acts on $I_n$ and
$[\omega]$ acts by $L_0 = n$. Since $P([\omega]) = 0$, $P(n)I_n = 0$, so $P(n) = 0$ and $n\in\{0,1\}$. But $n=0$ would
put $\ltwelvevac$ in $I$ (as $V_0 = \C\ltwelvevac$) and then $I = V$; and $n = 1$ is impossible since $V_1=0$. \qed

\begin{status}
Proved; \texttt{\#print axioms} = \lean{[propext, Classical.choice, Quot.sound]}.
\end{status}

\begin{physics}
An ideal of the chiral algebra would be a set of null fields; its lowest state would be a primary of some weight
$n$ in the vacuum module. Zhu's algebra lists the allowed weights of primaries (roots of $P$); if none of them is a
positive integer $\ge2$, there is no room for a null primary, so the vacuum module is irreducible. This is the
chiral-algebra version of ``no null vectors at these levels''.
\end{physics}

%======================================================================
\section{\texorpdfstring{$C_2$}{C2}-cofiniteness and the Poisson algebra \texorpdfstring{$V/C_2(V)$}{V/C2(V)}}\label{l12:sec:c2}
%======================================================================

Zhu introduced a finiteness condition under which the characters of modules converge, and, together with rationality,
are modular. (The library proves convergence from $C_2$ alone, and modularity from $C_2$ together with rationality
(\lean{AdmissiblyRational}), or more generally with semisimplicity of $A(V)$; Chapter~\ref{ch:L15}.) It is the cleanest ``finitely many sectors'' condition that can be checked on the vacuum
module alone.

\begin{definition}[$C_2(V)$, $C_2$-cofiniteness; Zhu 1996, §4.4]\label{l12:def:c2}
$C_2(V) = \operatorname{span}\{u_{(-2)}v : u,v\in V\}$, $R_V = V/C_2(V)$, and $V$ is \emph{$C_2$-cofinite} if
$\dim R_V<\infty$.
\end{definition}
\noindent File \leanpath{VOA/\allowbreak Modules/\allowbreak C2Cofinite.lean}.
\begin{leancode}
def O₂ (W : VertexOperatorAlgebra 𝒱) : Submodule ℂ 𝒱 :=
  Submodule.span ℂ {x | ∃ u v : 𝒱, x = W.smul u (-2) v}

abbrev A₂ (W : VertexOperatorAlgebra 𝒱) : Type u := 𝒱 ⧸ O₂ W

def poissonMul (W : VertexOperatorAlgebra 𝒱) (u v : 𝒱) : 𝒱 := W.smul u (-1) v

def poissonBracket (W : VertexOperatorAlgebra 𝒱) (u v : 𝒱) : 𝒱 := W.smul u 0 v

def C2Cofinite (W : VertexOperatorAlgebra 𝒱) : Prop := FiniteDimensional ℂ (A₂ W)
\end{leancode}

\begin{dictionary}
\lean{O₂ W} is $C_2(V)$ (Zhu calls it $C_2(V)$; the library's ``$O_2$'' echoes Zhu's $O(V)$). \lean{A₂ W} is the
quotient $R_V = V/C_2(V)$. \lean{C2Cofinite W} is $\dim R_V<\infty$. Same as Zhu's definition, and as
Gaberdiel--Neitzke's and Arakawa's ($R_V$ is Zhu's $C_2$-algebra, Li's $V/C_2(V)$, Arakawa's
$R_V$). The standing CFT type is automatic.
\end{dictionary}

\begin{warning}
\textbf{Notation clash: \lean{A₂} is not the second higher Zhu algebra (documentation only).} In the literature, $A_n(V)$
(Dong--Li--Mason) denotes the $n$-th higher Zhu algebra (Section~\ref{l12:sec:an}), and $A_2(V)$ would be the
second one; Gaberdiel--Neitzke write $A_{[2]}(V)$ for $V/C_2(V)$, Arakawa writes $R_V$. The library's
\lean{A₂ W} is $V/C_2(V)$, the commutative Poisson algebra, while the higher Zhu algebras are \lean{An W n}. A
reader must not read \lean{A₂ W} as \lean{An W 2}.
\end{warning}

\begin{physics}
$C_2(V)$ is spanned by states containing at least one ``derivative'' or excited mode $u_{(-2)}$; the quotient
$R_V$ keeps only the ``classical'' part of the chiral algebra, the commutative algebra of fields modulo
derivatives. It is the coordinate ring of the \emph{associated variety}, the classical phase space of the chiral
algebra. $C_2$-cofiniteness (``lisse'') says that this phase space is a finite set of points (zero-dimensional):
the theory has no continuous moduli of vacua, no free bosons, and --- as Zhu, Gaberdiel--Neitzke, Abe--Buhl--Dong
and Miyamoto show --- finitely many irreducible sectors with convergent characters, modular in Zhu's sense when $A(V)$ is
semisimple and in Miyamoto's pseudo-trace sense in general. The free boson
fails it: $R_V = \C[J]$ is a line.
\end{physics}

\begin{theorem}[$R_V$ is a Poisson algebra; Zhu, Li]\label{l12:thm:poisson}
$C_2(V)$ is a two-sided ideal for both products $u_{(-1)}v$ and $u_{(0)}v$, and on $R_V$ the product
$[u][v] = [u_{(-1)}v]$ is commutative, associative and unital, the bracket $\{[u],[v]\} = [u_{(0)}v]$ is a Lie
bracket, and the Leibniz rule $\{a,bc\} = \{a,b\}c + b\{a,c\}$ holds.
\end{theorem}
\noindent Files \leanpath{VOA/\allowbreak Modules/\allowbreak C2Cofinite.lean},
\leanpath{VOA/\allowbreak Modules/\allowbreak PoissonDescent.lean}.
\begin{leancode}
theorem poissonBracket_leibniz (W : VertexOperatorAlgebra 𝒱) (u v w : 𝒱) :
    poissonBracket W u (poissonMul W v w) - poissonMul W v (poissonBracket W u w)
      = poissonMul W (poissonBracket W u v) w := by

noncomputable instance A₂CommRing (W : VertexOperatorAlgebra 𝒱) : CommRing (A₂ W) :=

noncomputable instance A₂LieRing (W : VertexOperatorAlgebra 𝒱) : LieRing (A₂ W) :=
\end{leancode}

\begin{dictionary}
\lean{CommRing (A₂ W)}: a commutative ring structure on the quotient whose product is $[u_{(-1)}v]$ and unit
$[\ltwelvevac]$ (the definitions inside the instance use \lean{A₂mul}, \lean{A₂bracket}, the descents of
\lean{poissonMul}, \lean{poissonBracket}). \lean{LieRing (A₂ W)}: a Lie bracket with \lean{A₂bracket}. The
Leibniz rule is stated in $V$ and holds \emph{exactly} there (it is the derivation property of $u_{(0)}$), hence
also in $R_V$. The library does not package a single ``Poisson algebra'' structure (Mathlib has no such class);
the two instances plus Leibniz are its content. Same as Li's Poisson structure on $V/C_2(V)$.
\end{dictionary}

\textbf{Proof on paper.} $u_{(-n)}v\in C_2(V)$ for every $n\ge2$: by Theorem~\ref{l12:thm:covariance} applied to the adjoint module, $(L_{-1}u)_{(-n)} = n\,u_{(-n-1)}$, so $u_{(-n-1)}v = \frac1n(L_{-1}u)_{(-n)}v$ and induction on $n$ applies. For the
ideal property, the iterate formula writes $(a_{(-2)}b)_{(m)}x$ ($m\le0$) as a sum of terms with an outer mode
$\le-2$, which lie in $C_2(V)$; commutativity $u_{(-1)}v - v_{(-1)}u\in C_2(V)$ and antisymmetry
$u_{(0)}v + v_{(0)}u\in C_2(V)$ follow from skew-symmetry $Y(u,z)v = e^{zL_{-1}}Y(v,-z)u$, since every correction
term contains $L_{-1}$, and $L_{-1}V\subseteq C_2(V)$ ($L_{-1}x = x_{(-2)}\ltwelvevac$). Associativity modulo $C_2$
is the iterate formula at $\ell = -1$; the Jacobi identity for the bracket and the Leibniz rule are the
statement that $u_{(0)}$ is a derivation of all products (commutator formula with $m = 0$). \qed

\begin{status}
Proved, unconditional; \texttt{\#print axioms} of \lean{VOA.A₂CommRing}, \lean{VOA.A₂LieRing},
\lean{VOA.poissonBracket\_leibniz}: standard three.
\end{status}

\begin{physics}
The Poisson bracket $\{u,v\} = u_{(0)}v$ is the classical limit of the OPE's simple pole, i.e.\ of the charge
commutators; the commutative product $u_{(-1)}v$ is the classical limit of normal ordering. $R_V$ is thus the
classical (semiclassical, $\hbar\to0$) limit of the chiral algebra, and its Poisson structure records the
leading-order Ward identities. Arakawa's associated variety is $\operatorname{Specm}R_V$ with this Poisson
structure.
\end{physics}

\begin{theorem}[$C_2$-cofinite implies finitely strongly generated; Zhu, Gannon 5.3.11a]\label{l12:thm:fingen}
Let $S\subseteq V$ be a set of homogeneous vectors with $V = \operatorname{span}S + C_2(V)$. Then $V$ is spanned by
the vectors $s^1_{(-m_1)}\cdots s^k_{(-m_k)}\ltwelvevac$ with $s^i\in S$, $m_i\ge1$, $k\ge0$. In particular, if $V$ is
$C_2$-cofinite, lifts of a basis of $R_V$ strongly generate $V$.
\end{theorem}
\noindent Files \leanpath{VOA/\allowbreak Modules/\allowbreak FiniteGeneration.lean},
\leanpath{VOA/\allowbreak Modules/\allowbreak FiniteGenerationProof.lean}.
\begin{leancode}
def IsHomog (W : VertexOperatorAlgebra 𝒱) (x : 𝒱) : Prop := ∃ n : ℤ, x ∈ W.grading n

inductive IsIterProd (W : VertexOperatorAlgebra 𝒱) (S : Set 𝒱) : 𝒱 → Prop
  | vacuum : IsIterProd W S W.vacuum
  | step (s : 𝒱) (hs : s ∈ S) (m : ℤ) (hm : 1 ≤ m) (x : 𝒱) (hx : IsIterProd W S x) :
      IsIterProd W S (W.smul s (-m) x)

def iterProdSpan (W : VertexOperatorAlgebra 𝒱) (S : Set 𝒱) : Submodule ℂ 𝒱 :=
  Submodule.span ℂ {x | IsIterProd W S x}

theorem iterProdSpan_eq_top (W : VertexOperatorAlgebra 𝒱) (S : Set 𝒱)
    (hShom : ∀ s ∈ S, IsHomog W s)
    (htop : (⊤ : Submodule ℂ 𝒱) = Submodule.span ℂ S ⊔ O₂ W) :
    iterProdSpan W S = ⊤ :=
\end{leancode}

\begin{dictionary}
\lean{IsIterProd W S} is the smallest set containing $\ltwelvevac$ and closed under $x\mapsto s_{(-m)}x$ for $s\in S$,
$m\ge1$ (an \lean{inductive} predicate: exactly the iterated words). \lean{⊔} is the sum of subspaces. The
hypothesis \lean{htop} is $V = \operatorname{span}S + C_2(V)$, which for $C_2$-cofinite $V$ holds with $S$ finite
(\lean{iterProdSpan_}\allowbreak\lean{eq_}\allowbreak\lean{top_}\allowbreak\lean{of_}\allowbreak\lean{finite}). The words are not required to be ordered, so this is the weak form of
strong generation (Gaberdiel--Neitzke prove an ordered PBW-type refinement; the library's ordered version,
needed for $C_m$, is in Chapter~\ref{ch:L15}). Same as Gannon 5.3.11a / Zhu Lemma~4.4.1.
\end{dictionary}

\textbf{Proof on paper.} Induction on the weight. Let $T = \operatorname{span}$ of the words. Weight-$n$ vectors are in
$\operatorname{span}S + (C_2(V)\cap V_n)$, so it suffices to show the homogeneous generators $a_{(-2)}b$ of weight $n$ lie in
$T$, assuming $V_k\subseteq T$ for $k<n$. Write $a$ itself modulo $C_2$ as a combination of $S$ plus elements of $C_2(V)$ of
the same weight $<n$, and peel the left factor with the iterate formula: an $S$-left factor gives a word step; a
$C_2$-left factor recurses on a strictly smaller weight, and all right factors that appear have weight $<n$.
(The library's double induction is ``\lean{genPeel_master}''.) \qed

\begin{status}
Proved, unconditional; \texttt{\#print axioms VOA.iterProdSpan\_eq\_top}: standard three. A companion,
\lean{iterProdSpan_smul_closed} (\leanpath{VOA/\allowbreak Modules/\allowbreak PairClosure.lean}), shows that if all
products $s_{(j)}t$, $j\ge0$, of generators of weight $\ge1$ lie in the span of words, then that span is closed under
\emph{all} products; and \leanpath{VOA/\allowbreak Modules/\allowbreak NilpotentGenerators.lean} proves the pure algebra
fact that a commutative algebra generated by finitely many nilpotents $g_i$ with $g_i^{m_i} = 0$ has dimension at most
$\prod m_i$ (\lean{NilpotentGenerators.finrank_le_prod}), the engine for proving $C_2$-cofiniteness of the triplet
algebra (Chapter~\ref{ch:L18}). Both axiom-clean.
\end{status}

\begin{physics}
A lisse chiral algebra is generated, in the sense of normal-ordered products and derivatives, by finitely many
fields (for the Virasoro minimal models, $T$ alone; for the triplet, $T$ and $W^{\pm,0}$). The OPEs among the
generators then determine everything, which is why a CFT is specified by a finite OPE table.
\end{physics}

\begin{example}[Which VOAs are $C_2$-cofinite, in the library]\label{l12:ex:c2}
\begin{itemize}
\item The rank-one Heisenberg (free boson) VOA is \emph{not} $C_2$-cofinite: \lean{heisenberg_not_C2Cofinite}
  (\leanpath{VOA/\allowbreak Examples/\allowbreak HeisenbergC2.lean}, Chapter~\ref{ch:L05}), proved by exhibiting the
  linearly independent classes $[(a_{(-1)})^n\ltwelvevac]$ in $R_V$.
\item The lattice triplet algebra $W(2)$ \emph{is} $C_2$-cofinite and is \emph{not} admissibly rational:
  \lean{VOA.}\allowbreak\lean{LatGen.}\allowbreak\lean{Wp.}\allowbreak\lean{tripletWVOA_}\allowbreak\lean{two_}\allowbreak\lean{c2Cofinite_}\allowbreak\lean{and_}\allowbreak\lean{not_}\allowbreak\lean{admissiblyRational}
  (Chapter~\ref{ch:L18}; \texttt{\#print axioms} run for this chapter: standard three).
\item Every VOA with $V = \C\ltwelvevac$ is $C_2$-cofinite (\lean{c2Cofinite_of_span_vacuum}).
\item Every lattice VOA $V_L$ ($L$ even, positive definite) is $C_2$-cofinite and strongly finite:
  \lean{latticeVOAGen_c2Cofinite}, \lean{latticeVOAGen_stronglyFinite}
  (\leanpath{VOA/\allowbreak Classification/\allowbreak COne/\allowbreak LatticeC2.lean}, Chapter~\ref{ch:L22}); in
  particular the rank-one $V_{\sqrt{2m}\Z}$ (\lean{V2m_c2Cofinite}, \lean{V2m_stronglyFinite}).
\end{itemize}
$C_2$-cofiniteness of $\Vnat$ is not proved in the library.
\end{example}

\begin{physics}
The triplet example is the physically important one: $W(2)$ (symplectic fermions' even part, $c=-2$) is the
prototype logarithmic CFT, with finitely many sectors (lisse) but indecomposable, non-semisimple sector
structure. It shows that ``finitely many sectors'' does not imply ``rational''.
\end{physics}

\begin{definition}[Zeroth Poisson homology; Etingof--Schedler, Arakawa--Kawasetsu]\label{l12:def:etingof}
$\ltwelveHP(R_V) = R_V/\{R_V,R_V\}$, where $\{R_V,R_V\}$ is the span of all Poisson brackets. $V$ is
\emph{Etingof-finite} if $\dim\ltwelveHP(R_V)<\infty$.
\end{definition}
\noindent File \leanpath{VOA/\allowbreak Modules/\allowbreak PoissonTrace.lean} (assigned to Chapter~\ref{ch:L16}).
\begin{leancode}
noncomputable def poissonSpan (W : VertexOperatorAlgebra 𝒱) : Submodule ℂ (A₂ W) :=
  bilinSpan (A₂bracket W)

noncomputable abbrev HP₀ (W : VertexOperatorAlgebra 𝒱) : Type u := A₂ W ⧸ poissonSpan W

def EtingofFinite (W : VertexOperatorAlgebra 𝒱) : Prop := FiniteDimensional ℂ (HP₀ W)

theorem EtingofFinite_of_C2Cofinite (W : VertexOperatorAlgebra 𝒱) (h : C2Cofinite W) :
    EtingofFinite W := by
\end{leancode}

\begin{dictionary}
\lean{bilinSpan (A₂bracket W)} is the span of the values $\{p,q\}$ of the descended bracket (library lemma
\lean{poissonSpan_def}). $\ltwelveHP$ is a vector space, not a ring. \textbf{Same as the literature?}
Arakawa--Kawasetsu call $V$ \emph{quasi-lisse} if the associated variety $\operatorname{Specm}R_V$ has finitely many
symplectic leaves; Etingof--Schedler show that this implies $\dim\ltwelveHP(R_V)<\infty$, which is the only consequence
Arakawa--Kawasetsu use for modularity. So \lean{EtingofFinite} is \emph{weaker} than (implied by) quasi-lisse. Under it
the library proves the Arakawa--Kawasetsu MLDE
(\lean{IsSqPrimary.}\allowbreak\lean{exists_}\allowbreak\lean{monicMLDEOn_}\allowbreak\lean{all_}\allowbreak\lean{analyticBlocks_}\allowbreak\lean{of_}\allowbreak\lean{etingofFinite_uncond},
Chapter~\ref{ch:L16}): conditional only on \lean{EtingofFinite} (its former input \lean{AKWeightGrading W} is now the theorem
\lean{akWeightGrading W}), one monic MLDE of unbounded order for all abstract analytic genus-one blocks. It has no Zhu-type
modular-invariance theorem under it beyond the trivial VOA. The library proves
$C_2$-cofinite $\Rightarrow$ Etingof-finite; strictness of the implication is shown in Chapter~\ref{ch:L16} (the
simple affine $\mathfrak{sl}_2$ quotient at an admissible level).
\end{dictionary}

\textbf{Proof on paper} (of the implication). $\ltwelveHP(R_V)$ is a quotient of $R_V$. \qed

\begin{status}
Proved; \texttt{\#print axioms VOA.EtingofFinite\_of\_C2Cofinite}: standard three.
\end{status}

\begin{physics}
Quasi-lisse chiral algebras (Arakawa--Kawasetsu) are those whose classical phase space has finitely many
symplectic leaves: the Higgs-branch chiral algebras of four-dimensional $\mathcal N=2$ superconformal theories
(Beem--Rastelli and collaborators) are the motivating examples. $\ltwelveHP(R_V)^*$ is the space of Poisson traces,
the classical shadow of torus one-point functions; its finite-dimensionality is what makes characters satisfy a
modular linear differential equation even when the theory has continuous families of states.
\end{physics}

%======================================================================
\section{Higher Zhu algebras and finiteness under \texorpdfstring{$C_2$}{C2}}\label{l12:sec:an}
%======================================================================

Zhu's algebra sees only the ground states of a sector. To see the first $n$ excited levels one uses the higher
Zhu algebras.

\begin{definition}[Higher Zhu algebras; Dong--Li--Mason 1998; the spaces $C_m(V)$]\label{l12:def:an}
For $n\in\N$ and homogeneous $a$,
\begin{equation*}
 a\ltwelvecirc_n b = \Res_zY(a,z)b\frac{(1+z)^{\wt a+n}}{z^{2n+2}} = \sum_{i\ge0}\binom{\wt a+n}{i}a_{(i-2n-2)}b,
\end{equation*}
extended bilinearly; $O_n^{\mathrm{lib}}(V) = \operatorname{span}\{a\ltwelvecirc_nb\}$ and
$A_n^{\mathrm{lib}}(V) = V/O_n^{\mathrm{lib}}(V)$. Also $C_m(V) = \operatorname{span}\{u_{(-m)}v\}$.
\end{definition}
\noindent Files \leanpath{VOA/\allowbreak Modules/\allowbreak ZhuVermaAn.lean} (Chapter~\ref{ch:L13}) and
\leanpath{VOA/\allowbreak Modularity/\allowbreak AnFiniteDim.lean} (Chapter~\ref{ch:L15}).
\begin{leancode}
noncomputable def circFamily (W : VertexOperatorAlgebra 𝒱) (n : ℕ) (b : 𝒱) (k : ℤ) :
    (W.grading k) →ₗ[ℂ] 𝒱 where
  toFun ak := ∑ i ∈ Finset.range ((k + n).toNat + 1),
    cbz (k + (n : ℤ)) i • W.smul (ak : 𝒱) ((i : ℤ) - 2 * (n : ℤ) - 2) b

noncomputable def circ (W : VertexOperatorAlgebra 𝒱) (n : ℕ) (a b : 𝒱) : 𝒱 :=
  letI := DirectSum.IsInternal.chooseDecomposition W.grading W.isInternal
  (DFinsupp.lsum ℂ (circFamily W n b)) (DirectSum.decompose W.grading a)

def On (W : VertexOperatorAlgebra 𝒱) (n : ℕ) : Submodule ℂ 𝒱 :=
  Submodule.span ℂ {x | ∃ a b : 𝒱, x = circ W n a b}

abbrev An (W : VertexOperatorAlgebra 𝒱) (n : ℕ) : Type u := 𝒱 ⧸ On W n

def CGen (W : VertexOperatorAlgebra 𝒱) (M : ℤ) : Submodule ℂ 𝒱 :=
  Submodule.span ℂ {x | ∃ u v : 𝒱, x = W.smul u M v}
\end{leancode}

\begin{dictionary}
\lean{circ W n a b} is $a\ltwelvecirc_nb$, assembled from homogeneous components exactly as \lean{zhuMul} (Chapter~\ref{ch:01}
explains the \lean{DFinsupp.lsum} pattern); the range $\{0,\dots,\wt a+n\}$ is where the binomial coefficient is
nonzero. \lean{CGen W (-m)} is $C_m(V)$; \lean{CGen_neg_two : CGen W (-2) = O₂ W} holds by definition. At $n=0$ the
library proves \lean{On_zero_eq_O : On W 0 = O W} (Chapter~\ref{ch:L13}), so $A_0 = A(V)$.
\end{dictionary}

\begin{warning}
\textbf{The library's $O_n(V)$ omits generators (moderate), and the attribution is to DLM (documentation only).} Dong--Li--Mason
(\emph{J. Algebra} 206, 1998), who introduced $A_n(V)$, define $O_n(V)$ as the span of the $a\ltwelvecirc_nb$
\emph{and} of all $(L_{-1}+L_0)a$. At $n = 0$ the extra generators are redundant ($(L_{-1}+L_0)a = a\ltwelvecirc_0\ltwelvevac$),
which is why \lean{On_zero_eq_O} holds; for $n\ge1$, $a\ltwelvecirc_n\ltwelvevac$ is a different vector and the library does
not prove that $(L_{-1}+L_0)a\in O_n^{\mathrm{lib}}(V)$ (its own comment in
\leanpath{VOA/\allowbreak Modules/\allowbreak ZhuVermaAnArch.lean} says this family ``is what \lean{On W n} omits'').
Hence $O_n^{\mathrm{lib}}\subseteq O_n^{\mathrm{DLM}}$ and $A_n^{\mathrm{lib}}(V)$ \emph{surjects onto} the literature's
$A_n(V)$; it may be larger. For the finiteness theorem below this is harmless --- finiteness of the larger space
implies finiteness of DLM's --- so that theorem is at least as strong as the literature's. But
\lean{An W n} must not be identified with DLM's $A_n(V)$ without an argument. The library's docstrings attribute
$A_n$ to Frenkel--Zhu; Frenkel--Zhu (1992) introduced the $A(V)$-bimodules $A(M)$, while the tower $A_n(V)$ is
Dong--Li--Mason's.
\end{warning}

\begin{physics}
$A_n(V)$ is the algebra of all modes that map the first $n$ excited levels of any sector to themselves; its
representations classify sectors together with their low-lying descendants. Physically it records the
level-truncated OPE algebra: how many and which states appear up to energy $h+n$.
\end{physics}

\begin{theorem}[$C_2$ controls the whole tower; Gaberdiel--Neitzke, Zhu, Dong--Li--Mason]\label{l12:thm:an}
If $V$ is $C_2$-cofinite then: (i) $V/C_m(V)$ is finite-dimensional for every $m\ge1$; (ii) $A(V)$ is
finite-dimensional; (iii) $A_n^{\mathrm{lib}}(V)$ is finite-dimensional for every $n\in\N$.
\end{theorem}
\noindent File \leanpath{VOA/\allowbreak Modules/\allowbreak C2ImpliesZhuFinite.lean}.
\begin{leancode}
theorem finiteDimensional_quot_CGen_of_C2Cofinite (W : VertexOperatorAlgebra 𝒱)
    (hC2 : C2Cofinite W) (m : ℕ) (hm : 1 ≤ m) :
    FiniteDimensional ℂ (𝒱 ⧸ CGen W (-(m : ℤ))) := by

theorem zhuAlgebra_finiteDimensional_of_C2Cofinite (W : VertexOperatorAlgebra 𝒱)
    (hC2 : C2Cofinite W) : FiniteDimensional ℂ (A W) :=
  finiteDimensional_A_of_C2Cofinite W hC2

theorem finiteDimensional_An_of_C2Cofinite (W : VertexOperatorAlgebra 𝒱)
    (hC2 : C2Cofinite W) (n : ℕ) : FiniteDimensional ℂ (An W n) := by
\end{leancode}

\begin{dictionary}
Hypothesis: only \lean{C2Cofinite W}; the conclusions are literal finite-dimensionality of the quotient spaces just
defined. (i) at $m = 1$ is trivial ($C_1(V) = V$, since $v = \ltwelvevac_{(-1)}v$), at $m = 2$ it is the hypothesis, and for $m\ge3$ it is the content (since $C_{m+1}\subseteq C_m$, larger $m$ is
\emph{harder}). (ii) is Zhu's (and Gaberdiel--Neitzke's) theorem; the library's proof is the one-line alias of
\lean{finiteDimensional_}\allowbreak\lean{A_}\allowbreak\lean{of_}\allowbreak\lean{C2Cofinite} (\leanpath{VOA/\allowbreak Modularity/\allowbreak ZhuAlgebraFinite.lean},
Chapter~\ref{ch:L15}). (iii) concerns the library's $A_n$ (see the warning above); since it may be bigger than DLM's,
the statement is at least as strong as the literature's ``$C_2$-cofinite $\Rightarrow\dim A_n(V)<\infty$''. The same file
also proves finiteness of the Dong--Jiang bimodules $A_{n,0}(V)$ (\lean{finiteDimensional_}\allowbreak\lean{AnmFull_}\allowbreak\lean{of_}\allowbreak\lean{C2Cofinite}); their
definition belongs to Chapter~\ref{ch:L13} and we do not present it.
\end{dictionary}

\textbf{Proof on paper.} (ii): modulo $O(V)$, $u_{(-2)}v$ equals $u\circ v$ minus terms $u_{(i-2)}v$, $i\ge1$, of strictly
smaller weight; so by induction on weight, a finite set spanning $V$ modulo $C_2(V)$ spans $V$ modulo $O(V)$. (i):
Gaberdiel--Neitzke's Proposition~8 refines Theorem~\ref{l12:thm:fingen}: $V$ is spanned by ordered words
$s^1_{(-n_1)}\cdots s^k_{(-n_k)}\ltwelvevac$ with $n_1>n_2>\dots>n_k>0$ (strict decrease, achieved by the commutator formula
and $C_2$-cofiniteness); a word with $n_1\ge m$ lies in $C_m(V)$, and strictly decreasing words with all $n_i<m$ have length
$<m$, so finitely many of them span $V$ modulo $C_m(V)$ (their Theorem~11). (iii): $a\ltwelvecirc_nb = a_{(-2n-2)}b + {}$terms of
lower weight, so the same induction as in (ii), with $C_{2n+2}$ in place of $C_2$, shows that a finite set spanning
$V$ modulo $C_{2n+2}(V)$ spans $V$ modulo $O_n^{\mathrm{lib}}(V)$; apply (i) with $m = 2n+2$. \qed

\begin{status}
Proved, unconditional given \lean{hC2}; \texttt{\#print axioms} of all three (run for this chapter): standard three, no
\lean{VOA.pbw}.
\end{status}

\begin{physics}
In a lisse theory each energy window of each sector is described by a finite-dimensional algebra. Physically: finitely many
generators with a finite OPE table can only produce finitely many independent ``level-$n$ multiplets'', so the spectrum of
sectors up to any fixed excitation is finite. This is the input to Zhu's proof that torus one-point functions satisfy a
finite system of modular differential equations.
\end{physics}

%======================================================================
\section{The contragredient module}\label{l12:sec:contra}
%======================================================================

\begin{definition}[Restricted dual and contragredient action; FHL §5.2]\label{l12:def:contra}
Let $M$ be an ordinary module. Its \emph{restricted dual} is $M' = \bigoplus_\alpha(M_\alpha)^*$. For $u\in V_{\wt u}$,
$f\in M'$ and $x\in M$, the contragredient action is defined by
\begin{equation*}
 \langle Y'(u,z)f, x\rangle = \big\langle f, Y\big(e^{zL_1}(-z^{-2})^{L_0}u, z^{-1}\big)x\big\rangle ,
\end{equation*}
whose mode form is
\begin{equation*}
 \langle u'_{(n)}f, x\rangle = (-1)^{\wt u}\sum_{j\ge0}\frac1{j!}\big\langle f,\ (L_1^ju)_{(2\wt u-n-2-j)}x\big\rangle .
\end{equation*}
\end{definition}
\noindent Files \leanpath{VOA/\allowbreak Modules/\allowbreak Contragredient.lean},
\leanpath{VOA/\allowbreak Modules/\allowbreak ContragredientAction.lean},
\leanpath{VOA/\allowbreak Modules/\allowbreak ContragredientJacobi.lean}.
\begin{leancode}
def RestrictedDual (N : OrdVModule W M) : Type u := DirectSum ℂ (fun α => Module.Dual ℂ (N.weight α))

noncomputable def RestrictedDual.weight (N : OrdVModule W M) (α : ℂ) :
    Submodule ℂ (RestrictedDual N) :=
  LinearMap.range (DirectSum.lof ℂ ℂ (fun β => Module.Dual ℂ (N.weight β)) α)

noncomputable def dualCompMap (N : OrdVModule W M) (wt : ℤ) (u : 𝒱) (n : ℤ) (γ : ℂ) :
    Module.Dual ℂ (N.weight γ) →ₗ[ℂ] Module.Dual ℂ M :=
  (-1 : ℂ) ^ wt • ∑ j ∈ Finset.range (wt.toNat + 1),
    ((j.factorial : ℂ)⁻¹) • LinearMap.lcomp ℂ ℂ (dualTargetOp N wt u n γ j)

noncomputable def dualModeEndo (N : OrdVModule W M) (wt : ℤ) (u : 𝒱) (n : ℤ) :
    RestrictedDual N →ₗ[ℂ] RestrictedDual N :=
  DirectSum.toModule ℂ ℂ (RestrictedDual N) (fun γ => dualPieceMap N wt u n γ)

theorem dualModeEndo_pairing_out (N : OrdVModule W M) (wt : ℤ) (u : 𝒱) (n : ℤ) (γ : ℂ)
    (g : Module.Dual ℂ (N.weight γ)) {x : M} (hx : x ∈ N.weight (dualOutIdx wt n γ)) :
    dualPairing N x
        (dualModeEndo N wt u n
          (DirectSum.lof ℂ ℂ (fun β => Module.Dual ℂ (N.weight β)) γ g))
      = dualCompMap N wt u n γ g x := by

noncomputable def YMdual (N : OrdVModule W M) : 𝒱 →ₗ[ℂ] (End ℂ (RestrictedDual N))⟦z^±⟧ where
  toFun u := fun p => dualModeAt N (-p - 1) u
\end{leancode}

\begin{dictionary}
\textbf{Decoding.} \lean{Module.Dual ℂ X} is the dual space $X^*$; \lean{DirectSum ℂ F} is the external direct sum
$\bigoplus_\alpha F(\alpha)$, so \lean{RestrictedDual N} is $M' = \bigoplus_\alpha M_\alpha^*$ (not the full dual $M^*$, which is
the direct \emph{product}). \lean{DirectSum.lof ... α} is the inclusion of the $\alpha$-summand, so
\lean{RestrictedDual.weight N α} is $(M')_\alpha = M_\alpha^*$. \lean{dualTargetOp N wt u n γ j} is $x\mapsto$ the
$M_\gamma$-component of $(L_1^ju)_{(2\,\mathrm{wt}-n-2-j)}x$, and \lean{LinearMap.lcomp} precomposes a functional
$g\in M_\gamma^*$ with it. So \lean{dualCompMap N wt u n γ g} is the functional
$x\mapsto(-1)^{\wt u}\sum_{j=0}^{\wt u}\frac1{j!}\,g\big(\mathrm{pr}_\gamma(L_1^ju)_{(2\wt u-n-2-j)}x\big)$, exactly the
mode formula (the sum stops at $j = \wt u$ because $L_1^ju\in V_{\wt u-j}$ vanishes for $j>\wt u$ by CFT type).
\lean{dualModeEndo} assembles these over all $\gamma$, putting the image of $g\in M_\gamma^*$ into the slot
$\gamma+\wt u-n-1$ (\lean{dualOutIdx}); \lean{dualModeEndo_pairing_out} certifies that pairing the result with
$x$ in that slot gives back the formula, so nothing is lost by the reindexing. \lean{YMdual} is the field
$Y'(u,z) = \sum_n u'_{(n)}z^{-n-1}$ (\lean{dualModeAt N n u} is $u'_{(n)}$, assembled over homogeneous components of
$u$ as for \lean{zhuMul}).

\textbf{Derivation of the mode formula.} For $u$ homogeneous, $(-z^{-2})^{L_0}u = (-1)^{\wt u}z^{-2\wt u}u$, and
$e^{zL_1} = \sum_j z^jL_1^j/j!$. With $Y(w,z^{-1}) = \sum_m w_{(m)}z^{m+1}$, the coefficient of $z^{-n-1}$ of
$(-1)^{\wt u}\sum_j\frac{z^{j-2\wt u}}{j!}Y(L_1^ju,z^{-1})$ comes from $j-2\wt u+m+1 = -n-1$, i.e.\
$m = 2\wt u-n-2-j$. This is the Lean formula. So the definition is \emph{the same} as FHL's (5.2.4) and Gannon's
(5.1.11a), not a variant.
\end{dictionary}

\begin{physics}
The contragredient sector $M'$ is the \emph{charge conjugate} (antiparticle) sector: its states are the bras
dual to the kets of $M$, and $Y'$ is the field obtained by the BPZ conjugation $z\mapsto 1/z$ of the field acting
on kets, with the factor $e^{zL_1}(-z^{-2})^{L_0}$ being the conformal transformation law of a quasi-primary under
inversion. In a rational CFT the fusion rules satisfy $N_{MM'}^V = 1$: $M'$ is the unique sector that fuses with
$M$ to the vacuum.
\end{physics}

\begin{theorem}[$M'$ is an ordinary module; FHL Theorem~5.2.1]\label{l12:thm:contra}
For every ordinary module $M$, $(M',Y')$ is an ordinary $V$-module with $(M')_\alpha = M_\alpha^*$; in particular
$\chi_{M'} = \chi_M$. Moreover, for quasi-primary $v\in V_n$ (that is, $L_1v = 0$),
$\chi_{M'}(v;q) = (-1)^n\chi_M(v;q)$.
\end{theorem}
\noindent Files \leanpath{VOA/\allowbreak Modules/\allowbreak ContragredientJacobi.lean},
\leanpath{VOA/\allowbreak Modules/\allowbreak ContragredientClose.lean},
\leanpath{VOA/\allowbreak Modules/\allowbreak Contragredient.lean},
\leanpath{VOA/\allowbreak Modules/\allowbreak ContragredientCharacterClose.lean}.
\begin{leancode}
def ContragredientJacobi (N : OrdVModule W M) : Prop :=
  ∀ (u v : 𝒱) (x : RestrictedDual N) (ℓ m n : ℤ),
    (∑ᶠ i : ℤ, (negOnePow i * cbz ℓ i) •
        (mode (YMdual N u) (ℓ + m - i) (mode (YMdual N v) (n + i) x)
          - negOnePow ℓ • mode (YMdual N v) (ℓ + n - i) (mode (YMdual N u) (m + i) x)))
      = ∑ᶠ i : ℤ, cbz m i • mode (YMdual N (W.smul u (ℓ + i) v)) (m + n - i) x

theorem contragredientJacobi_holds (N : OrdVModule W M) : ContragredientJacobi N :=

noncomputable def contragredientWeakModule (N : OrdVModule W M)
    (hJac : ContragredientJacobi N) : WeakModule W (RestrictedDual N) where
  YM := YMdual N
  vacuum_YM := YMdual_vacuum N
  truncation := YMdual_truncation N
  jacobi := hJac

noncomputable def contragredientOrdVModule (N : OrdVModule W M) (hJac : ContragredientJacobi N) :
    OrdVModule W (RestrictedDual N) where
  toWeakModule := contragredientWeakModule N hJac
  weight := RestrictedDual.weight N

noncomputable def contragredientModule (N : OrdVModule W M) : OrdVModule W (RestrictedDual N) :=
  contragredient_of_dualAdjResidual N (dualAdjResidual_of_mvec N (dualAdjResidualMvec_holds N))

theorem contragredientModule_eq_contragredient (N : OrdVModule W M) :
    contragredientModule N = contragredient N (contragredientJacobi_holds N) := rfl

theorem contragredient_gradedDim (N : OrdVModule W M) :
    contragredientGradedDim N = gradedDim N := by

theorem onePoint_contragredient_holds (N : OrdVModule W M) {n : ℤ}
    (hvn : v ∈ W.grading n) (hv : W.L 1 v = 0) :
    onePoint (contragredientFromResidual N) v = (-1 : ℂ) ^ n.toNat • onePoint N v :=
\end{leancode}

\begin{dictionary}
\textbf{The Jacobi identity is proved, not assumed.} \lean{ContragredientJacobi N} is word for word the module
Jacobi axiom for the field \lean{YMdual N}, and \lean{contragredientJacobi_holds} proves it for every $N$ with no
hypothesis. \lean{contragredientModule N} is the unconditional ordinary module; the \lean{rfl} lemma
\lean{contragredientModule_}\allowbreak\lean{eq_}\allowbreak\lean{contragredient} shows it is \emph{definitionally} the module whose field is \lean{YMdual N}
and whose weight spaces are $M_\alpha^*$ (the remaining fields of \lean{contragredientOrdVModule}: $L_0$ acts by $\alpha$
on $M_\alpha^*$, finite-dimensionality, lower bound --- proved in the same file).
\lean{contragredientFromResidual N} is another name for \lean{contragredientModule N}.

\textbf{Characters.} \lean{contragredientGradedDim N} is $e\mapsto\dim(M_{e+c/24})^*$, and the theorem says this equals
$\dim M_{e+c/24}$ --- a statement about vector spaces, independent of the action (so it holds equally for the module just
built). \lean{onePoint_}\allowbreak\lean{contragredient_}\allowbreak\lean{holds}: for $v\in V_n$ with $L_1v = 0$, the one-point function of $M'$ is
$(-1)^n$ times that of $M$ (\lean{n.toNat} equals $n$ for $n\ge0$; for $n<0$, $v = 0$ by CFT type and both sides vanish).
For general $v$ the formula has $L_1$-correction terms, which the library does not state.

\textbf{Same as the literature?} Yes: FHL Theorem~5.2.1 (and the evident equality $\chi_{M'} = \chi_M$, Gannon (5.3.4b)), and Zhu's
$\tr_{M'}o(v)q^{L_0} = (-1)^n\tr_Mo(v)q^{L_0}$ for quasi-primary $v$ (the transpose of $o(v)$ is $(-1)^n o(v)$ up to
$L_1$-terms). What the library does \emph{not} prove: $M''\cong M$, and that $M'$ is irreducible when $M$ is.
\end{dictionary}

\textbf{Proof on paper.} Vacuum and truncation are immediate from the formula ($\ltwelvevac'_{(n)}$ is the transpose of
$\ltwelvevac_{(-n-2)} = \delta_{n,-1}$; and $u'_{(n)}$ maps $M_\gamma^*$ to $M^*_{\gamma+\wt u-n-1}$, which is $0$ for $n$ large by the
lower bound of $M$). The Jacobi identity is FHL's argument: one must show that the conjugation
$u\mapsto e^{zL_1}(-z^{-2})^{L_0}u$, $z\mapsto z^{-1}$ transforms the Jacobi identity of $Y$ into that of $Y'$. The key
identities are the conjugation formulas $e^{zL_1}(-z^{-2})^{L_0}\,Y(u,z_0)\,(-z^{-2})^{-L_0}e^{-zL_1} =
Y\big(e^{(z+z_0)L_1}(-(z+z_0)^{-2})^{L_0}u, -z_0/(z(z+z_0))\big)$ and the change of variables in the formal delta
functions. The library carries this out mode by mode: it reduces the identity to a residual (\lean{DualAdjResidual}),
proves the quasi-primary case, and then the general case by eliminating the comultiplication coefficient of
$e^{zL_1}$ in closed form and a Vandermonde convolution (\leanpath{ContragredientClose.lean}). For the one-point
function: on $M_\alpha^*$ the zero mode of a quasi-primary $v\in V_n$ is $(-1)^n$ times the transpose of $o(v)|_{M_\alpha}$,
and a transpose has the same trace. \qed

\begin{status}
Proved, unconditional. \texttt{\#print axioms} (run for this chapter) of \lean{VOA.contragredientModule},
\lean{VOA.contragredientJacobi\_holds}, \lean{VOA.contragredient\_gradedDim},
\lean{VOA.onePoint\_contragredient\_holds}: standard three; no \lean{VOA.pbw}. The file
\leanpath{Contragredient.lean} still describes the action as ``frontier'' in one of its section comments; the
status paragraph at the top of the same file says it is unconditional, which is correct.
\end{status}

\begin{physics}
Charge conjugation preserves the spectrum: $M'$ has the same energies and degeneracies as $M$. The one-point function of
a quasi-primary field of weight $n$ picks up $(-1)^n$, the parity of the field under the inversion that exchanges
particles and antiparticles; for a current ($n = 1$) it is the statement that conjugate sectors carry opposite charges.
\end{physics}

\begin{warning}
\textbf{\lean{SelfContragredient} is vacuous (serious, recorded in the library).} In
\leanpath{VOA/\allowbreak Modules/\allowbreak Contragredient.lean}:
\begin{leancode}
def SelfContragredient (N : OrdVModule W M) : Prop :=
  Nonempty (M ≃ₗ[ℂ] RestrictedDual N)
\end{leancode}
This asks only for a \emph{linear} isomorphism $M\cong\bigoplus_\alpha M_\alpha^*$, which exists for every ordinary module
(choose a basis of each finite-dimensional $M_\alpha$). The library proves it:
\lean{selfContragredient_of_ordVModule (N : OrdVModule W M) : SelfContragredient N}
(\leanpath{VOA/\allowbreak Modules/\allowbreak StronglyRational.lean}; \texttt{\#print axioms}: standard three). The
literature's ``$M\cong M'$'' means an isomorphism \emph{of $V$-modules} (with \lean{contragredientModule N}), which the
library does not define. The predicate has no consumer; the library's strong-rationality notions deliberately use
\lean{SelfDual} (Section~\ref{l12:sec:strong}) instead.
\end{warning}

\begin{proposition}[\lean{Unitary} is unsatisfiable in dimension $\ge2$]\label{l12:prop:unitary}
The library's ``unitary module'' structure on an ordinary module $M$ asks for a $\C$-bilinear form $(\cdot|\cdot)$ on $M$
with $(u_{(n)}x|y) = (x|u_{(n)}y)$ for all $u,n,x,y$ and $(x|x)\neq0$ for every $x\neq0$. If $M$ contains two linearly
independent vectors, no such structure exists.
\end{proposition}
\noindent File \leanpath{VOA/\allowbreak Modules/\allowbreak Contragredient.lean}; the proposition is the library theorem
\lean{VOA.}\allowbreak\lean{old_}\allowbreak\lean{unitary_}\allowbreak\lean{isEmpty} (\leanpath{VOA/\allowbreak Modules/\allowbreak CoreFixUnitary.lean}).
\begin{leancode}
structure Unitary [StarRing ℂ] (N : OrdVModule W M) where
  form : M →ₗ[ℂ] M →ₗ[ℂ] ℂ
  invariant : ∀ (u : 𝒱) (n : ℤ) (x y : M),
    form (N.toWeakModule.act u n x) y = form x (N.toWeakModule.act u n y)
  posDef : ∀ x : M, x ≠ 0 → form x x ≠ 0
\end{leancode}

\begin{dictionary}
\textbf{Decoding.} \lean{M →ₗ[ℂ] M →ₗ[ℂ] ℂ} is a \emph{bilinear} (not sesquilinear) form. The hypothesis
\lean{[StarRing ℂ]} asks for a star operation on $\C$ but nothing in the structure uses it; the docstring's ``star
involution $u\mapsto u^*$'' does not appear in the fields. \lean{posDef} is not positivity (which would need a real
part) but non-isotropy: $(x|x)\neq0$.

\textbf{What is wrong (serious).} Over $\C$, a quadratic form in two or more variables always has a nonzero zero:
$q(ae_1+e_2) = a^2(e_1|e_1) + a\big((e_1|e_2)+(e_2|e_1)\big) + (e_2|e_2)$ is a quadratic polynomial in $a$, which has a
root if $(e_1|e_1)\neq0$, and if $(e_1|e_1) = 0$ then $x = e_1$ itself is isotropic. So \lean{posDef} fails on every
space of dimension $\ge2$, whatever the action. Independently, the \lean{invariant} field is not the unitarity
condition of the literature, which is $(u_{(n)}x|y) = (x|\theta(u)^\dagger_{(n)}y)$ with an \emph{antilinear}
involution $\theta$ and, for the Virasoro modes, $(L_nx|y) = (x|L_{-n}y)$; the library's condition pairs $u_{(n)}$
with itself. Unitarity in the sense of Dong--Lin (Hermitian, positive definite, $\theta$-twisted invariance) is
therefore not represented by this structure. It has no consumers (the only use is an \lean{example} that it
typechecks); graded unitarity in the sense of the rest of the library is a different notion
(Chapter~\ref{ch:L19}).
\end{dictionary}

\textbf{Proof on paper.} The isotropic-vector argument just given; the Lean proof formalises it with Mathlib's
quadratic-formula lemma \lean{exists_quadratic_eq_zero} and the existence of square roots in $\C$. \qed

\begin{status}
Proved in the library (\lean{old_unitary_isEmpty}). When this chapter was written the statement was checked in Lean
outside the library, compiled with \texttt{lake env lean -DmaxSynthPendingDepth=3}; axioms \lean{[propext, Classical.choice,
Quot.sound]}. \texttt{\#print axioms} for the library theorem was not re-run for this edition.
\end{status}

\begin{physics}
In physics a unitary sector carries a positive-definite \emph{Hermitian} inner product for which $L_n^\dagger = L_{-n}$
(and $J^a_n{}^\dagger = J^a_{-n}$ for currents); unitarity of the minimal models singles out the FQS series. The library's
structure, being bilinear and asking each mode to be its own adjoint, does not express this, and over the complex
numbers no bilinear form can be ``positive'' on a space of dimension two or more.
\end{physics}

%======================================================================
\section{Intertwining operators and the fusion coefficient}\label{l12:sec:fusion}
%======================================================================

\begin{definition}[The library's intertwining operators and intertwiner space]\label{l12:def:intertwiner}
Let $M,N,P$ be ordinary $V$-modules. \emph{Mode data} of type $\binom PM\,N$ is a family of bilinear maps
$\ltwelveY_s:M\otimes N\to P$, $s\in\C$ (the coefficient of $z^{-s-1}$ of a would-be field $\ltwelveY(u,z)x$ with arbitrary
complex powers of $z$). Its \emph{intertwining defect} is
\begin{equation*}
 \Delta(a,p,s,u,x) = a^P_{(p)}\ltwelveY_s(u)x - \ltwelveY_s(u)a^N_{(p)}x - \sum_{j\ge0}\binom pj\ltwelveY_{s+p-j}(a^M_{(j)}u)x ,
\end{equation*}
for $a\in V$, $p\in\Z$, $s\in\C$, $u\in M$, $x\in N$. The library's \emph{intertwiner space} $\mathrm{Int}(M,N;P)$ is the subspace of mode data
with $\Delta\equiv0$; an \lean{Intertwiner} is an element of it which in addition is lower truncated ($\ltwelveY_s(u)x = 0$ for
$\mathrm{Re}\,s\gg0$); and $\mathrm{fusionCoeff} = \dim\mathrm{Int}(M,N;P)$.
\end{definition}
\noindent File \leanpath{VOA/\allowbreak Modules/\allowbreak Fusion.lean}.
\begin{leancode}
abbrev IntertwinerModes (M N P : Type u) [AddCommGroup M] [Module ℂ M] [AddCommGroup N]
    [Module ℂ N] [AddCommGroup P] [Module ℂ P] : Type u := ℂ → (M →ₗ[ℂ] N →ₗ[ℂ] P)

noncomputable def intertwinerSummand (A : OrdVModule W M) (Y : IntertwinerModes M N P)
    (a : 𝒱) (p : ℤ) (s : ℂ) (u : M) (x : N) (j : ℕ) : P :=
  cbz p j • Y (s + p - j) (A.toWeakModule.act a (j : ℤ) u) x

noncomputable def intertwiningDefect (A : OrdVModule W M) (B : OrdVModule W N) (C : OrdVModule W P) :
    IntertwinerModes M N P →ₗ[ℂ] (𝒱 → ℤ → ℂ → M → N → P) where
  toFun Y := fun a p s u x =>
    C.toWeakModule.act a p (Y s u x) - Y s u (B.toWeakModule.act a p x)
      - ∑ᶠ j : ℕ, intertwinerSummand A Y a p s u x j

noncomputable def IntertwinerSpace (A : OrdVModule W M) (B : OrdVModule W N) (C : OrdVModule W P) :
    Submodule ℂ (IntertwinerModes M N P) :=
  LinearMap.ker (intertwiningDefect A B C)

structure Intertwiner (A : OrdVModule W M) (B : OrdVModule W N) (C : OrdVModule W P) where
  modes : IntertwinerModes M N P
  truncation : ∀ (u : M) (x : N), ∃ s₀ : ℝ, ∀ s : ℂ, s₀ ≤ s.re → modes s u x = 0
  isIntertwiner : modes ∈ IntertwinerSpace A B C
\end{leancode}
(\lean{VOA.IntertwinerModes}, \lean{VOA.intertwinerSummand}, \lean{VOA.intertwiningDefect}, \lean{VOA.IntertwinerSpace},
\lean{VOA.Intertwiner}; the two linearity fields of \lean{intertwiningDefect} are proofs and are omitted. The last
declaration, \lean{fusionCoeff}, is displayed in Theorem~\ref{l12:thm:fusionzero}.)

\begin{dictionary}
\lean{IntertwinerModes M N P} is the type of functions $\C\to(M\to_{\C}N\to_{\C}P)$; \lean{Y s u x} is $\ltwelveY_s(u)x$, the
coefficient of $z^{-s-1}$. \lean{A.toWeakModule.act a j u} is $a^M_{(j)}u$ (Definition~\ref{l12:def:weakmodule}), and
\lean{intertwinerSummand} is the single term $\binom pj\ltwelveY_{s+p-j}(a^M_{(j)}u)x$ of the commutator sum; the sum over $j\in\N$ is a
\lean{finsum}, which is finite because $a^M_{(j)}u = 0$ for $j\gg0$ (the library proves
\lean{intertwinerSummand_support_finite}), so it is the true sum. The defect is exactly the commutator formula
$[a_{(p)},\ltwelveY(u,z)] = \sum_j\binom pj z^{p-j}\ltwelveY(a_{(j)}u,z)$ written in modes: it is the $\ell=0$ slice of the Jacobi
identity for intertwining operators of Frenkel--Huang--Lepowsky. \lean{IntertwinerSpace A B C} is the kernel of the linear map
\lean{intertwiningDefect A B C}, so it is a subspace of the mode data.

\textbf{Same, stronger or weaker.} \emph{Weaker} than FHL's Definition~5.4.1: there is no iterate (associativity) part of the Jacobi
identity, no $L_{-1}$-derivative property $\ltwelveY(L_{-1}u,z) = \partial_z\ltwelveY(u,z)$, and \lean{IntertwinerSpace} does not
even include the lower truncation (only the bundled \lean{Intertwiner} does). So the library's intertwiner space \emph{contains} the space
of FHL intertwiners, possibly strictly. \textbf{Could the definition be wrong or vacuous?} It is not vacuous as a definition: the Moonshine
construction's skew intertwiner $W$ is an \lean{Intertwiner} (Chapter~\ref{ch:09}). But the number \lean{fusionCoeff} built from it is
\emph{not} a fusion rule, and in fact is $0$ for all modules (Theorem~\ref{l12:thm:fusionzero}); the discussion in that theorem applies.
\end{dictionary}

\begin{physics}
An intertwining operator is the algebraic form of a chiral vertex operator of CFT: the operator $\ltwelveY(u,z)x$ that inserts a primary field
of sector $M$ and acts on states of sector $N$, landing in sector $P$. The operator product of fields from two sectors
decomposes as the fusion product $M\times N = \bigoplus_PN_{MN}^P\,P$, and $N^P_{MN}$ is the dimension of the space of intertwiners. The library's notion
keeps only the commutator relations with the modes (the Ward identities for $V$) and drops associativity and the derivative property; so
its dimension is an upper bound on the true fusion multiplicity and is not the fusion rule. The OPE of sectors becomes rigid only
after the full Jacobi identity is imposed, which is what leads to the Verlinde formula.
\end{physics}

\textbf{Proof on paper.} There is nothing to prove: these are definitions. The one claim made, that the vanishing of the defect is the commutator formula in modes,
is the following: for an honest intertwining operator the $\ell=0$ slice of its Jacobi identity reads
$[a_{(p)},\ltwelveY(u,z)] = \sum_{j\ge0}\binom pj z^{p-j}\ltwelveY(a_{(j)}u,z)$, and comparing the coefficients of $z^{-s-1}$ on both sides gives
$a^P_{(p)}\ltwelveY_s(u)x - \ltwelveY_s(u)a^N_{(p)}x = \sum_j\binom pj\ltwelveY_{s+p-j}(a^M_{(j)}u)x$, i.e.\ $\Delta = 0$. \qed

\begin{status}
Definitions. \texttt{\#print axioms} of \lean{VOA.IntertwinerModes}, \lean{VOA.intertwiningDefect}, \lean{VOA.IntertwinerSpace} is
\lean{[propext, Classical.}\allowbreak\lean{choice, Quot.}\allowbreak\lean{sound]} (run for this chapter). \textbf{Warning:}
the derived number \lean{fusionCoeff} is identically $0$ (Theorem~\ref{l12:thm:fusionzero}); do not read it as a fusion rule.
\end{status}

Definition~\ref{l12:def:intertwiner} introduced the library's intertwining operators: mode data
$\ltwelveY_s:M\otimes N\to P$, $s\in\C$, subject only to the commutator condition
\begin{equation*}
 a^P_{(p)}\ltwelveY_s(u)x - \ltwelveY_s(u)a^N_{(p)}x = \sum_{j\ge0}\binom pj\ltwelveY_{s+p-j}(a^M_{(j)}u)x
\end{equation*}
for all $a\in V$, $p\in\Z$, $s\in\C$, $u\in M$, $x\in N$ (the kernel \lean{IntertwinerSpace A B C} of the linear map
\lean{intertwiningDefect}), and \lean{fusionCoeff A B C := Module.finrank ℂ (IntertwinerSpace A B C)}. In the
literature (FHL Definition~5.4.1) an intertwining operator of type $\binom{P}{M\,N}$ must in addition satisfy lower
truncation, the full Jacobi identity (with its iterate part) and the $L_{-1}$-derivative property
$\ltwelveY(L_{-1}u,z) = \partial_z\ltwelveY(u,z)$; the fusion rule $N_{MN}^P$ is the dimension of that space.

\begin{theorem}[\lean{fusionCoeff} is identically zero]\label{l12:thm:fusionzero}
For all ordinary modules $A,B,C$, \lean{fusionCoeff A B C = 0}. More precisely, \lean{IntertwinerSpace A B C} is either
the zero space or infinite-dimensional.
\end{theorem}
\noindent Library file \leanpath{VOA/\allowbreak Modules/\allowbreak Fusion.lean}; the result is the library theorem
\lean{VOA.}\allowbreak\lean{fusionCoeff_}\allowbreak\lean{eq_}\allowbreak\lean{zero} (\leanpath{VOA/\allowbreak Modules/\allowbreak CoreFixFusion.lean}).
\begin{leancode}
noncomputable def fusionCoeff (A : OrdVModule W M) (B : OrdVModule W N) (C : OrdVModule W P) : ℕ :=
  Module.finrank ℂ (IntertwinerSpace A B C)
\end{leancode}

\begin{dictionary}
\textbf{Status of the claim.} Chapter~\ref{ch:01} and Chapter~\ref{ch:15} argued this on paper, and it
had been recorded as a defect that was not yet formalised. It is now formalised, and the defect is \emph{confirmed}. When this chapter was written it was
checked in Lean outside the library, with the standard three axioms; \texttt{\#print axioms} for the library theorem was not
re-run for this edition.

\textbf{Why this is a defect and not a theorem about fusion (serious).} The true fusion rule $N_{VV}^V$ of a simple VOA is $1$,
and $N_{MM'}^V = 1$; the library's number is $0$ in all cases. The cause is the missing $L_{-1}$-derivative property and
grading: without them the exponents $s$ can be shifted by any complex number.

\textbf{Does this make any theorem vacuous?} We searched the library: \lean{fusionCoeff} and \lean{IntertwinerSpace} are used
nowhere outside \leanpath{Fusion.lean} (the identically named \lean{fusionCoeff} in
\leanpath{VOA/\allowbreak Examples/\allowbreak VirasoroFusion*.lean} is an unrelated local power-series coefficient). The
bundled \lean{Intertwiner} structure is used by the Moonshine construction (\lean{skewIntertwiner},
\lean{leechSkewIntertwiner}, Chapter~\ref{ch:09}); its fields are satisfiable (the skew intertwiner is a nonzero
example), so no theorem stated with it is vacuous, and the Moonshine chapters prove the $L_{-1}$-property and all
uniqueness statements separately. The ``Verlinde apex'' in the docstring of \leanpath{Fusion.lean} (``\lean{fusionCoeff} is
exactly the interface they wire to'') is not implemented and, with this definition, could not be (documentation only).
\end{dictionary}

\textbf{Proof on paper.} Let $\ltwelveY\neq0$ be in the space; pick $s_0,u,x$ with $\ltwelveY_{s_0}(u)x\neq0$. For $c\in\C$ let
$\ltwelveY^{(c)}_s = \ltwelveY_{s-c}$ if $s\in s_0+c+\Z$ and $0$ otherwise (restrict to one coset, then shift). The defect
condition at $s$ involves only $\ltwelveY_s$ and $\ltwelveY_{s+p-j}$ with $p-j\in\Z$, i.e.\ values on the coset of $s$, with
coefficients independent of $s$; so on the coset $s_0+c+\Z$ the condition for $\ltwelveY^{(c)}$ at $s$ is the condition for
$\ltwelveY$ at $s-c$, and off that coset every term is $0$. Hence $\ltwelveY^{(c)}$ lies in the space. Take $c_k = ik$,
$k\in\N$. Then $\ltwelveY^{(c_l)}_{s_0+ik}(u)x$ equals $\ltwelveY_{s_0}(u)x\neq0$ if $k = l$ and $0$ if $k\neq l$ (the imaginary
parts differ by a nonzero integer, so the points are in different cosets of $\Z$). Evaluating a vanishing finite linear
combination $\sum_l g_l\ltwelveY^{(c_l)}$ at $s_0+ik$ gives $g_k = 0$: the family is linearly independent and infinite.
By Mathlib's convention the \lean{finrank} of an infinite-dimensional space is $0$; and if the space is zero its
\lean{finrank} is $0$ too. \qed

\begin{status}
Proved in the library (\lean{fusionCoeff_eq_zero}). When this chapter was written the statement was checked in Lean
outside the library (\texttt{lake env lean -DmaxSynthPendingDepth=3}), axioms standard three. The definitions
\lean{IntertwinerSpace}, \lean{Intertwiner}, \lean{fusionCoeff} themselves are axiom-clean.
\end{status}

\begin{physics}
The fusion rule $N_{MN}^P$ is the multiplicity of the sector $P$ in the operator product of fields from sectors $M$ and $N$:
$\phi_M\times\phi_N = \sum_PN_{MN}^P\phi_P$. It is finite because a chiral vertex operator is pinned down by its leading
coefficient: the $L_{-1}$-derivative property fixes the $z$-dependence ($\ltwelveY(u,z)\sim z^{h_P-h_M-h_N}$ times a series),
and the grading fixes the exponents. If one keeps only the commutation relations with the modes, the exponent $z^{s}$ can be
multiplied by an arbitrary $z^c$ --- a ``fake'' intertwiner for each $c$ --- and the count becomes infinite. That is what the
theorem detects.
\end{physics}

%======================================================================
\section{The rationality hierarchy}\label{l12:sec:rationality}
%======================================================================

A finite group has finitely many irreducible representations and every representation is completely reducible.
``Rational'' VOAs are those whose module theory looks like that. Unfortunately the literature uses the word with
several inequivalent meanings, depending on which module category is required to be semisimple and which
finiteness conditions are added. The library names each predicate after the category it quantifies over and
records, for each implication, whether it is \emph{proved}, \emph{carried as a named citation}, or \emph{not
available}. This section presents the predicates, the proved implications, and the comparison with the literature.

\subsection{Simplicity and the set \texorpdfstring{$\Phi(V)$}{Phi(V)} of irreducible modules}

\begin{definition}[Simple VOA; irreducible modules and their isomorphism classes]\label{l12:def:phi}
$V$ is \emph{simple} if $V\neq0$ and its only ideals are $0$ and $V$. An isomorphism of ordinary modules is a linear
isomorphism intertwining all modes. $\Phi(V)$ is the set of isomorphism classes of irreducible ordinary modules.
\end{definition}
\noindent Files \leanpath{VOA/\allowbreak Core/\allowbreak BilinearForm.lean} (Chapter~\ref{ch:L04}),
\leanpath{VOA/\allowbreak Modules/\allowbreak Rationality.lean}, \leanpath{VOA/\allowbreak Modules/\allowbreak StronglyRational.lean}.
\begin{leancode}
def IsSimple (W : VertexOperatorAlgebra 𝒱) : Prop :=
  (⊤ : Submodule ℂ 𝒱) ≠ ⊥ ∧ ∀ I : Submodule ℂ 𝒱, W.IsIdeal I → I = ⊥ ∨ I = ⊤

def IsSimple (W : VertexOperatorAlgebra 𝒱) : Prop :=
  OrdVModule.IsIrreducible (adjointModule W)

theorem isSimple_iff_adjointModule_isIrreducible (W : VertexOperatorAlgebra 𝒱) :
    W.IsSimple ↔ OrdVModule.IsIrreducible (adjointModule W) := by

structure IrrModule (W : VertexOperatorAlgebra 𝒱) where
  carrier : Type u
  [addCommGroup : AddCommGroup carrier]
  [module : Module ℂ carrier]
  isMod : OrdVModule W carrier
  irreducible : OrdVModule.IsIrreducible isMod

structure ModuleEquiv (W : VertexOperatorAlgebra 𝒱) {M N : Type u} [AddCommGroup M]
    [Module ℂ M] [AddCommGroup N] [Module ℂ N] (A : OrdVModule W M) (B : OrdVModule W N) where
  toLinearEquiv : M ≃ₗ[ℂ] N
  map_act : ∀ (u : 𝒱) (n : ℤ) (x : M),
    toLinearEquiv (A.toWeakModule.act u n x) = B.toWeakModule.act u n (toLinearEquiv x)

def Φ (W : VertexOperatorAlgebra 𝒱) : Type (u + 1) := Quotient (IrrModule.setoid W)
\end{leancode}

\begin{dictionary}
\textbf{Two \lean{IsSimple}s.} The first (\lean{VOA.VertexOperatorAlgebra.IsSimple}, written \lean{W.IsSimple}) is the
ideal-theoretic definition; the second (\lean{VOA.IsSimple}, in \leanpath{Rationality.lean}, marked
\lean{@[deprecated]}) is ``$V$ is irreducible as a module over itself''. The library proves them equivalent
(\lean{isSimple_iff_adjointModule_isIrreducible}): a weak submodule of the adjoint module is exactly an ideal. All
rationality predicates below use \lean{W.IsSimple}.

\textbf{$\Phi(V)$.} \lean{IrrModule.setoid} is ``there is a \lean{ModuleEquiv}''; \lean{Quotient} is the set of classes. The
carriers range over \lean{Type u}, the universe of $V$; this is the standard way to make ``the set of all modules up to
isomorphism'' a set, and $\Phi(V)$ itself lives one universe up (\lean{Type (u + 1)}). A \lean{ModuleEquiv} imposes no
condition on gradings because commuting with $L_0$ already forces it to preserve weight spaces
(\lean{ModuleEquiv.maps_weight}). Same as Gannon's $\Phi(V)$.
\end{dictionary}

\begin{status}
Definitions; \texttt{\#print axioms VOA.isSimple\_iff\_adjointModule\_isIrreducible}: standard three.
\end{status}

\begin{physics}
Simplicity of $V$ means the vacuum sector has no null fields: every nonzero field generates all others, the
vacuum is ``cyclic and separating''. $\Phi(V)$ is the set of primary fields (equivalently, of irreducible sectors) of
the chiral CFT --- for the Ising model, $\{1,\sigma,\varepsilon\}$; for $\Vnat$, a single element.
\end{physics}

\subsection{Admissible rationality, and the submodule convention}\label{l12:warn:admcr}

\begin{definition}[Rational; Zhu 1996, Dong--Li--Mason 1997--98]\label{l12:def:admrat}
$V$ is \emph{(admissibly) rational} if every admissible $V$-module is completely reducible.
\end{definition}
\noindent File \leanpath{VOA/\allowbreak Modules/\allowbreak AdmissibleModule.lean}.
\begin{leancode}
def AdmissiblyRational (W : VertexOperatorAlgebra 𝒱) : Prop :=
  ∀ {M : Type u} [AddCommGroup M] [Module ℂ M] (N : AdmissibleModule W M),
    N.IsCompletelyReducible

abbrev IsRational (W : VertexOperatorAlgebra 𝒱) : Prop := AdmissiblyRational W
\end{leancode}

\begin{dictionary}
Quantifies over all admissible modules with carrier in \lean{Type u}, and asks every \emph{weak} submodule to have a weak
complement (Definition~\ref{l12:def:admissible}). \lean{IsRational} is a pure alias. \textbf{Literature.} Dong--Li--Mason's
rationality (DLM, \emph{Math.\ Ann.} 310, and \emph{J. Algebra} 206) is the single clause ``every admissible module is a direct
sum of (graded-)simple admissible modules''. Zhu's own Definition~1.2.4 of ``rational'' had three clauses (finitely many
irreducibles, finite-dimensional graded pieces, complete reducibility); DLM Theorem~8.1 shows the extra two follow from
the single clause. The library follows DLM.
\end{dictionary}

\begin{physics}
A rational CFT is one whose Hilbert space decomposes into finitely many irreducible sectors of the chiral algebra,
$\mathcal H = \bigoplus_{i,j}N_{ij}M_i\otimes\overline{M_j}$; complete reducibility of all (admissible) modules is the
algebraic statement that there are no indecomposable, non-split sectors --- the hallmark of non-logarithmic theories.
\end{physics}

\begin{warning}
\textbf{Weak versus graded submodules; a self-contradictory docstring (documentation only).} Because complete reducibility is phrased
with weak submodules, \lean{AdmissiblyRational} is not literally DLM's notion. Consider the trivial VOA $V = \C\ltwelvevac$
and $M = \C e_0\oplus\C e_1$ with $M(0) = \C e_0$, $M(1) = \C e_1$. The line $\C(e_0+e_1)$ is a weak submodule that is not
graded. The library's predicate demands a complement for it (it has one); DLM's never mentions it. Which is stronger? The
library's docstrings disagree with each other: \leanpath{AdmissibleRationality.lean} says, in one paragraph, that the
weak reading ``\emph{implies}'' the graded one (``Corrected''), and two paragraphs later that ``both halves of that
incomparability are real'', and \leanpath{AdmissibleModule.lean} repeats the ``implies'' claim. \emph{One direction of
the comparison with DLM's definition is now proved in Lean}: for $C_2$-cofinite $V$, DLM-rational $\Rightarrow$
\lean{AdmissiblyRational} (Theorem~\ref{l12:thm:dlmrat} below). The converse is open. What is proved in the other
direction runs from the library's notion towards DLM's:
\lean{admissiblyRationalGraded_of_admissiblyRational} (\lean{AdmissiblyRational} implies the level-graded reading
\lean{AdmissiblyRationalGraded}: every level-graded weak submodule has a level-graded complement), and, under $C_2$,
\lean{AdmissiblyRational.}\allowbreak\lean{exists_}\allowbreak\lean{iSupIndep_}\allowbreak\lean{irreducible_}\allowbreak\lean{ordinary}
(every admissible module is an internal direct sum of irreducible ordinary modules, the conclusion of DLM Theorem~8.1).
That the graded reading implies DLM's ``every admissible module is a direct sum of irreducibles'', and the converse
\lean{AdmissiblyRationalGraded} $\Rightarrow$ \lean{AdmissiblyRational}, are not formalised. On paper: if DLM-rationality holds, then by DLM Theorem~8.1 every admissible module is a direct sum of simple
\emph{ordinary} modules, each weak-simple, hence semisimple as a module over the algebra of modes, hence every weak submodule
is complemented --- so DLM-rational $\Rightarrow$ \lean{AdmissiblyRational}, using DLM's theorem. The Lean proof of
Theorem~\ref{l12:thm:dlmrat} avoids DLM's theorem and uses $C_2$ instead. The converse requires producing
graded simple summands from weak complements, which is not formalised. Practical consequences: using
\lean{AdmissiblyRational} as a \emph{hypothesis} is now justified for every $C_2$-cofinite DLM-rational $V$ (by
Theorem~\ref{l12:thm:dlmrat}); \emph{applying DLM's theorems} to an \lean{AdmissiblyRational} $V$ still needs the open
converse, and a proved \lean{¬ AdmissiblyRational W} (e.g.\ for the triplet
$W(2)$) is formally a statement about the library's notion. For the triplet the irrationality comes from an
indecomposable \emph{ordinary} module, where weak and graded submodules coincide (Section~\ref{l12:sec:categories}), so
it is irrational in every sense.
\end{warning}

\begin{theorem}[DLM rationality implies \lean{AdmissiblyRational}, under $C_2$]\label{l12:thm:dlmrat}
Let $V$ be $C_2$-cofinite and rational in the sense of Dong--Li--Mason: every admissible module is the internal direct sum
of graded submodules that are minimal among nonzero graded submodules. Then $V$ is \lean{AdmissiblyRational}.
\end{theorem}
\noindent File \leanpath{VOA/\allowbreak Classification/\allowbreak COne/\allowbreak DLMRationality.lean}, namespace
\lean{VOA.Classification.COne.DLMRat}.
\begin{leancode}
def GradedMinimal (N : AdmissibleModule W M) (P : WeakModule.Submodule' N.toWeakModule) : Prop :=
  N.IsLevelGraded P ∧ P.carrier ≠ ⊥ ∧
    ∀ Q : WeakModule.Submodule' N.toWeakModule, N.IsLevelGraded Q → Q.carrier ≤ P.carrier →
      Q.carrier = ⊥ ∨ Q.carrier = P.carrier

def DLMRational (W : VertexOperatorAlgebra 𝒱) : Prop :=
  ∀ {M : Type u} [AddCommGroup M] [Module ℂ M] (N : AdmissibleModule W M),
    ∃ (ι : Type u) (P : ι → WeakModule.Submodule' N.toWeakModule),
      iSupIndep (fun i => (P i).carrier) ∧ ⨆ i, (P i).carrier = ⊤ ∧ ∀ i, GradedMinimal N (P i)

theorem admissiblyRational_of_dlmRational (hC2 : C2Cofinite W) (h : DLMRational W) :
    AdmissiblyRational W := by
\end{leancode}

\begin{dictionary}
\lean{DLMRational} is a faithful transcription of DLM's definition: admissible submodules are the graded ones, and
irreducibility of a summand in its induced grading is exactly graded minimality inside $N$ (independent release audit,
\texttt{Notes/\allowbreak Textbook/\allowbreak referee/\allowbreak RELEASE\_AUDIT\_QUEUE2.md}). It does not require
$M(0)\neq0$, which is harmless, and quantifies over carriers in \lean{Type u}, like \lean{AdmissiblyRational}.
\textbf{Proof idea.} The cyclic submodule of a homogeneous vector is level-graded (\lean{cyclic_isLevelGraded}, no $C_2$),
so a graded-minimal $P$ is generated by any nonzero homogeneous vector in it. Under $C_2$, $L(0)$ acts on level $j$ by
$h+j$, so a weak submodule of $P$ is $L(0)$-stable with distinct eigenvalues on the levels, hence level-graded, hence
$0$ or $P$: graded-minimal implies weakly simple (\lean{isSimpleSub_of_gradedMinimal}). A direct sum of weakly simple
modules is a semisimple module over the mode algebra, so every weak submodule is complemented. \textbf{Not an
equivalence:} the converse \lean{AdmissiblyRational} $\Rightarrow$ \lean{DLMRational} is open. \textbf{Extra hypothesis:}
DLM's notion needs no $C_2$; the theorem does.
\end{dictionary}

\begin{status}
Proved, conditional on \lean{hC2}. \texttt{\#print axioms}: \lean{[propext, Classical.}\allowbreak\lean{choice, Quot.}\allowbreak\lean{sound]},
recorded in \texttt{Notes/\allowbreak release/\allowbreak WS-RELEASE2.md} (commit \texttt{34b74b85}); the aggregate
\texttt{lake build} at the release commit is pending.
\end{status}

\begin{physics}
Physicists use ``rational'' loosely for ``finitely many primaries, every sector a direct sum of irreducibles''. The two
mathematical readings compared here differ only in which sub-sectors one asks to split off: DLM ask it of energy-graded
sub-sectors, the library of all subspaces closed under the modes. The theorem says that once $V$ is $C_2$-cofinite
(lisse: its semiclassical phase space is a point, Chapter~\ref{ch:L16}) the two agree in one direction: a theory with no
non-split graded sectors has no non-split sectors at all, because $L_0$ separates the energy levels inside an irreducible
sector. The converse direction is open, and without $C_2$ nothing is claimed.
\end{physics}

\begin{definition}[Gannon-rational and holomorphic; Gannon, Def.~5.3.2]\label{l12:def:gannon}
$V$ is \emph{Gannon-rational} if (1) every ordinary module is completely reducible, (2) $\Phi(V)$ is finite, and (3) every
irreducible admissible module is ordinary (on the same weak module). $V$ is \emph{holomorphic} if it is Gannon-rational
and $|\Phi(V)| = 1$.
\end{definition}
\noindent File \leanpath{VOA/\allowbreak Modules/\allowbreak Rationality.lean}.
\begin{leancode}
def GannonAdmissibleIsOrdinary (W : VertexOperatorAlgebra 𝒱) : Prop :=
  ∀ {M : Type u} [AddCommGroup M] [Module ℂ M] (N : AdmissibleModule W M),
    N.IsIrreducible → ∃ O : OrdVModule W M, O.toWeakModule = N.toWeakModule

structure GannonRational (W : VertexOperatorAlgebra 𝒱) : Prop where
  completelyReducible : ∀ {M : Type u} [AddCommGroup M] [Module ℂ M] (N : OrdVModule W M),
    OrdVModule.IsCompletelyReducible N
  finitelyManyIrred : Finite (Φ W)
  admissibleIsOrdinary : GannonAdmissibleIsOrdinary W

def Holomorphic (W : VertexOperatorAlgebra 𝒱) : Prop :=
  GannonRational W ∧ Nat.card (Φ W) = 1
\end{leancode}

\begin{dictionary}
\lean{∃ O : OrdVModule W M, O.toWeakModule = N.toWeakModule}: the \emph{given} weak module carries an ordinary grading.
\lean{Finite (Φ W)}: finitely many classes. \lean{Nat.card (Φ W) = 1}: exactly one class (\lean{Nat.card} of an infinite
type is $0$, so this also forces finiteness). Gannon calls this ``weakly rational''; his ``weak module'' (Def.~5.3.1) is
$\N$-graded, i.e.\ admissible, so clause (3) is over admissible modules --- the library repaired an earlier transcription
(\lean{WeaklyRational}, retired) that used ungraded weak modules and a detached conclusion. \textbf{Comparison:} clause (1)
concerns only ordinary modules, so \lean{GannonRational} does \emph{not} imply \lean{AdmissiblyRational} (Gannon himself
adds admissible complete reducibility in his stronger Def.~6.2.3). \lean{Holomorphic} has a single consumer, a structure
(\lean{MoonshineHypotheses}) with no consumers; the Moonshine theorem of this book does not use it. Note that
``$\Vnat$ is holomorphic'' is not proved in the library.
\end{dictionary}

\begin{physics}
A holomorphic VOA is a chiral CFT with a single sector, the vacuum: its character is then a modular function (up to a
phase), as for $\Vnat$ whose character is $J(q)$. Gannon's three clauses say: sectors don't mix, there are finitely many
primaries, and every reasonable representation has a genuine energy grading.
\end{physics}

\begin{definition}[Regular; Dong--Li--Mason 1997, Def.~3.1]\label{l12:def:regular}
$V$ is \emph{regular} if every weak module is completely reducible and every irreducible weak module is ordinary.
$\Phi_{\mathrm{adm}}(V)$ is the set of isomorphism classes of irreducible admissible modules (isomorphism of the underlying
weak modules).
\end{definition}
\noindent File \leanpath{VOA/\allowbreak Modules/\allowbreak RationalityLattice.lean}.
\begin{leancode}
structure Regular (W : VertexOperatorAlgebra 𝒱) : Prop where
  weakCompletelyReducible : ∀ {M : Type u} [AddCommGroup M] [Module ℂ M] (A : WeakModule W M)
    (P : WeakModule.Submodule' A), ∃ Q : WeakModule.Submodule' A, IsCompl P.carrier Q.carrier
  irreducibleWeakIsOrdinary : ∀ {M : Type u} [AddCommGroup M] [Module ℂ M] (A : WeakModule W M),
    (Nontrivial M ∧ ∀ P : WeakModule.Submodule' A, P.carrier = ⊥ ∨ P.carrier = ⊤) →
    ∃ N : OrdVModule W M, N.toWeakModule = A

structure IrrAdmModule (W : VertexOperatorAlgebra 𝒱) where
  carrier : Type u
  [addCommGroup : AddCommGroup carrier]
  [module : Module ℂ carrier]
  isMod : AdmissibleModule W carrier
  irreducible : isMod.IsIrreducible

def Φadm (W : VertexOperatorAlgebra 𝒱) : Type (u + 1) := Quotient (IrrAdmModule.setoid W)
\end{leancode}

\begin{dictionary}
DLM define regular as ``every weak module is a direct sum of irreducible ordinary modules''. The library's two clauses are
equivalent to that: complete reducibility of a weak module (every submodule complemented) makes it semisimple over the
algebra of modes, hence a direct sum of irreducible weak submodules, each ordinary by the second clause; conversely a
direct sum of irreducibles has complemented submodules. Same as DLM.
\end{dictionary}

\begin{physics}
Regularity says that \emph{every} representation of the chiral algebra, even without any energy grading, splits into
physical sectors. It is the strongest semisimplicity one can ask, and for CFT-type VOAs it is equivalent to rationality
plus $C_2$-cofiniteness (Abe--Buhl--Dong).
\end{physics}

\subsection{The literature theorems, stated as named predicates}\label{l12:sec:litnamed}

The library states each of the deep structure theorems of Dong--Li--Mason, Abe--Buhl--Dong and Li that it uses as a
\lean{Prop}-valued definition, so that a theorem using it can carry it as an explicit hypothesis (never an axiom). This was
how all seven entered the library. Six of them have since been \emph{proved} (Theorems~\ref{l12:thm:zhuss}
and~\ref{l12:thm:litproved} below; two of the six only for $C_2$-cofinite $V$), and only \lean{LiC2CofiniteOfRegular} is
still a genuine, unproved input. The definitions themselves are unchanged.
\begin{leancode}
def DLMSimpleAdmissibleOrdinary (W : VertexOperatorAlgebra 𝒱) : Prop :=
  AdmissiblyRational W →
    ∀ {M : Type u} [AddCommGroup M] [Module ℂ M] (N : AdmissibleModule W M),
      N.IsIrreducible → ∃ O : OrdVModule W M, O.toWeakModule = N.toWeakModule

def DLMFinitelyManyIrredAdmissible (W : VertexOperatorAlgebra 𝒱) : Prop :=
  AdmissiblyRational W → Finite (Φadm W)

def ABDWeakIsAdmissible (W : VertexOperatorAlgebra 𝒱) : Prop :=
  C2Cofinite W → ∀ {M : Type u} [AddCommGroup M] [Module ℂ M] (A : WeakModule W M),
    ∃ N : AdmissibleModule W M, N.toWeakModule = A

def ABDIrreducibleWeakIsOrdinary (W : VertexOperatorAlgebra 𝒱) : Prop :=
  C2Cofinite W → ∀ {M : Type u} [AddCommGroup M] [Module ℂ M] (A : WeakModule W M),
    (Nontrivial M ∧ ∀ P : WeakModule.Submodule' A, P.carrier = ⊥ ∨ P.carrier = ⊤) →
    ∃ N : OrdVModule W M, N.toWeakModule = A

def ABDRegularOfC2AndRational (W : VertexOperatorAlgebra 𝒱) : Prop :=
  C2Cofinite W → AdmissiblyRational W → Regular W

def LiC2CofiniteOfRegular (W : VertexOperatorAlgebra 𝒱) : Prop :=
  Regular W → C2Cofinite W

def ZhuAdmissibleSemisimple (W : VertexOperatorAlgebra 𝒱) : Prop :=
  AdmissiblyRational W → letI := instRingA W; IsSemisimpleRing (A W)
\end{leancode}
(Files \leanpath{VOA/\allowbreak Modules/\allowbreak RationalityLattice.lean} and, for the last,
\leanpath{VOA/\allowbreak Modules/\allowbreak AdmissibleRationality.lean}.)

\begin{dictionary}
Each \lean{def} is the implication asserted by a published theorem, specialised to $V$: DLM \emph{Math.\ Ann.}~310,
Thm.~8.1(c) and 8.1(b) (rational $\Rightarrow$ irreducible admissibles are ordinary; finitely many of them);
Abe--Buhl--Dong, \emph{Trans.\ AMS} 356, Prop.~5.7, Cor.~5.8 and Thm.~4.5; H.~Li, \emph{J. Algebra} 212, Thm.~3.8; and
the semisimplicity of $A(V)$ for rational $V$ (DLM \emph{J.~Algebra} 206, Thm.~4.10; Zhu Thm.~2.2.3 for Zhu's stronger
notion). Fidelity points: (a) the library's \lean{IsIrreducible} for admissible modules is \emph{weak}-simplicity, a subclass of
DLM's graded-simple modules, so \lean{DLMSimpleAdmissibleOrdinary} and \lean{DLMFinitelyManyIrredAdmissible} are
\emph{implied by} DLM's theorems \emph{provided} the library's \lean{AdmissiblyRational} implies DLM-rationality (see the
warning above); (b) all carriers are in \lean{Type u}, which only restricts the conclusions. As long as such a predicate is
an \emph{unproved} hypothesis, a theorem that assumes it is ``conditional on the literature'', and its licence depends on
proviso (a). (c) The same proviso applies to \lean{ZhuAdmissibleSemisimple}. (d) The $C_2$-hypothesis of the three ABD
predicates sits \emph{inside} the predicate (\lean{C2Cofinite W → …}), so proving the predicate for every $V$ is the
literature's ``for every $C_2$-cofinite $V$''; the two DLM predicates have only the rationality antecedent inside, so
proving them under an external hypothesis \lean{hC2 : C2Cofinite W} covers only the $C_2$-cofinite VOAs.

\textbf{Status of each predicate (this revision).}
\begin{itemize}
\item \lean{ZhuAdmissibleSemisimple}: \textbf{proved} for every $V$ (Theorem~\ref{l12:thm:zhuss}; also restated as
  \lean{zhuAdmissibleSemisimple_holds}, Theorem~\ref{l12:thm:litproved}). The first release revision of this chapter called it
  an unproved bridge; that was wrong.
\item \lean{ABDWeakIsAdmissible}, \lean{ABDIrreducibleWeakIsOrdinary}, \lean{ABDRegularOfC2AndRational}: \textbf{proved} for
  every $V$ (\lean{abdWeakIsAdmissible}, \lean{abdIrreducibleWeakIsOrdinary_holds}, \lean{abdRegularOfC2AndRational_holds};
  Theorem~\ref{l12:thm:litproved}).
\item \lean{DLMSimpleAdmissibleOrdinary}, \lean{DLMFinitelyManyIrredAdmissible}: \textbf{proved for every $C_2$-cofinite $V$}
  (\lean{dlmSimpleAdmissibleOrdinary_of_c2Cofinite}, \lean{dlmFinitelyManyIrredAdmissible_of_c2Cofinite};
  Theorem~\ref{l12:thm:litproved}), hence for every strongly rational $V$. \textbf{Not proved} for an admissibly rational $V$
  that is not known to be $C_2$-cofinite; that case is tied to the open conjecture ``rational $\Rightarrow C_2$-cofinite''.
\item \lean{LiC2CofiniteOfRegular}: \textbf{unproved}. No proof of it exists in the library (the only instance is the
  trivial one-dimensional VOA). It is used by exactly one theorem of this chapter,
  \lean{regular_iff_c2Cofinite_and_admissiblyRational} (Theorem~\ref{l12:thm:lattice}, item~(7)), which therefore remains
  conditional on it (its other hypothesis, \lean{hABD}, is now discharged by \lean{abdRegularOfC2AndRational_holds}). No
  other declaration of the library takes \lean{LiC2CofiniteOfRegular} as a hypothesis (search of \texttt{VOA/} at
  the third release revision); the headline theorems of the book assume $C_2$-cofiniteness directly instead.
\end{itemize}
Because the six proved predicates are proved directly from the library's own definitions (\lean{AdmissiblyRational},
\lean{C2Cofinite}), their truth does not depend on the (half-open) comparison of \lean{AdmissiblyRational} with
DLM-rationality.
\end{dictionary}

\begin{physics}
These are the structure theorems of rational CFT: in a rational theory every primary has an honest energy grading, there
are finitely many primaries, lisse + rational is the same as ``every representation splits'', and the zero-mode algebra is a
finite direct sum of matrix algebras --- one block per primary, of size the ground-state degeneracy.
\end{physics}

\subsection{Zhu's semisimplicity theorem, proved}

The first of the predicates above that the library proves is Zhu's semisimplicity theorem; the others follow in
Section~\ref{l12:sec:litproved}.

\begin{theorem}[Zhu 1996, Theorem~2.2.3; Dong--Li--Mason: rational $\Rightarrow A(V)$ semisimple]\label{l12:thm:zhuss}
Let $V$ be a vertex operator algebra of the library. If $V$ is admissibly rational (Definition~\ref{l12:def:admrat}), then
Zhu's algebra $A(V)$, with Zhu's product, is a semisimple ring. Equivalently, the predicate \lean{ZhuAdmissibleSemisimple W}
holds for every $V$.
\end{theorem}
\noindent File \leanpath{VOA/\allowbreak Huang/\allowbreak ZhuSemisimple.lean} (Chapter~\ref{ch:L21}), namespace
\lean{VOA.Huang.ZhuSemisimple}. The section variables are \lean{{W : VertexOperatorAlgebra 𝒱}} and
\lean{(Pr : BottomProjection}\allowbreak\lean{ (parabolicRepOfAMod W (regularAMod W)))}, and each theorem is preceded by
\lean{include Pr in}.
\begin{leancode}
theorem isSemisimpleRing_A_of_admissiblyRational (hrat : AdmissiblyRational W) :
    letI := instRingA W; IsSemisimpleRing (A W) := by

theorem zhuAdmissibleSemisimple_of_Pr : ZhuAdmissibleSemisimple W :=
  fun hrat => isSemisimpleRing_A_of_admissiblyRational Pr hrat
\end{leancode}

\begin{dictionary}
\textbf{Hypotheses.} \lean{hrat : AdmissiblyRational W} is Definition~\ref{l12:def:admrat}: every weak submodule of every
admissible module has a weak complement. \lean{include Pr in} makes the section variable \lean{Pr} an explicit first
argument. \lean{Pr} is the bottom-projection datum of Chapter~\ref{ch:L13}: the regular $A(V)$-module embeds as the degree-zero
part of the module induced from it. It is supplied for every $V$ by \lean{voaPBWBottomProjection}, from the PBW theorem
\lean{VOA.pbw} of Chapter~\ref{ch:L02}, which is a theorem. So the only real hypothesis is \lean{hrat}.

\textbf{Conclusion.} \lean{letI := instRingA W} installs the ring structure on \lean{A W} given by Zhu's product (the ring of
Theorem~\ref{l12:thm:zhuring}), and \lean{IsSemisimpleRing (A W)} is Mathlib's notion: $A(V)$ is semisimple as a left module
over itself, i.e.\ every left ideal has a complementary left ideal. The second theorem is the first, packaged as this chapter's
predicate \lean{ZhuAdmissibleSemisimple W} $=$ \lean{AdmissiblyRational W → letI := instRingA W; IsSemisimpleRing (A W)}.

\textbf{Same as the literature?} Same conclusion as Zhu's Theorem~2.2.3 and Dong--Li--Mason's version, with ``rational'' read
as \lean{AdmissiblyRational}. DLM-rationality implies \lean{AdmissiblyRational}, in Lean for $C_2$-cofinite $V$
(Theorem~\ref{l12:thm:dlmrat}) and on paper in general (warning of Section~\ref{l12:warn:admcr}), so with that step the
Lean theorem implies the literature's statement. The Lean proof uses no
theorem of Dong--Li--Mason as a hypothesis, so the unsettled converse plays no role in it.

\textbf{Could it be vacuous?} \lean{AdmissiblyRational} is satisfiable: the library proves it for the trivial VOA and for the
lattice VOAs $V_{\sqrt{2m}\Z}$ (\lean{V2m_admissiblyRational}, Chapter~\ref{ch:L22}).
\end{dictionary}

\textbf{Proof on paper.} Let $M$ be the generalised Verma module induced from the regular module $A(V)$; it is admissible and,
by \lean{Pr}, its degree-zero part is $A(V)$. Given a left ideal $T$, the submodule $P$ generated by $T$ is graded with
$P\cap M(0)=T$ (bottom recovery). Rationality gives a complement $Q$ of $P$. The degree-zero components of the elements of $Q$
that differ from their degree-zero component by an element of $P$ form a left ideal $T'$, and $A(V)=T\oplus T'$, because zero
modes commute with the projection to degree zero and act on $M(0)$ by left multiplication. The full argument is written out in
Chapter~\ref{ch:L15}, Theorem~\ref{l15:thm:zhuss}. \qed

\begin{physics}
In a rational CFT every representation of the chiral algebra splits into sectors. The theorem says the splitting is
already visible on ground states: the zero-mode algebra $A(V)$ is a direct sum of matrix algebras, one per sector. This is
what Zhu's modularity theorem needs in order to count the torus conformal blocks by the number of sectors.
\end{physics}

\begin{status}
\texttt{\#print axioms} of \lean{VOA.Huang.ZhuSemisimple.}\allowbreak\lean{isSemisimpleRing_}\allowbreak\lean{A_of_}\allowbreak\lean{admissiblyRational} and
\lean{VOA.Huang.ZhuSemisimple.}\allowbreak\lean{zhuAdmissibleSemisimple_of_Pr}: \texttt{[propext, Classical.choice, Quot.sound]}, run
at an earlier commit (not re-run for this edition) with \texttt{lake env lean -DmaxSynthPendingDepth=3} (scratch file; the olean of
\leanpath{ZhuSemisimple.lean} is newer than its source). Proved; conditional on \lean{Pr}, which \lean{voaPBWBottomProjection}
supplies (standard three, same run). In the same run, composing the two gave \lean{ZhuAdmissibleSemisimple W} for every $V$ with
the standard three axioms.
\end{status}

\subsection{The Abe--Buhl--Dong and Dong--Li--Mason predicates, proved}\label{l12:sec:litproved}

Five more of the named predicates of Section~\ref{l12:sec:litnamed} are theorems of the library. Three are the theorems of
Abe--Buhl--Dong on $C_2$-cofinite VOAs, proved for every $V$ (the $C_2$-hypothesis is inside each predicate); two are the
Dong--Li--Mason theorems on rational VOAs, proved for every $C_2$-cofinite $V$. The file
\leanpath{VOA/\allowbreak Classification/\allowbreak COne/\allowbreak LiteratureInputsProved.lean} collects them under single
names. The consequences for strongly rational VOAs are Theorem~\ref{l12:thm:srlit} in Section~\ref{l12:sec:strong}, after
strong rationality has been defined.

\begin{theorem}[Abe--Buhl--Dong Prop.~5.7, Cor.~5.8, Thm.~4.5; Dong--Li--Mason Thm.~8.1(b),(c); Zhu]\label{l12:thm:litproved}
Let $V$ be a vertex operator algebra of the library (of CFT type, Chapter~\ref{ch:01}).
\begin{enumerate}
\item \emph{(ABD Prop.~5.7)} If $V$ is $C_2$-cofinite, then every weak $V$-module $M$ carries an admissible-module structure
  (Definition~\ref{l12:def:admissible}) whose underlying weak module is $M$ itself, with the same vertex operators.
\item \emph{(ABD Cor.~5.8)} If $V$ is $C_2$-cofinite, then every irreducible weak $V$-module $M$ (that is, $M\neq0$ and its only
  weak submodules are $0$ and $M$) carries an ordinary-module structure (Definition~\ref{l12:def:ordmodule}) with the same
  vertex operators.
\item \emph{(ABD Thm.~4.5, hard direction)} If $V$ is $C_2$-cofinite and admissibly rational, then $V$ is regular.
\item \emph{(Zhu)} If $V$ is admissibly rational, then $A(V)$ is a semisimple ring (Theorem~\ref{l12:thm:zhuss}, restated).
\item \emph{(DLM Thm.~8.1(c))} If $V$ is $C_2$-cofinite, then every irreducible admissible $V$-module carries an
  ordinary-module structure with the same vertex operators. No rationality is assumed.
\item \emph{(DLM Thm.~8.1(b))} If $V$ is $C_2$-cofinite and admissibly rational, then $V$ has only finitely many isomorphism
  classes of irreducible admissible modules: $\Phi_{\mathrm{adm}}(V)$ is finite.
\end{enumerate}
\end{theorem}
\noindent File \leanpath{VOA/\allowbreak Classification/\allowbreak COne/\allowbreak ABDWeakAdmissibleComplete.lean}, namespace
\lean{VOA.Classification.COne.ABDWeakAdmissible}, with \lean{W} an implicit section variable (item~1):
\begin{leancode}
theorem abdWeakIsAdmissible : ABDWeakIsAdmissible W :=
  abdWeakIsAdmissible_of_degreeNormalForm_boundedBelow
    weakDegreeNormalForm weakSpectrumBoundedBelow
\end{leancode}
\noindent File \leanpath{VOA/\allowbreak Classification/\allowbreak COne/\allowbreak AdmissibleOrdinaryC2.lean}, namespace
\lean{VOA.Classification.COne.AdmissibleOrdinary}, with \lean{W} implicit (item~5):
\begin{leancode}
theorem irreducible_admissible_isOrdinary (hC2 : C2Cofinite W)
    (N : AdmissibleModule W M) (hirr : N.IsIrreducible) :
    ∃ O : OrdVModule W M, O.toWeakModule = N.toWeakModule := by

theorem dlmSimpleAdmissibleOrdinary_of_c2Cofinite (hC2 : C2Cofinite W) :
    DLMSimpleAdmissibleOrdinary W := by
\end{leancode}
\noindent File \leanpath{VOA/\allowbreak Classification/\allowbreak COne/\allowbreak LiteratureInputsProved.lean}, namespace
\lean{VOA.Classification.COne.LiteratureInputs} (items 2, 3, 4, 6):
\begin{leancode}
variable {𝒱 : Type u} [AddCommGroup 𝒱] [Module ℂ 𝒱] {W : VertexOperatorAlgebra 𝒱}

theorem abdIrreducibleWeakIsOrdinary_holds : ABDIrreducibleWeakIsOrdinary W :=
  abdIrreducibleWeakIsOrdinary_of_weakIsAdmissible abdWeakIsAdmissible

theorem abdRegularOfC2AndRational_holds : ABDRegularOfC2AndRational W :=
  abdRegularOfC2AndRational_of_components abdWeakIsAdmissible
    abdIrreducibleWeakIsOrdinary_holds

theorem zhuAdmissibleSemisimple_holds : ZhuAdmissibleSemisimple W :=
  VOA.Huang.ZhuSemisimple.zhuAdmissibleSemisimple_of_Pr
    (voaPBWBottomProjection (parabolicRepOfAMod W (regularAMod W)))

theorem dlmFinitelyManyIrredAdmissible_of_c2Cofinite (hC2 : C2Cofinite W) :
    DLMFinitelyManyIrredAdmissible W := fun hAR =>
  c2SemisimpleAdmissibleFiniteness_of_bottomProjection
    (voaPBWBottomProjection (parabolicRepOfAMod W (regularAMod W))) hC2
    (VOA.Huang.ZhuSemisimple.isSemisimpleRing_A_of_admissiblyRational
      (voaPBWBottomProjection (parabolicRepOfAMod W (regularAMod W))) hAR)
\end{leancode}

\begin{dictionary}
\textbf{Reading the statements.} Each conclusion is one of the named predicates of Section~\ref{l12:sec:litnamed},
applied to the implicit VOA \lean{W}; unfolding the predicate gives the mathematical statement:
\begin{itemize}
\item \lean{ABDWeakIsAdmissible W} $=$ \lean{C2Cofinite W → ∀ M (A : WeakModule W M), ∃ N : AdmissibleModule W M,
  N.toWeakModule = A}: item~1. The equation \lean{N.toWeakModule = A} says that the admissible structure \emph{extends} the
  given weak action (same carrier, same vertex operators); it is not merely an isomorphic module.
\item \lean{ABDIrreducibleWeakIsOrdinary W}: under \lean{C2Cofinite W}, a weak module \lean{A} with \lean{Nontrivial M}
  (``$M\neq0$'') and every weak submodule \lean{P} equal to $\bot$ or $\top$ (irreducibility in the weak sense) has an
  ordinary structure \lean{N : OrdVModule W M} with \lean{N.toWeakModule = A}: item~2.
\item \lean{ABDRegularOfC2AndRational W} $=$ \lean{C2Cofinite W → AdmissiblyRational W → Regular W}: item~3, with
  \lean{Regular} of Definition~\ref{l12:def:regular}.
\item \lean{ZhuAdmissibleSemisimple W}: item~4, as in Theorem~\ref{l12:thm:zhuss}; here the bottom projection \lean{Pr} is
  supplied by \lean{voaPBWBottomProjection} (Chapter~\ref{ch:L13}), built from the PBW theorem \lean{VOA.pbw} of
  Chapter~\ref{ch:L02}, which is a theorem, not an axiom.
\item \lean{irreducible_admissible_isOrdinary}: item~5 in its direct form (the carrier \lean{M} is an implicit section
  variable \lean{{M : Type u} [AddCommGroup M] [Module ℂ M]}), for one admissible module \lean{N} with
  \lean{N.IsIrreducible} (weak simplicity, Definition~\ref{l12:def:admissible}). \lean{dlmSimpleAdmissibleOrdinary_of_c2Cofinite}
  packages it as the predicate \lean{DLMSimpleAdmissibleOrdinary W}, whose \lean{AdmissiblyRational W} antecedent is simply
  not used.
\item \lean{DLMFinitelyManyIrredAdmissible W} $=$ \lean{AdmissiblyRational W → Finite (Φadm W)}: item~6. The proof term is
  displayed because it is informative: given the antecedent \lean{hAR}, it feeds the semisimplicity of \lean{A W}
  (item~4, with the same bottom projection) and \lean{hC2} into \lean{c2SemisimpleAdmissibleFiniteness_of_bottomProjection}
  (\leanpath{AdmissibleOrdinaryFinite.lean}).
\end{itemize}

\textbf{Hypotheses.} The four \lean{_holds} theorems and \lean{abdWeakIsAdmissible} have \emph{no} hypotheses; \lean{W} is
implicit and arbitrary. The two DLM theorems have the single hypothesis \lean{hC2 : C2Cofinite W} (Definition~\ref{l12:def:c2}),
outside the predicate.

\textbf{Same as the literature?} Items~1--3 are ABD's statements exactly, for VOAs of CFT type (which every library VOA is),
with the library's \lean{AdmissiblyRational} as ``rational''; ABD themselves use DLM-rationality, which implies
\lean{AdmissiblyRational} on paper (warning of Section~\ref{l12:warn:admcr}), so the Lean items~1--3 imply ABD's. Item~4 is
Theorem~\ref{l12:thm:zhuss}. Items~5 and~6 are \emph{weaker} than DLM's Theorem~8.1 in one respect and \emph{stronger} in
another: DLM need no $C_2$-cofiniteness (rationality alone suffices for them), while the library needs \lean{hC2}; on the other
hand item~5 drops rationality altogether. For the VOAs the book cares about (strongly rational ones) this makes no difference.
For an admissibly rational $V$ not known to be $C_2$-cofinite the two DLM predicates remain \textbf{unproved}; since
``rational $\Rightarrow C_2$-cofinite'' is an open conjecture, that case is genuinely open in the library. Fidelity point~(a)
of Section~\ref{l12:sec:litnamed} still applies to item~5: \lean{IsIrreducible} for admissible modules is weak simplicity.

\textbf{Could it be vacuous?} No. \lean{C2Cofinite} is satisfied by the lattice VOAs $V_{\sqrt{2m}\Z}$ and by the triplet
algebra $W(2)$ (Example~\ref{l12:ex:c2}); for $W(2)$, which is irrational, items~1, 2 and~5 are informative, and items~3
and~6 are vacuously true there because their rationality antecedent fails. For $V_{\sqrt{2m}\Z}$, which is strongly
rational (Chapter~\ref{ch:L22}), all six items apply. The conclusions are not trivially satisfiable either: an admissible
or ordinary structure must reproduce the given vertex operators exactly.
\end{dictionary}

\begin{proof}[Proof on paper]
Throughout, $V$ is $C_2$-cofinite, so Zhu's algebra $A(V)$ is finite-dimensional (Theorem~\ref{l12:thm:an}).

\emph{A finite spectrum.} Let $M$ be any weak module and $\alpha$ a generalised eigenvalue of $L_0$ on $M$; pick an honest
eigenvector $x$. The weak submodule generated by $x$ is spanned, by $C_2$-cofiniteness (finite strong generation,
Theorem~\ref{l12:thm:fingen}), by monomials in finitely many generators with bounded modes applied to $x$. Since a mode
$u_{(m)}$ of a vector $u\in V_k$ shifts $L_0$-eigenvalues by $k-m-1$ (the commutator formula of
Theorem~\ref{l12:thm:modcomm}), this submodule is an ordinary module whose weights lie in $\alpha+\Z$ and are bounded
below. Its lowest weight space is a module for the finite-dimensional algebra $A(V)$, on which $L_0$ acts by the image of
the central element $[\omega]$. A central element of a finite-dimensional algebra has a minimal polynomial, and on every
module its eigenvalues are roots of that polynomial. So the lowest weight lies in a finite set $S\subset\C$ that depends
only on $V$, and $\alpha\in S+\N$. In particular the real parts of all $L_0$-eigenvalues on all weak modules are bounded
below by one constant.

\emph{Item 1.} With this bound, every weak module is the sum of its generalised $L_0$-eigenspaces (Buhl's spanning set makes
$L_0$ locally finite; in the library this is the ``degree normal form'' of the proof file), and grading $M$ by ``eigenvalue minus the least element of $S+\N$ in the same class mod~$\Z$'' gives a
$\N$-grading with $u_{(m)}M(n)\subseteq M(n+k-m-1)$: the admissible structure, with the original vertex operators.

\emph{Item 5.} Let $N$ be irreducible admissible. Its lowest nonzero level $N(0)$ is a simple $A(V)$-module, hence
finite-dimensional, and the central $[\omega]$ acts on it by a scalar $h$. Irreducibility makes $N$ generated by $N(0)$, so
$N(j)$ is spanned by finitely many monomials applied to $N(0)$, is finite-dimensional, and $L_0$ acts on it by $h+j$. These
are exactly the axioms of an ordinary module, with weight spaces $N_{h+j}=N(j)$.

\emph{Item 2} is item~1 followed by item~5: the admissible structure of an irreducible weak module is irreducible (the weak
submodules are the same), hence ordinary. \emph{Item 3:} by item~1 every weak module is admissible, so rationality makes it
completely reducible; by item~2 its irreducible summands are ordinary; that is regularity. \emph{Item 4} is
Theorem~\ref{l12:thm:zhuss}.

\emph{Item 6.} By the finiteness theorem for $C_2$-cofinite VOAs (Chapter~\ref{ch:L21}, Theorem~\ref{l21:thm:finitephi}; in the
library this route uses the semisimplicity of $A(V)$ supplied by item~4 and the bottom projection), $\Phi(V)$ is finite. By
item~5 every irreducible admissible class has an ordinary representative, so the injection $\Phi(V)\hookrightarrow\Phi_{\mathrm{adm}}(V)$
of Theorem~\ref{l12:thm:lattice}(3) is also surjective, and $\Phi_{\mathrm{adm}}(V)$ is finite.
\end{proof}

\begin{physics}
These are the basic facts of \emph{lisse} ($C_2$-cofinite) chiral algebras. Item~1: every representation of a lisse chiral
algebra, even one given without any energy operator, has an energy spectrum bounded below, so it is a legitimate state
space of a chiral CFT. Items~2 and~5: every irreducible representation is a genuine sector, with finite degeneracy at each
energy and hence a well-defined character $\operatorname{tr}q^{L_0-c/24}$. Item~3: a lisse theory in which every graded
representation splits splits completely, i.e.\ it is a rational CFT in the strongest sense. Item~6: such a theory has
finitely many primary fields. A physicist takes all of this for granted in a rational CFT; the theorem says the
library no longer has to.
\end{physics}

\begin{status}
\texttt{\#print axioms}, run for this revision at an earlier commit (not re-run for this edition) with \texttt{lake env lean -DmaxSynthPendingDepth=3} from the
repository root: \lean{abdWeakIsAdmissible}, \lean{dlmSimpleAdmissibleOrdinary_of_c2Cofinite},
\lean{abdIrreducibleWeakIsOrdinary_holds}, \lean{abdRegularOfC2AndRational_holds}, \lean{zhuAdmissibleSemisimple_holds},
\lean{dlmFinitelyManyIrredAdmissible_of_c2Cofinite}: \texttt{[propext, Classical.choice, Quot.sound]} each. The olean of
\leanpath{LiteratureInputsProved.lean} was not present in the main checkout, so the scratch file was a verbatim copy of that
file followed by the \texttt{\#print axioms} commands, compiled against the existing oleans of its imports (exit status~0). The
coordinator reports the umbrella build of \lean{VOA.Classification}, which imports the file, as green (12,179 jobs, standard
axioms); that build was not re-run for this chapter. \textbf{Proved, unconditional} for items~1--4 (no hypotheses);
items~5--6 are proved under the stated hypothesis \lean{hC2}. Items~4 and~6 use the theorem \lean{VOA.pbw} through
\lean{voaPBWBottomProjection}.
\end{status}

\subsection{The proved implications}

\begin{theorem}[The rationality lattice]\label{l12:thm:lattice}
For every VOA $V$:
\begin{enumerate}
\item regular $\Rightarrow$ admissibly rational $\Rightarrow$ every ordinary module is completely reducible;
\item regular $\Rightarrow$ Gannon clause (3);
\item $\Phi(V)\hookrightarrow\Phi_{\mathrm{adm}}(V)$, so $\Phi_{\mathrm{adm}}(V)$ finite $\Rightarrow\Phi(V)$ finite;
\item admissibly rational $+$ [DLM 8.1(b)] $+$ [DLM 8.1(c)] $\Rightarrow$ Gannon-rational;
\item regular $+$ [DLM 8.1(b)] $\Rightarrow$ Gannon-rational;
\item [ABD Prop.~5.7] $+$ [ABD Cor.~5.8] $\Rightarrow$ [ABD Thm.~4.5, hard direction];
\item [ABD Thm.~4.5] $+$ [Li Thm.~3.8] $\Rightarrow$ (regular $\iff$ $C_2$-cofinite and admissibly rational).
\end{enumerate}
\end{theorem}
\noindent File \leanpath{VOA/\allowbreak Modules/\allowbreak RationalityLattice.lean}.
\begin{leancode}
theorem admissiblyRational (h : Regular W) : AdmissiblyRational W :=
  fun N P => h.weakCompletelyReducible N.toWeakModule P

theorem OrdVModule.isCompletelyReducible_of_admissiblyRational {W : VertexOperatorAlgebra 𝒱}
    {M : Type u} [AddCommGroup M] [Module ℂ M] (N : OrdVModule W M)
    (hAR : AdmissiblyRational W) : OrdVModule.IsCompletelyReducible N :=
  N.isCompletelyReducible_of_hasAdmissibleRefinement hAR N.hasAdmissibleRefinement

theorem gannonAdmissibleIsOrdinary (h : Regular W) : GannonAdmissibleIsOrdinary W :=
  fun N hN => h.irreducibleWeakIsOrdinary N.toWeakModule hN

theorem finite_Φ_of_finite_Φadm (W : VertexOperatorAlgebra 𝒱) (h : Finite (Φadm W)) :
    Finite (Φ W) :=

theorem gannonRational_of_admissiblyRational (hAR : AdmissiblyRational W)
    (hFin : DLMFinitelyManyIrredAdmissible W) (hOrd : DLMSimpleAdmissibleOrdinary W) :
    GannonRational W where

theorem gannonRational_of_regular (h : Regular W)
    (hFin : DLMFinitelyManyIrredAdmissible W) : GannonRational W where

theorem abdRegularOfC2AndRational_of_components (hWA : ABDWeakIsAdmissible W)
    (hIrr : ABDIrreducibleWeakIsOrdinary W) : ABDRegularOfC2AndRational W := by

theorem regular_iff_c2Cofinite_and_admissiblyRational (hABD : ABDRegularOfC2AndRational W)
    (hLi : LiC2CofiniteOfRegular W) :
    Regular W ↔ (C2Cofinite W ∧ AdmissiblyRational W) :=
\end{leancode}

\begin{dictionary}
The first and third displayed one-liners show how thin some implications are: regular $\Rightarrow$ admissibly rational is
literally ``an admissible module is a weak module'' (both use the same weak-submodule notion), and
regular $\Rightarrow$ Gannon~(3) is ``an irreducible admissible module is an irreducible weak module''. The second uses
Theorem~\ref{l12:thm:ordadm}. Hypotheses in brackets in the mathematical statement are the named predicates of
Section~\ref{l12:sec:litnamed}; every other hypothesis is shown. All of them except [Li Thm.~3.8] are now theorems of the
library, the two DLM ones for $C_2$-cofinite $V$ (Theorem~\ref{l12:thm:litproved} and the status box below). \textbf{Not available} (no proof, no citation): admissibly rational
$\Rightarrow$ every irreducible \emph{weak} module is ordinary (needs $C_2$), Gannon-rational $\Rightarrow$ admissibly
rational, and rational $\Rightarrow C_2$-cofinite (a famous open conjecture). \textbf{False}: $C_2$-cofinite
$\Rightarrow$ rational (triplet $W(p)$, Adamovi\'c--Milas; in the library, $W(2)$, Example~\ref{l12:ex:c2}).
\end{dictionary}

\textbf{Proof on paper.} (1)--(3) as in the dictionary; (3): an irreducible ordinary module is an irreducible admissible
module via Theorem~\ref{l12:thm:ordadm}, and two ordinary modules are isomorphic iff their weak modules are, so the map
on classes is injective. (4): clause (1) by (1), clause (2) by (3) and DLM 8.1(b), clause (3) is DLM 8.1(c). (5): clause (3)
by (2). (6): by ABD 5.7 every weak module carries an admissible structure, so rationality makes it completely reducible,
and ABD 5.8 makes its irreducibles ordinary. (7): combine (1), (6) and Li. \qed

\begin{status}
All proved; \texttt{\#print axioms} (run for this chapter) of \lean{VOA.Regular.admissiblyRational},
\lean{VOA.OrdVModule.isCompletelyReducible\_of\_admissiblyRational}, \lean{VOA.finite\_Φ\_of\_finite\_Φadm},
\lean{VOA.gannonRational\_of\_admissiblyRational}, \lean{VOA.gannonRational\_of\_regular},
\lean{VOA.abdRegularOfC2AndRational\_of\_components}, \lean{VOA.regular\_iff\_c2Cofinite\_and\_admissiblyRational}:
standard three. As \emph{stated}, items (4)--(7) take the named literature predicates as hypotheses. After this revision
those hypotheses are dischargeable as follows (Theorem~\ref{l12:thm:litproved}):
\begin{itemize}
\item item (6) has no remaining content as a conditional: its conclusion is \lean{abdRegularOfC2AndRational_holds};
\item items (4) and (5) become unconditional \emph{for $C_2$-cofinite $V$}: [DLM 8.1(b)] is
  \lean{dlmFinitelyManyIrredAdmissible_of_c2Cofinite hC2} and [DLM 8.1(c)] is \lean{dlmSimpleAdmissibleOrdinary_of_c2Cofinite hC2};
  for strongly rational $V$ the composite is the named theorem \lean{StronglyRational.gannonRational'}
  (Theorem~\ref{l12:thm:srlit}). For admissibly rational $V$ not known to be $C_2$-cofinite they stay conditional;
\item item (7) remains \textbf{conditional on \lean{LiC2CofiniteOfRegular}} (H.~Li, Thm.~3.8), which is \textbf{unproved}; its
  other hypothesis \lean{hABD} is discharged by \lean{abdRegularOfC2AndRational_holds}. The direction ``$C_2$-cofinite and
  admissibly rational $\Rightarrow$ regular'' is therefore a theorem; the direction ``regular $\Rightarrow C_2$-cofinite'' is not.
\end{itemize}
No single library declaration states item (7) with \lean{hABD} discharged; the substitution above is a one-line composition
we describe, not a quoted theorem.
\end{status}

\begin{physics}
The lattice is the map of ``how nice'' a chiral CFT is. Lisse + rational (= regular) theories --- lattice CFTs, WZW models
at positive integer level, minimal models, $\Vnat$ --- are the modular tensor category world, where Verlinde's formula
holds. Lisse but irrational theories (the triplet algebras) are logarithmic CFTs. Rational but not known to be lisse: no
example is known, which is the content of the open conjecture.
\end{physics}

\subsection{Self-duality, strongly finite and strongly rational}\label{l12:sec:strong}

\begin{definition}[Self-dual VOA; McRae §2.2, Dong--Mason §2]\label{l12:def:selfdual}
$V$ is \emph{self-dual} if it carries a nondegenerate fully invariant bilinear form, i.e.\ a form with
$(Y(u,z)a|b) = (a|Y(e^{zL_1}(-z^{-2})^{L_0}u,z^{-1})b)$ and zero radical.
\end{definition}
\noindent File \leanpath{VOA/\allowbreak Modules/\allowbreak StronglyRational.lean}.
\begin{leancode}
def SelfDual (W : VertexOperatorAlgebra 𝒱) : Prop :=
  ∃ B : InvariantForm W, B.IsFullyInvariant ∧ InvariantForm.leftRadical B = ⊥

theorem selfDual_iff (hsimple : W.IsSimple) :
    W.SelfDual ↔ ∀ v, v ∈ W.grading 1 → W.L 1 v = 0 :=
\end{leancode}

\begin{dictionary}
\lean{InvariantForm W}, \lean{IsFullyInvariant}, \lean{leftRadical} are defined in Chapter~\ref{ch:L04}: a bilinear form with
$(L_na|b) = (a|L_{-n}b)$, the full FHL invariance (mode form of the displayed identity), and $\{a : (a|\cdot) = 0\}$. For a
module, ``$V\cong V'$'' is equivalent to the existence of such a form (FHL Prop.~5.3.6); the library uses the form version,
which avoids the vacuous \lean{SelfContragredient}. \lean{selfDual_iff}: for simple $V$, self-dual iff $L_1V_1 = 0$
(``strong CFT type''); this is Dong--Mason §2 / H.~Li's theorem, here proved in both directions.
\end{dictionary}

\textbf{Proof on paper.} If $L_1v = \lambda\ltwelvevac$ with $\lambda\neq0$ for some $v\in V_1$, then
$\lambda(\ltwelvevac|\ltwelvevac) = (L_1v|\ltwelvevac) = (v|L_{-1}\ltwelvevac) = 0$, so $(\ltwelvevac|\ltwelvevac) = 0$; with $V_m\perp V_n$ for
$m\neq n$ and $V_0 = \C\ltwelvevac$ this puts $\ltwelvevac$ in the radical. Conversely, if $L_1V_1 = 0$, Li's canonical form (the coefficient of
$\ltwelvevac$ in $Y(e^{zL_1}(-z^{-2})^{L_0}a,z^{-1})b$, suitably normalised) gives a fully invariant form whose radical is an ideal not containing $\ltwelvevac$,
hence $0$ by simplicity (Chapter~\ref{ch:L04}). \qed

\begin{status}
Proved; \texttt{\#print axioms VOA.VertexOperatorAlgebra.selfDual\_iff}: standard three.
\end{status}

\begin{physics}
A nondegenerate invariant form is the two-point function $\langle a(z)b(w)\rangle$ of the CFT on the sphere; self-duality
is the statement that the vacuum sector is its own charge conjugate, i.e.\ that every field has a nonzero two-point
function with some field. $L_1V_1 = 0$ says the currents are quasi-primary, which is what makes their two-point functions
conformally covariant.
\end{physics}

\begin{definition}[Strongly finite, strongly rational; Creutzig--Gannon, Dong--Mason, McRae]\label{l12:def:strong}
$V$ is \emph{strongly finite} if it is simple, self-dual and $C_2$-cofinite (and of CFT type, which is automatic here);
\emph{strongly rational} if moreover it is admissibly rational.
\end{definition}
\noindent File \leanpath{VOA/\allowbreak Modules/\allowbreak StronglyRational.lean}.
\begin{leancode}
structure StronglyFinite (W : VertexOperatorAlgebra 𝒱) : Prop where
  simple : W.IsSimple
  selfDual : W.SelfDual
  c2Cofinite : C2Cofinite W

structure StronglyRational (W : VertexOperatorAlgebra 𝒱) : Prop extends StronglyFinite W where
  rational : AdmissiblyRational W

theorem stronglyFinite_iff {W : VertexOperatorAlgebra 𝒱} (hsimple : W.IsSimple)
    (hc2 : C2Cofinite W) :
    StronglyFinite W ↔ ∀ v, v ∈ W.grading 1 → W.L 1 v = 0 :=

theorem stronglyRational_of_regular {W : VertexOperatorAlgebra 𝒱} (h : StronglyFinite W)
    (hreg : Regular W) : StronglyRational W where

theorem stronglyRational_of_span_vacuum {W : VertexOperatorAlgebra 𝒱} (hne : W.vacuum ≠ 0)
    (hspan : Submodule.span ℂ ({W.vacuum} : Set 𝒱) = ⊤) : StronglyRational W where
\end{leancode}

\begin{dictionary}
\textbf{Literature comparison.} Creutzig--Gannon (2017, §2.2): strongly finite = simple, $C_2$-cofinite, CFT type,
$V\cong V'$; strongly rational = strongly finite + \emph{regular}. Dong--Mason (2004, §3): strongly rational = strong CFT
type ($L_1V_1 = 0$), $C_2$-cofinite, rational. McRae (§2.2): CFT type, $V\cong V'$, $C_2$-cofinite, rational. The library:
\lean{StronglyFinite} = Creutzig--Gannon's (with $V\cong V'$ in the form version, and CFT type built in);
\lean{StronglyRational} = McRae's/Dong--Mason's definition with ``rational'' read as the library's
\lean{AdmissiblyRational} (DLM-rational $\Rightarrow$ \lean{AdmissiblyRational} is proved for $C_2$-cofinite $V$,
Theorem~\ref{l12:thm:dlmrat}, so every McRae/Dong--Mason strongly rational VOA satisfies the library's rationality; the
converse is open), plus an
explicit simplicity conjunct. By
\lean{selfDual_iff} and \lean{stronglyFinite_iff}, given simplicity, self-duality is exactly $L_1V_1 = 0$, so the
Dong--Mason and McRae versions agree; simplicity is implied in McRae's setting by CFT type + $V\cong V'$ (an argument the
library does not formalise, hence the explicit conjunct). Creutzig--Gannon's strongly rational (with regular) implies the
library's (\lean{stronglyRational_}\allowbreak\lean{of_}\allowbreak\lean{regular}); the converse needs ABD Thm.~4.5 (\lean{StronglyRational.}\allowbreak\lean{regular_}\allowbreak\lean{of_}\allowbreak\lean{abd}),
which is now a theorem, so the converse holds without hypothesis (\lean{StronglyRational.regular}, Theorem~\ref{l12:thm:srlit}).
So: implied by McRae/Dong--Mason (given simplicity), the same \emph{modulo the open converse
\lean{AdmissiblyRational} $\Rightarrow$ DLM-rational}, and
\emph{equivalent} to Creutzig--Gannon's (given
\lean{StronglyFinite}, admissible rationality and regularity are interchangeable in Lean).

\textbf{Non-vacuity.} \lean{stronglyRational_}\allowbreak\lean{of_}\allowbreak\lean{span_}\allowbreak\lean{vacuum}: any VOA spanned by a nonzero vacuum is strongly rational
(instantiated at the trivial VOA in Chapter~\ref{ch:L15}). More interesting, the lattice triplet algebra $W(2)$ is
\emph{strongly finite and irrational}:
\lean{VOA.}\allowbreak\lean{LatGen.}\allowbreak\lean{Wp.}\allowbreak\lean{tripletWVOA_}\allowbreak\lean{two_}\allowbreak\lean{stronglyFinite_}\allowbreak\lean{and_}\allowbreak\lean{irrational} (Chapter~\ref{ch:L18}; \texttt{\#print axioms} run for
this chapter: standard three). Nontrivial strongly rational VOAs now exist in the library: the rank-one lattice VOAs
$V_{\sqrt{2m}\Z}$ (\lean{V2m_stronglyRational}, Chapter~\ref{ch:L22}). Strong rationality gives Zhu's modularity theorem for
them, through Theorem~\ref{l12:thm:zhuss}: this is \lean{zhuModularity_V2m}, and more generally
\lean{zhuModularity_}\allowbreak\lean{canonical_of_}\allowbreak\lean{stronglyRational} (Chapter~\ref{ch:L15}). (The first release
revision of this chapter said the opposite; that was wrong.) $\Vnat$ is not proved strongly rational in the library.
\end{dictionary}

\textbf{Proof on paper.} \lean{stronglyFinite_iff} is \lean{selfDual_iff}. For $V = \C\ltwelvevac$: every subspace of every weak
module is a weak submodule (the only modes are multiples of $\ltwelvevac_{(-1)} = \mathrm{id}$), so complements exist; $V$ is
simple, $L_1 = 0$, and $\dim V/C_2(V)\le1$. \qed

\begin{status}
Proved; \texttt{\#print axioms} of \lean{VOA.stronglyFinite\_iff}, \lean{VOA.stronglyRational\_of\_regular},
\lean{VOA.stronglyRational\_of\_span\_vacuum}: standard three.
\end{status}

\begin{physics}
Strongly rational VOAs are the chiral algebras of the ``nicest'' CFTs: finitely many sectors, all of them ordinary,
semisimple fusion, a nondegenerate sphere two-point function and no negative-energy states. Huang's theorem (whose
formalisation in the library, and its hypotheses, are the subject of Chapter~\ref{ch:L21}) shows that their module categories are modular tensor categories, which is the mathematical form of the
Moore--Seiberg axioms of rational CFT. Strongly finite but irrational theories are the logarithmic CFTs whose
representation categories are only log-modular.
\end{physics}

\begin{theorem}[Strongly rational VOAs: regular, finitely many admissible irreducibles, Gannon-rational]\label{l12:thm:srlit}
Let $V$ be strongly rational (Definition~\ref{l12:def:strong}). Then
\begin{enumerate}
\item $V$ is regular (Definition~\ref{l12:def:regular});
\item $V$ satisfies the predicate \lean{DLMFinitelyManyIrredAdmissible}, and $V$ has only finitely many isomorphism classes
  of irreducible admissible modules;
\item $V$ is Gannon-rational (Definition~\ref{l12:def:gannon}).
\end{enumerate}
No named literature hypothesis is assumed.
\end{theorem}
\noindent File \leanpath{VOA/\allowbreak Modules/\allowbreak StronglyRational.lean} (the older forms, which carry the
literature predicates as hypotheses):
\begin{leancode}
theorem StronglyRational.regular_of_abd {W : VertexOperatorAlgebra 𝒱} (h : StronglyRational W)
    (hABD : ABDRegularOfC2AndRational W) : Regular W :=
  hABD h.c2Cofinite h.rational

theorem gannonRational (h : StronglyRational W) (hFin : DLMFinitelyManyIrredAdmissible W)
    (hOrd : DLMSimpleAdmissibleOrdinary W) : GannonRational W :=
  gannonRational_of_admissiblyRational h.rational hFin hOrd
\end{leancode}
\noindent File \leanpath{VOA/\allowbreak Classification/\allowbreak COne/\allowbreak LiteratureInputsProved.lean} (the
discharged forms; same section variables as in Theorem~\ref{l12:thm:litproved}):
\begin{leancode}
theorem _root_.VOA.StronglyRational.regular (h : StronglyRational W) : Regular W :=
  h.regular_of_abd abdRegularOfC2AndRational_holds

theorem _root_.VOA.StronglyRational.dlmFinitelyManyIrredAdmissible (h : StronglyRational W) :
    DLMFinitelyManyIrredAdmissible W :=
  dlmFinitelyManyIrredAdmissible_of_c2Cofinite h.c2Cofinite

theorem _root_.VOA.StronglyRational.finite_Φadm (h : StronglyRational W) : Finite (Φadm W) :=
  h.dlmFinitelyManyIrredAdmissible h.rational

theorem _root_.VOA.StronglyRational.gannonRational' (h : StronglyRational W) :
    GannonRational W :=
  h.gannonRational h.dlmFinitelyManyIrredAdmissible
    (dlmSimpleAdmissibleOrdinary_of_stronglyRational h)
\end{leancode}

\begin{dictionary}
\textbf{Names.} \lean{_root_.VOA.StronglyRational.regular} declares the theorem with the full name
\lean{VOA.StronglyRational.regular}, although the file's namespace is \lean{VOA.Classification.COne.LiteratureInputs}; this
places it in the namespace of the structure \lean{StronglyRational}, so that for \lean{h : StronglyRational W} one may write
\lean{h.regular} (dot notation). The same holds for the other three. The older \lean{gannonRational} is likewise
\lean{VOA.StronglyRational.gannonRational} (it sits inside \lean{namespace StronglyRational} in its file).

\textbf{Hypotheses.} Only \lean{h : StronglyRational W}: simple, self-dual, $C_2$-cofinite (\lean{h.c2Cofinite}) and
admissibly rational (\lean{h.rational}). The older forms carry, in addition, \lean{hABD}, \lean{hFin}, \lean{hOrd}: exactly the
literature predicates [ABD Thm.~4.5], [DLM 8.1(b)], [DLM 8.1(c)]. The new forms discharge them by
\lean{abdRegularOfC2AndRational_holds}, \lean{dlmFinitelyManyIrredAdmissible_of_c2Cofinite h.c2Cofinite} and
\lean{dlmSimpleAdmissibleOrdinary_of_stronglyRational h} (the $C_2$-locus theorem of Theorem~\ref{l12:thm:litproved}
item~5, applied to \lean{h.c2Cofinite}); the one-line proof terms are shown because they are exactly these compositions.

\textbf{Conclusions.} \lean{Regular W}, \lean{Finite (Φadm W)} (with \lean{Φadm} the set of isomorphism classes of
irreducible admissible modules) and \lean{GannonRational W} of Definitions~\ref{l12:def:regular} and~\ref{l12:def:gannon}.
\lean{StronglyRational.dlmFinitelyManyIrredAdmissible} is the predicate form (\lean{AdmissiblyRational W → Finite (Φadm W)});
\lean{finite_Φadm} applies it to \lean{h.rational}.

\textbf{Same as the literature?} Yes, with ``rational'' read as \lean{AdmissiblyRational}: for strongly rational VOAs in
Dong--Mason's sense these are ABD Thm.~4.5 and DLM Thm.~8.1, and with \lean{StronglyRational.regular} the library's strong rationality coincides with Creutzig--Gannon's
(strongly finite + regular). \textbf{Vacuous?} No: \lean{V2m_stronglyRational} (Chapter~\ref{ch:L22}) supplies nontrivial
instances.
\end{dictionary}

\textbf{Proof on paper.} A strongly rational VOA is $C_2$-cofinite and admissibly rational; apply items~3, 5 and~6 of
Theorem~\ref{l12:thm:litproved}, and for (3) Theorem~\ref{l12:thm:lattice}(4). \qed

\begin{physics}
For the chiral algebra of a rational, lisse CFT with a nondegenerate two-point function, every structural property a physicist assumes ---
every representation splits into sectors, there are finitely many primaries, and all sectors have finite degeneracies --- is
now a theorem in Lean, from the four defining properties alone. This is the input that Zhu's modular-invariance theorem
(Chapter~\ref{ch:L15}) and the modular tensor category of Chapter~\ref{ch:L21} rely on.
\end{physics}

\begin{status}
Proved, \textbf{unconditional} given \lean{h : StronglyRational W}. \texttt{\#print axioms} of
\lean{VOA.StronglyRational.regular}, \lean{.dlmFinitelyManyIrredAdmissible}, \lean{.finite_Φadm}, \lean{.gannonRational'}:
\texttt{[propext, Classical.choice, Quot.sound]} (same scratch run as Theorem~\ref{l12:thm:litproved}). The last three use the theorem \lean{VOA.pbw} through \lean{voaPBWBottomProjection}.
\end{status}

\subsection{Summary table}

\begin{center}
\small
\begin{tabular}{p{3.1cm}p{4.6cm}p{5.2cm}}
\toprule
Lean predicate & quantifies over / content & literature; same / weaker / stronger \\
\midrule
\lean{AdmissiblyRational} & admissible modules, weak complements & DLM ``rational''; DLM $\Rightarrow$ library proved under $C_2$ (\lean{admissiblyRational_of_dlmRational}); converse
open; library $\Rightarrow$ graded reading proved \\
\lean{GannonRational} & ordinary c.r.\ + $|\Phi|<\infty$ + adm.\ irreducibles ordinary & Gannon Def.~5.3.2 (his ``weakly
rational''); same \\
\lean{Holomorphic} & Gannon-rational, $|\Phi| = 1$ & Gannon; same \\
\lean{Regular} & all weak modules c.r., irreducibles ordinary & DLM 1997 Def.~3.1; same \\
\lean{C2Cofinite} & $\dim V/C_2(V)<\infty$ & Zhu; same \\
\lean{EtingofFinite} & $\dim\ltwelveHP(R_V)<\infty$ & weaker than AK quasi-lisse \\
\lean{StronglyFinite} & simple, self-dual, $C_2$ & Creutzig--Gannon; same \\
\lean{StronglyRational} & strongly finite + adm.\ rational & McRae / Dong--Mason with \lean{AdmissiblyRational} (see
first row). Creutzig--Gannon's implies it \\
\lean{SelfContragredient} & a linear iso $M\cong M'$ & \textbf{vacuous}; not ``$M\cong M'$ as modules'' \\
\lean{Unitary} & bilinear, non-isotropic, mode-symmetric form & \textbf{unsatisfiable} in $\dim\ge2$; not unitarity \\
\bottomrule
\end{tabular}
\end{center}

%======================================================================
\section{Recorded traps and dead ends}\label{l12:sec:traps}
%======================================================================

For the convenience of the referee we collect the traps of this chapter in one place, with their severity on a four-step scale:
serious (a statement that is wrong or vacuous), moderate (weaker or stronger than advertised, or a redundant
hypothesis), open gap (a conditional result whose hypothesis is not discharged, or a missing piece), and documentation
only (docstrings and notes).
\begin{enumerate}
\item \lean{fusionCoeff} $\equiv0$ (serious; Theorem~\ref{l12:thm:fusionzero}, proved in the library). No consumer.
\item \lean{Unitary} unsatisfiable in dimension $\ge2$ and not the literature's unitarity (serious;
  Proposition~\ref{l12:prop:unitary}, proved in the library). No consumer.
\item \lean{SelfContragredient} holds for every module (serious, proved in the library by
  \lean{selfContragredient_}\allowbreak\lean{of_}\allowbreak\lean{ordVModule}). No consumer.
\item \lean{InducedAssignment} uninhabited; \lean{forwardMap_invFun}, \lean{forwardMap_}\allowbreak\lean{comp_}\allowbreak\lean{invFun},
  \lean{forwardMap_}\allowbreak\lean{invFun_}\allowbreak\lean{ofPointwise} vacuous (serious, proved in the library by \lean{not_}\allowbreak\lean{nonempty_}\allowbreak\lean{inducedAssignment});
  stale docstrings in \leanpath{ZhuInduced.lean}, \leanpath{ZhuBijection.lean} (documentation only).
\item \lean{IsLowestWeight} allows $M_h = 0$ and requires a single coset (moderate; Section~\ref{l12:sec:char}).
\item \lean{On W n} omits DLM's generators $(L_{-1}+L_0)a$ for $n\ge1$; \lean{An} attributed to Frenkel--Zhu (moderate; documentation only;
  Section~\ref{l12:sec:an}).
\item \lean{A₂ W} is $V/C_2(V)$, not the second higher Zhu algebra (documentation only; Section~\ref{l12:sec:c2}).
\item \lean{OrdVModule.boundedBelow} and \lean{LogVModule.boundedBelow} are global bounds, stronger than FHL/HLZ per-coset
  truncation (moderate; Section~\ref{l12:sec:categories}).
\item \lean{conformalWeight_}\allowbreak\lean{char_}\allowbreak\lean{periodicity} carries two unused hypotheses (moderate, marked by the library with
  underscores).
\item \lean{isSimpleAMod_}\allowbreak\lean{of_}\allowbreak\lean{irreducible_}\allowbreak\lean{of_}\allowbreak\lean{semisimple} assumes semisimplicity of $A(V)$ and the PBW datum \lean{Pr}, where Zhu's
  theorem needs neither (moderate; Theorem~\ref{l12:thm:simplebottom}). The gap is closed elsewhere: the hypothesis-free
  \lean{isSimpleAMod_toAMod_of_irreducible} (Chapter~\ref{ch:L13}).
\item The docstrings of \leanpath{AdmissibleRationality.lean} contradict each other on whether the weak-submodule reading of
  complete reducibility implies the graded one (documentation only; Section~\ref{l12:warn:admcr}); DLM $\Rightarrow$
\lean{AdmissiblyRational} is now formalised under $C_2$ (Theorem~\ref{l12:thm:dlmrat}), the converse is not.
\item \leanpath{Fusion.lean} docstring: ``built here faithfully'' and ``the Verlinde apex \ldots wires to \lean{fusionCoeff}'' (documentation only).
  \leanpath{Contragredient.lean}: the closing section comment still calls the contragredient action a frontier (documentation only).
\item The module docstrings of \leanpath{RationalityLattice.lean} and \leanpath{Rationality.lean}, and the docstrings of
  \lean{StronglyRational} and \lean{StronglyRational.regular_of_abd}, still call the ABD and DLM predicates ``cited'' or
  ``carried as a citation''; six of the seven are theorems (Theorem~\ref{l12:thm:litproved}). Documentation only.
\item Open gap: \lean{LiC2CofiniteOfRegular} is unproved, so \lean{regular_iff_c2Cofinite_and_admissiblyRational} is conditional
  on it (Section~\ref{l12:sec:litnamed}).
\end{enumerate}
Retired dead ends recorded by the library itself: the predicate \lean{WeaklyRational} (wrong module category in its third
clause, detached conclusion; replaced by \lean{GannonRational} with no alias), and Gannon's conjecture ``weakly rational
$\iff$ $C_2$-cofinite'' (false in the direction $C_2\Rightarrow$ rational, by the triplet algebras).

%======================================================================
\section{Referee's checklist}\label{l12:sec:checklist}
%======================================================================

\begingroup\small
\begin{longtable}{p{4.2cm}p{4.3cm}p{3.3cm}p{1.6cm}}
\toprule
Lean declaration & intended meaning & hypotheses; status & check \\
\midrule
\endhead
\lean{AdmissibleModule}, \lean{LogVModule} & admissible / logarithmic modules & definitions & Defs.~\ref{l12:def:admissible}, \ref{l12:def:log}\\
\lean{WeakModule}, \lean{WeakModule.act} & module $Y_M$ with vacuum, truncation and Borcherds identity & definition; \lean{∑ᶠ} finite & Def.~\ref{l12:def:weakmodule}\\
\lean{mode_commutator} (module) & $[u^M_{(m)},v^M_{(n)}] = \sum_i\binom mi(u_{(i)}v)^M_{(m+n-i)}$ & weak module; unconditional & Thm.~\ref{l12:thm:modcomm}\\
\lean{OrdVModule}, \lean{OrdVModule.}\allowbreak\lean{IsIrreducible}, \lean{IsCompletelyReducible} & ordinary module; irreducible; completely reducible (weak submodules) & definitions; \textbf{global} lower bound on weights & Def.~\ref{l12:def:ordmodule}\\
\lean{adjointWeakModule}, \lean{adjointModule} & $V$ as a module over itself & unconditional & Def.~\ref{l12:def:adjoint}\\
\lean{act_}\allowbreak\lean{T_}\allowbreak\lean{covariance} & $Y_M(L_{-1}a,z) = \partial_zY_M(a,z)$ & none; unconditional & Thm.~\ref{l12:thm:covariance}\\
\lean{LogVModule.}\allowbreak\lean{weight_}\allowbreak\lean{eq_}\allowbreak\lean{maxGenEigenspace}, \lean{isL0Semisimple_}\allowbreak\lean{iff_}\allowbreak\lean{L0nil_}\allowbreak\lean{eq_}\allowbreak\lean{zero} & grading canonical; ordinary $\iff L_0^{\mathrm{nil}} = 0$ & none; unconditional & Thm.~\ref{l12:thm:loggrading}\\
\lean{OrdVModule.}\allowbreak\lean{hasAdmissibleRefinement} & ordinary $\subseteq$ admissible & none; unconditional & Thm.~\ref{l12:thm:ordadm}\\
\lean{ShortExact.}\allowbreak\lean{gradedDim_}\allowbreak\lean{eq_}\allowbreak\lean{add} & characters additive & exactness; unconditional & Prop.~\ref{l12:prop:ses}\\
\lean{restrictWeakModule}, \lean{restrictOrdVModule} & restriction to a conformal sub-VOA & definitions & Def.~\ref{l12:def:restrict}\\
\lean{IsLowestWeight} & weights in $h+\N$ & \textbf{$M_h$ may be $0$} & Def.~\ref{l12:def:lowest}\\
\lean{char_}\allowbreak\lean{T_}\allowbreak\lean{transform} & $T$-periodicity of $\chi_M$ & conformal weight $h$ & Thm.~\ref{l12:thm:Tperiod}\\
\lean{zhuRingData_}\allowbreak\lean{holds}, \lean{instRingA} & $A(V)$ associative unital algebra & none; unconditional & Thm.~\ref{l12:thm:zhuring}\\
\lean{oMod_zhuMul_lowest_all}, \lean{rhoLW}, \lean{O_le_ker_rhoLW}, \lean{Pres_mem_O} & $o(u*v) = o(u)o(v)$ on $M_h$; $O(V)$ acts by $0$; residues lie in $O(V)$ & \lean{IsLowestWeight} (\textbf{$M_h$ may be $0$}); unconditional & Thm.~\ref{l12:thm:zhuaction}\\
\lean{moduleMh} & $M_h$ is an $A(V)$-module & \lean{IsLowestWeight}; unconditional & Thm.~\ref{l12:thm:moduleMh}\\
\lean{forwardMap} & $[M]\mapsto[M_h]$ & definition & Def.~\ref{l12:def:lwmodule}\\
\lean{InducedAssignment.}\allowbreak\lean{forwardMap_}\allowbreak\lean{invFun} & round trip of Zhu's bijection & \textbf{vacuous} (hypothesis uninhabited) & §\ref{l12:sec:zhu}\\
\lean{isSimpleAMod_}\allowbreak\lean{of_}\allowbreak\lean{irreducible_}\allowbreak\lean{of_}\allowbreak\lean{semisimple} & bottoms of irreducibles simple & \lean{Pr}, $A(V)$ semisimple; conditional (superseded by the unconditional Chapter~\ref{ch:L13} theorem) & Thm.~\ref{l12:thm:simplebottom}\\
\lean{DLMRational}, \lean{admissiblyRational_}\allowbreak\lean{of_}\allowbreak\lean{dlmRational} & DLM rationality; DLM $\Rightarrow$ library rationality & \lean{hC2}; proved; converse open & Thm.~\ref{l12:thm:dlmrat}\\
\lean{ZhuAdmissibleSemisimple} & rational $\Rightarrow A(V)$ semisimple & named predicate; \textbf{proved} for every $V$ (next row) & §\ref{l12:sec:rationality}\\
\lean{isSemisimpleRing_}\allowbreak\lean{A_of_}\allowbreak\lean{admissiblyRational}, \lean{zhuAdmissibleSemisimple_}\allowbreak\lean{of_Pr} & Zhu's semisimplicity theorem & \lean{Pr} (from \lean{VOA.pbw}), \lean{AdmissiblyRational}; proved; used by the Zhu headline & Thm.~\ref{l12:thm:zhuss}\\
\lean{isSimple_}\allowbreak\lean{of_}\allowbreak\lean{zhuPolynomial} & Zhu-polynomial simplicity test & $V\neq0$, $V_1 = 0$, relation, roots & Thm.~\ref{l12:thm:zhupoly}\\
\lean{C2Cofinite}, \lean{A₂CommRing}, \lean{A₂LieRing} & $C_2$; Poisson algebra $R_V$ & definitions; unconditional & Def.~\ref{l12:def:c2}, Thm.~\ref{l12:thm:poisson}\\
\lean{iterProdSpan_}\allowbreak\lean{eq_}\allowbreak\lean{top} & finite strong generation & $V = \mathrm{span}S + C_2$ & Thm.~\ref{l12:thm:fingen}\\
\lean{EtingofFinite_}\allowbreak\lean{of_}\allowbreak\lean{C2Cofinite} & $C_2\Rightarrow$ Etingof-finite & $C_2$ & Def.~\ref{l12:def:etingof}\\
\lean{finiteDimensional_}\allowbreak\lean{An_}\allowbreak\lean{of_}\allowbreak\lean{C2Cofinite} (+ \lean{A}, \lean{CGen}) & $C_2\Rightarrow\dim A_n,\dim A,\dim V/C_m<\infty$ & $C_2$; unconditional & Thm.~\ref{l12:thm:an}\\
\lean{contragredientModule}, \lean{contragredientJacobi_}\allowbreak\lean{holds} & $M'$ is an ordinary module (FHL) & none; unconditional & Thm.~\ref{l12:thm:contra}\\
\lean{onePoint_}\allowbreak\lean{contragredient_}\allowbreak\lean{holds} & $\chi_{M'}(v) = (-1)^n\chi_M(v)$ & $v\in V_n$, $L_1v = 0$ & Thm.~\ref{l12:thm:contra}\\
\lean{SelfContragredient}, \lean{Unitary} & self-dual / unitary module & \textbf{vacuous / unsatisfiable} & §\ref{l12:sec:contra}\\
\lean{fusionCoeff} & fusion rule & \textbf{$\equiv0$} (proved in the library) & Thm.~\ref{l12:thm:fusionzero}\\
\lean{IntertwinerSpace}, \lean{intertwiningDefect}, \lean{Intertwiner} & kernel of the commutator-only defect; plus lower truncation & definitions; \textbf{weaker than FHL} (no iterate, no $L_{-1}$ axiom) & Def.~\ref{l12:def:intertwiner}\\
\lean{AdmissiblyRational}, \lean{GannonRational}, \lean{Regular}, \lean{Holomorphic} & rationality notions & definitions & Defs.~\ref{l12:def:admrat}--\ref{l12:def:regular}\\
\lean{gannonRational_}\allowbreak\lean{of_}\allowbreak\lean{admissiblyRational} etc. & the lattice & named DLM/ABD/Li predicates as hypotheses; all dischargeable for $C_2$-cofinite $V$ except Li's & Thm.~\ref{l12:thm:lattice}\\
\lean{abdWeakIsAdmissible}, \lean{abdIrreducible}\allowbreak\lean{WeakIsOrdinary_holds}, \lean{abdRegularOf}\allowbreak\lean{C2AndRational_holds}, \lean{zhuAdmissible}\allowbreak\lean{Semisimple_holds} & ABD 5.7, 5.8, 4.5; Zhu & none; \textbf{proved}, unconditional & Thm.~\ref{l12:thm:litproved}\\
\lean{dlmSimpleAdmissible}\allowbreak\lean{Ordinary_of_c2Cofinite}, \lean{dlmFinitelyMany}\allowbreak\lean{IrredAdmissible_of_c2Cofinite} & DLM 8.1(c), 8.1(b) & \lean{hC2}; proved; open without $C_2$ & Thm.~\ref{l12:thm:litproved}\\
\lean{LiC2CofiniteOfRegular} & regular $\Rightarrow C_2$ (Li 3.8) & \textbf{unproved}; hypothesis of \lean{regular_iff_}\allowbreak\lean{c2Cofinite_and_}\allowbreak\lean{admissiblyRational} only & §\ref{l12:sec:litnamed}\\
\lean{selfDual_}\allowbreak\lean{iff} & self-dual $\iff L_1V_1 = 0$ & simple; unconditional & Def.~\ref{l12:def:selfdual}\\
\lean{StronglyFinite}, \lean{StronglyRational} & Creutzig--Gannon / McRae notions & definitions; satisfiable & Def.~\ref{l12:def:strong}\\
\lean{StronglyRational.regular}, \lean{.finite_Φadm}, \lean{.gannonRational'} & strongly rational $\Rightarrow$ regular, $|\Phi_{\mathrm{adm}}|<\infty$, Gannon-rational & \lean{StronglyRational}; proved, no literature hypothesis & Thm.~\ref{l12:thm:srlit}\\
\bottomrule
\end{longtable}
\endgroup

The \texttt{\#print axioms} results of the original revision of this chapter were obtained at an earlier commit (not re-run for this edition)
by running a scratch file \texttt{import VOA} with \texttt{lake env lean -DmaxSynthPendingDepth=3} (never a bare
\texttt{lake env lean}); later runs are dated where they occur below. No full \texttt{lake build} was run for this chapter.
Every headline in this chapter reports exactly \lean{[propext, Classical.choice, Quot.sound]}; \lean{VOA.pbw} (now a
theorem) is used by \lean{zhuBijective_of_voaPBW} (Chapter~\ref{ch:L13}) and, through \lean{voaPBWBottomProjection}, by
\lean{zhuAdmissibleSemisimple_holds}, \lean{dlmFinitelyManyIrredAdmissible_of_c2Cofinite} and the strongly-rational
corollaries of Theorem~\ref{l12:thm:srlit}.

\paragraph{Release audit 2026-10-27: RELEASE\_AUDIT\_ZHU.md.} The release audit of the Zhu headlines
(\texttt{Notes/\allowbreak Textbook/\allowbreak referee/\allowbreak RELEASE\_AUDIT\_ZHU.md}, finding F1) is the source of the
corrections in this revision: the modularity theorem assumes semisimplicity of $A(V)$, not rationality, and the bridge
\lean{ZhuAdmissibleSemisimple} is unproved (this last statement was wrong; see the next paragraph). It also notes that the docstring of \lean{StronglyRational.rational}
(\leanpath{StronglyRational.lean}: ``Rational in the standard Zhu / Dong--Li--Mason sense'') contradicts the ``unsettled'' note
of \leanpath{AdmissibleRationality.lean} (documentation only). The two \texttt{\#print axioms} runs added in this revision
(\lean{VOA.Huang.LWBottomRecovery.isSimpleAMod_toAMod_of_irreducible}, \lean{VOA.Huang.LWBottomRecovery.allBottomRecovers}) were
made at an earlier commit (not re-run for this edition) with \texttt{lake env lean -DmaxSynthPendingDepth=3}: standard three. No \texttt{lake build} was run.

\paragraph{Second release revision 2026-10-28: RELEASE\_AUDIT\_NEW\_DECLS.md, finding F2.} That audit found that the
implication called unproved above, ``rational $\Rightarrow A(V)$ semisimple'', is proved in the library
(\lean{VOA.Huang.ZhuSemisimple.}\allowbreak\lean{isSemisimpleRing_}\allowbreak\lean{A_of_}\allowbreak\lean{admissiblyRational}). So the first
release revision of this chapter, and the module docstrings that said the same, were wrong; the docstrings have been
corrected in the library. This revision adds Theorem~\ref{l12:thm:zhuss} and corrects the findings, the dictionary of the
named literature predicates, the strong-rationality comparison and the checklist. A scratch file was run at
an earlier commit (not re-run for this edition) with \texttt{lake env lean -DmaxSynthPendingDepth=3} from the repository root (the run is recorded in full in
Chapter~\ref{ch:L15}, end of Section~\ref{l15:sec:checklist}). The lines relevant here:
\begin{plaincode}
'VOA.Huang.ZhuSemisimple.isSemisimpleRing_A_of_admissiblyRational' depends on axioms: [propext,
 Classical.choice,
 Quot.sound]
'VOA.Huang.ZhuSemisimple.zhuAdmissibleSemisimple_of_Pr' depends on axioms: [propext, Classical.choice, Quot.sound]
'VOA.Classification.COne.V2m_stronglyRational' depends on axioms: [propext, Classical.choice, Quot.sound]
'VOA.VHat.DLM3.voaPBWBottomProjection' depends on axioms: [propext, Classical.choice, Quot.sound]
\end{plaincode}
No \texttt{lake build} was run. \texttt{Notes/RELEASE\_PLAN.md} records the umbrella builds of \lean{VOA.Modularity} and
\lean{VOA.Classification} as green (12,271 jobs, standard axioms) for the release follow-ups that use this theorem.

\paragraph{Third release revision 2026-10-28: the literature predicates.} At that revision the file
\leanpath{LiteratureInputsProved.lean} records that six of the seven named literature predicates are theorems (Theorems~\ref{l12:thm:litproved}
and~\ref{l12:thm:srlit}); the investigation behind this revision is \texttt{Notes/\allowbreak Textbook/\allowbreak referee/\allowbreak
LITERATURE\_INPUTS\_STATUS.md}. The olean of the new file was not present in the main checkout, so a scratch file consisting of a
verbatim copy of it followed by \texttt{\#print axioms} commands was compiled with \texttt{lake env lean -DmaxSynthPendingDepth=3}
from the repository root (exit status~0). Output (line breaks of the long lines joined):
\begin{plaincode}
'VOA.Classification.COne.LiteratureInputs.abdIrreducibleWeakIsOrdinary_holds' depends on axioms: [propext, Classical.choice, Quot.sound]
'VOA.Classification.COne.LiteratureInputs.abdRegularOfC2AndRational_holds' depends on axioms: [propext, Classical.choice, Quot.sound]
'VOA.Classification.COne.LiteratureInputs.zhuAdmissibleSemisimple_holds' depends on axioms: [propext, Classical.choice, Quot.sound]
'VOA.Classification.COne.LiteratureInputs.dlmFinitelyManyIrredAdmissible_of_c2Cofinite' depends on axioms: [propext, Classical.choice, Quot.sound]
'VOA.StronglyRational.regular' depends on axioms: [propext, Classical.choice, Quot.sound]
'VOA.StronglyRational.dlmFinitelyManyIrredAdmissible' depends on axioms: [propext, Classical.choice, Quot.sound]
'VOA.StronglyRational.finite_Φadm' depends on axioms: [propext, Classical.choice, Quot.sound]
'VOA.StronglyRational.gannonRational'' depends on axioms: [propext, Classical.choice, Quot.sound]
'VOA.Classification.COne.ABDWeakAdmissible.abdWeakIsAdmissible' depends on axioms: [propext, Classical.choice, Quot.sound]
'VOA.Classification.COne.AdmissibleOrdinary.dlmSimpleAdmissibleOrdinary_of_c2Cofinite' depends on axioms: [propext, Classical.choice, Quot.sound]
'VOA.regular_iff_c2Cofinite_and_admissiblyRational' depends on axioms: [propext, Classical.choice, Quot.sound]
\end{plaincode}
The last line is \texttt{\#print axioms}-clean but \emph{conditional} on its hypotheses \lean{hABD} (now dischargeable) and
\lean{hLi : LiC2CofiniteOfRegular W} (\textbf{unproved}). The umbrella build of \lean{VOA.Classification} including the new file
is reported by the coordinator as green (12,179 jobs, standard axioms); it was not re-run for this chapter. The library
docstrings of \leanpath{RationalityLattice.lean} (module header) and \leanpath{Rationality.lean} (``CITED'' list), and those of
\lean{StronglyRational.regular_of_abd} and \lean{StronglyRational} in \leanpath{StronglyRational.lean} (``carried as a
citation''), still describe the ABD/DLM predicates as citations (documentation only).

The open points of this chapter are now the converse comparison \lean{AdmissiblyRational} $\Rightarrow$ Dong--Li--Mason
rationality (the forward direction is Theorem~\ref{l12:thm:dlmrat}, under $C_2$),
\lean{LiC2CofiniteOfRegular}, and the two DLM predicates for admissibly rational VOAs not known to be $C_2$-cofinite.

%======================================================================
\section{Notes and references}\label{l12:sec:notes}
%======================================================================

\textbf{Modules, contragredients, intertwining operators.} I.~B.~Frenkel, Y.-Z.~Huang and J.~Lepowsky, ``On axiomatic
approaches to vertex operator algebras and modules'', \emph{Memoirs of the AMS} 104 (1993), no.~494 (contragredient module:
§5.2; invariant forms: §5.3; intertwining operators and fusion rules: §5.4). J.~Lepowsky and H.~Li, \emph{Introduction to
Vertex Operator Algebras and Their Representations}, Birkh\"auser (2004), Ch.~4--5. T.~Gannon, \emph{Moonshine beyond the
Monster}, Cambridge University Press (2006), §§5.3, 6.2, whose Definitions 5.3.1--5.3.2 and Lemma 5.3.3 the library
transcribes. Logarithmic (grading-restricted generalised) modules: Y.-Z.~Huang, J.~Lepowsky and L.~Zhang, ``Logarithmic
tensor category theory for generalized modules for a conformal vertex algebra, I'', arXiv:1012.4193.

\textbf{Zhu's algebra and modular invariance.} Y.~Zhu, ``Modular invariance of characters of vertex operator algebras'',
\emph{J. Amer.\ Math.\ Soc.} 9 (1996), 237--302 (Zhu's algebra: §2; $C_2$: §4.4). I.~B.~Frenkel and Y.~Zhu, ``Vertex
operator algebras associated to representations of affine and Virasoro algebras'', \emph{Duke Math.\ J.} 66 (1992),
123--168 (the bimodules $A(M)$ and the fusion-rule formula). C.~Dong, H.~Li and G.~Mason, ``Vertex operator algebras and
associative algebras'', \emph{J. Algebra} 206 (1998), 67--96 (the higher Zhu algebras $A_n(V)$; admissible modules,
Def.~2.2; semisimplicity of $A(V)$, Thm.~4.10). C.~Dong and C.~Jiang, ``Bimodules associated to vertex operator algebras'',
\emph{Math.\ Z.} 259 (2008), 799--826 (the bimodules $A_{n,m}(V)$).

\textbf{Rationality and regularity.} C.~Dong, H.~Li and G.~Mason, ``Regularity of rational vertex operator algebras'',
\emph{Adv.\ Math.} 132 (1997), 148--166, arXiv:q-alg/9508018; ``Twisted representations of vertex operator algebras'',
\emph{Math.\ Ann.} 310 (1998), 571--600, arXiv:q-alg/9509005 (Theorem~8.1); ``Modular-invariance of trace functions in orbifold
theory and generalized Moonshine'', \emph{Comm.\ Math.\ Phys.} 214 (2000), 1--56. T.~Abe, G.~Buhl and C.~Dong, ``Rationality,
regularity, and $C_2$-cofiniteness'', \emph{Trans.\ AMS} 356 (2004), 3391--3402, arXiv:math/0204021. H.~Li, ``Some finiteness
properties of regular vertex operator algebras'', \emph{J. Algebra} 212 (1999), 495--514, arXiv:math/9807077; ``Symmetric
invariant bilinear forms on vertex operator algebras'', \emph{J. Pure Appl.\ Algebra} 96 (1994), 279--297. C.~Dong and
G.~Mason, ``Rational vertex operator algebras and the effective central charge'', \emph{IMRN} 2004, arXiv:math/0201318.
T.~Creutzig and T.~Gannon, ``Logarithmic conformal field theory, log-modular tensor categories and modular forms'',
\emph{J. Phys.\ A} 50 (2017), 404004, arXiv:1605.04630. R.~McRae, ``On rationality for $C_2$-cofinite vertex operator
algebras'', arXiv:2108.01898. Y.-Z.~Huang, ``Rigidity and modularity of vertex tensor categories'', \emph{Commun.\ Contemp.\
Math.} 10 (2008), 871--911, arXiv:math/0502533, and ``Vertex operator algebras, the Verlinde conjecture and modular tensor
categories'', \emph{PNAS} 102 (2005), 5352--5356 (modular tensor categories from strongly rational VOAs; not formalised).

\textbf{$C_2$-cofiniteness, spanning sets, Poisson structures.} M.~R.~Gaberdiel and A.~Neitzke, ``Rationality, quasirationality
and finite W-algebras'', \emph{Comm.\ Math.\ Phys.} 238 (2003), 305--331, arXiv:hep-th/0009235 (Proposition~8 and Theorem~11
used in Theorem~\ref{l12:thm:an}). G.~Buhl, ``A spanning set for VOA modules'', \emph{J. Algebra} 254 (2002), 125--151.
M.~Miyamoto, ``Modular invariance of vertex operator algebras satisfying $C_2$-cofiniteness'', \emph{Duke Math.\ J.} 122
(2004), 51--91. D.~Adamovi\'c and A.~Milas, ``On the triplet vertex algebra $W(p)$'', \emph{Adv.\ Math.} 217 (2008),
2664--2699, arXiv:0707.1857 ($C_2$-cofinite, irrational). T.~Arakawa and K.~Kawasetsu, ``Quasi-lisse vertex algebras and
modular linear differential equations'', in \emph{Lie Groups, Geometry, and Representation Theory}, Progr.\ Math.\ 326 (2018),
arXiv:1610.05865. P.~Etingof and T.~Schedler, ``Poisson traces and D-modules on Poisson varieties'', \emph{GAFA} 20 (2010),
958--987, arXiv:0908.3868. T.~Arakawa, ``A remark on the $C_2$-cofiniteness condition on vertex algebras'', \emph{Math.\ Z.}
270 (2012), 559--575.

\textbf{Physics.} Fusion rules and the Verlinde formula: E.~Verlinde, ``Fusion rules and modular transformations in 2D
conformal field theory'', \emph{Nucl.\ Phys.\ B} 300 (1988), 360--376; G.~Moore and N.~Seiberg, ``Classical and quantum
conformal field theory'', \emph{Comm.\ Math.\ Phys.} 123 (1989), 177--254. Logarithmic CFT: V.~Gurarie, ``Logarithmic operators
in conformal field theory'', \emph{Nucl.\ Phys.\ B} 410 (1993), 535--549; M.~Flohr, ``Bits and pieces in logarithmic conformal
field theory'', \emph{Int.\ J. Mod.\ Phys.\ A} 18 (2003), 4497--4592; M.~R.~Gaberdiel and H.~G.~Kausch, ``A rational logarithmic
conformal field theory'', \emph{Phys.\ Lett.\ B} 386 (1996), 131--137. Chiral algebras of four-dimensional theories:
C.~Beem, M.~Lemos, P.~Liendo, W.~Peelaers, L.~Rastelli and B.~C.~van~Rees, ``Infinite chiral symmetry in four dimensions'',
\emph{Comm.\ Math.\ Phys.} 336 (2015), 1359--1433, arXiv:1312.5344. Standard CFT background: P.~Di~Francesco, P.~Mathieu and
D.~S\'en\'echal, \emph{Conformal Field Theory}, Springer (1997).

\textbf{Project records.} The retirement of \lean{WeaklyRational} and the rationality lattice are documented in the module
docstrings of \leanpath{Rationality.lean}, \leanpath{AdmissibleModule.lean}, \leanpath{AdmissibleRationality.lean} and
\leanpath{RationalityLattice.lean}; the refutation of \lean{InducedAssignment} in
\leanpath{VOA/\allowbreak Modules/\allowbreak ZhuVermaAnArchZhu534.lean}; the defects of \lean{fusionCoeff} and \lean{Unitary} in
\leanpath{VOA/\allowbreak Modules/\allowbreak CoreFixFusion.lean} and \leanpath{VOA/\allowbreak Modules/\allowbreak CoreFixUnitary.lean}.

\endgroup
